% Se pre-carga la información del estudiante sólo para poder emplear el macro de
% selección de versión (digital o impresa)
% ===============================================================================
% El estudiante debe llenar sus datos en esta sección para que la plantilla los 
% auto-importe y genere automáticamente las páginas de portada y de firmas 
% autorizadas.
% ===============================================================================
% Datos del estudiante:
% -------------------------------------------------------------------------------
% Nombre completo
\def \nombreestudiante {Sebastián Huertas Gómez}
% Carné
\def \uvgcarne {22295}
% Facultad
\def \uvgfacultad {Ingeniería}
% Carrera
\def \uvgcarrera {Ingeniería en Ciencia de la Computación y Tecnologías de la Información}

% Datos del trabajo:
% -------------------------------------------------------------------------------
% Título completo (tomado literalmente de protocolo_sebastian_v1 (3).pdf, sección 1)
\def \titulotesis {logSguarDian: Desarrollo de una librería para la detección de amenazas de seguridad en Endpoints de aplicaciones}
% Mes y año de entrega (confirmado por el estudiante: noviembre 2026)
\def \mesentrega {noviembre}
\def \anoentrega {2026}
% Asesor (según protocolo_sebastian_v1 (3).pdf, portada)
\def \nombreasesor {Ing. Erick Francisco Marroquín Rodríguez}
% El catedrático del curso (Ing. Bacilio Alexander Bolaños Lima) figura como
% primer examinador en la hoja de firmas, definido más abajo.

% Datos del tribunal examinador:
% -------------------------------------------------------------------------------
% Nombre del primer examinador (segunda línea de la hoja de firmas)
\def \nombreprimerex {Ing. Ma. Bacilio Alexander Bolaños Lima}

% Capítulos pre-definidos
% -------------------------------------------------------------------------------
% Comentar las líneas de las secciones que desean omitirse, por defecto se
% se incluyen todas.
% Se descomentan conforme se redacta cada capítulo (todavía sin archivo de
% contenido: prefacio, alcance, anexos, glosario).
\def \CAPprefacio {Prefacio}
\def \CAPantecedentes {Antecedentes}
\def \CAPalcance {Alcance}
\def \CAPmetodologia {Metodología}
\def \CAPresultados {Resultados}
\def \CAPdiscusion {Análisis de resultados}
\def \CAPanexos {Anexos}
\def \CAPglosario {Glosario}

% Formato y estilo de la plantilla
% -------------------------------------------------------------------------------
% Modo impresión: Puede des-comentar la siguiente línea para generar un documento pdf sin la portada, para cuando se desee imprimir el documento para encuadernación
% NOTA: la plantilla original traía esta línea activa por defecto, lo cual
% suprime la portada verde (main.tex hace \let\CAPportada\undefined cuando
% \printver está definido). Se comenta para que la versión digital, con
% portada, sea la que se genera por defecto.
%\def \printver {Versión del documento para impresión}

% Portada: Puede cambiarse la imagen en la portada al cambiar el nombre del 
% archivo siguiente. NOTA: debe tener la suficiente resolución para cubrir el área
% designada
\def \imagenportada {plantilla/portadacit.jpg}

% Referencias: Puede des-comentar la siguiente línea para utilizar el formato de referencias APA
%\def \usarAPA {Usar formato APA}

% Párrafo: Puede comentar la siguiente línea si desea emplear un formato de 
% párrafo distinto al establecido por defecto
\def \parpordefecto {Formato de párrafo por defecto}

% Capítulos y secciones: Puede des-comentar la siguiente línea para establecer el
% formato de los capítulos y secciones bajo el estándar original de UVG para
% trabajos de graduación. Este incluye: capítulos con numeración romana, secciones
% con letras mayúsculas, sub-secciones con números y sub-sub-secciones con letras
% minúsculas
%\def \capsecuvg {Formato UVG para capítulos y secciones}

\ifdefined\printver
    \documentclass[11pt, letterpaper, twoside, openright]{report}
\else
    \documentclass[11pt, letterpaper]{report}
\fi

% Eliminar la opción de twoside y openright si se desea generar la versión
% digital del documento en lugar de la versión impresa
%\documentclass[11pt, letterpaper, twoside, openright]{report}
\usepackage[spanish, es-nodecimaldot, es-noquoting]{babel}
% cambiar a spanish, mexico si se quiere emplear tabla en lugar de cuadro
\selectlanguage{spanish}
\usepackage[utf8]{inputenc}
\usepackage[T1]{fontenc}

\title{\titulotesis}
\author{\nombreestudiante}
\date{\today}

% Información del estudiante en el archivo datos_estudiante.tex
% ===============================================================================
% El estudiante debe llenar sus datos en esta sección para que la plantilla los 
% auto-importe y genere automáticamente las páginas de portada y de firmas 
% autorizadas.
% ===============================================================================
% Datos del estudiante:
% -------------------------------------------------------------------------------
% Nombre completo
\def \nombreestudiante {Sebastián Huertas Gómez}
% Carné
\def \uvgcarne {22295}
% Facultad
\def \uvgfacultad {Ingeniería}
% Carrera
\def \uvgcarrera {Ingeniería en Ciencia de la Computación y Tecnologías de la Información}

% Datos del trabajo:
% -------------------------------------------------------------------------------
% Título completo (tomado literalmente de protocolo_sebastian_v1 (3).pdf, sección 1)
\def \titulotesis {logSguarDian: Desarrollo de una librería para la detección de amenazas de seguridad en Endpoints de aplicaciones}
% Mes y año de entrega (confirmado por el estudiante: noviembre 2026)
\def \mesentrega {noviembre}
\def \anoentrega {2026}
% Asesor (según protocolo_sebastian_v1 (3).pdf, portada)
\def \nombreasesor {Ing. Erick Francisco Marroquín Rodríguez}
% El catedrático del curso (Ing. Bacilio Alexander Bolaños Lima) figura como
% primer examinador en la hoja de firmas, definido más abajo.

% Datos del tribunal examinador:
% -------------------------------------------------------------------------------
% Nombre del primer examinador (segunda línea de la hoja de firmas)
\def \nombreprimerex {Ing. Ma. Bacilio Alexander Bolaños Lima}

% Capítulos pre-definidos
% -------------------------------------------------------------------------------
% Comentar las líneas de las secciones que desean omitirse, por defecto se
% se incluyen todas.
% Se descomentan conforme se redacta cada capítulo (todavía sin archivo de
% contenido: prefacio, alcance, anexos, glosario).
\def \CAPprefacio {Prefacio}
\def \CAPantecedentes {Antecedentes}
\def \CAPalcance {Alcance}
\def \CAPmetodologia {Metodología}
\def \CAPresultados {Resultados}
\def \CAPdiscusion {Análisis de resultados}
\def \CAPanexos {Anexos}
\def \CAPglosario {Glosario}

% Formato y estilo de la plantilla
% -------------------------------------------------------------------------------
% Modo impresión: Puede des-comentar la siguiente línea para generar un documento pdf sin la portada, para cuando se desee imprimir el documento para encuadernación
% NOTA: la plantilla original traía esta línea activa por defecto, lo cual
% suprime la portada verde (main.tex hace \let\CAPportada\undefined cuando
% \printver está definido). Se comenta para que la versión digital, con
% portada, sea la que se genera por defecto.
%\def \printver {Versión del documento para impresión}

% Portada: Puede cambiarse la imagen en la portada al cambiar el nombre del 
% archivo siguiente. NOTA: debe tener la suficiente resolución para cubrir el área
% designada
\def \imagenportada {plantilla/portadacit.jpg}

% Referencias: Puede des-comentar la siguiente línea para utilizar el formato de referencias APA
%\def \usarAPA {Usar formato APA}

% Párrafo: Puede comentar la siguiente línea si desea emplear un formato de 
% párrafo distinto al establecido por defecto
\def \parpordefecto {Formato de párrafo por defecto}

% Capítulos y secciones: Puede des-comentar la siguiente línea para establecer el
% formato de los capítulos y secciones bajo el estándar original de UVG para
% trabajos de graduación. Este incluye: capítulos con numeración romana, secciones
% con letras mayúsculas, sub-secciones con números y sub-sub-secciones con letras
% minúsculas
%\def \capsecuvg {Formato UVG para capítulos y secciones}
\input{1-opciones_adicionales}

% Bandera que le indica a cada archivo de capítulo individual que está siendo
% ensamblado dentro de este documento global, para que no vuelva a abrir su
% propio \documentclass/\begin{document} (ver 0-preambulo_standalone.tex) ni
% imprima su propia bibliografía por capítulo.
\newcommand{\modoensamble}{}

% ==============================================================================
% DEFINICIÓN DE PAQUETES
% ==============================================================================
\usepackage{xcolor}
\usepackage{amsfonts}
\usepackage{amsmath}
\usepackage{amssymb}
\usepackage{amsthm}
\usepackage{amsfonts}
\usepackage{mathtools}
\usepackage{graphicx}
\usepackage{xfrac}
\usepackage{float}
\usepackage{mathtools}
\usepackage[hypertexnames=false]{hyperref}
% \usepackage{bookmark}
\usepackage[font=small]{caption}
\usepackage{subcaption}
% Cuadros: array (columnas L{} alineadas a la izquierda) y longtable (cuadros
% que ocupan más de una página). Se cargan aquí, no en
% 2-paquetes_y_comandos_usuario.tex, para que los capítulos de Metodología y
% Resultados no dependan de ese archivo.
\usepackage{array}
\usepackage{longtable}
\newcolumntype{L}[1]{>{\raggedright\arraybackslash}p{#1}}
\renewcommand{\arraystretch}{1.3}
\setlength{\LTpre}{0.3cm}
\setlength{\LTpost}{0.1cm}
% Margen de estiramiento para que los párrafos con palabras largas no
% generen líneas que sobresalgan del margen.
\setlength{\emergencystretch}{3em}
%\usepackage{csquotes}
\usepackage{xpatch}
\usepackage{emptypage}
% htt: permite dividir con guion el texto en \texttt (nombres de archivo,
% clases y campos), que de otro modo sobresale del margen derecho.
\usepackage[htt]{hyphenat}
\usepackage{fancyhdr}
\usepackage[backend=biber, style=ieee]{biblatex}
\ifdefined\usarAPA 
    \usepackage[backend=biber, style=apa]{biblatex}
\fi
\addbibresource{m-bibliografia.bib}

\usepackage[percent]{overpic}

\usepackage{chngcntr}

\ifdefined\CAPglosario
	%\usepackage[toc]{glossaries}
	\usepackage[numberedsection]{glossaries}
	\makeglossaries
    \newglossaryentry{rasp}
{
    name={RASP},
    description={Autoprotección de aplicaciones en tiempo de ejecución (\textit{Runtime Application Self-Protection}): mecanismo de seguridad que se ejecuta dentro de la propia aplicación y analiza sus peticiones para detectar y bloquear ataques.}
}

\newglossaryentry{waf}
{
    name={WAF},
    description={Cortafuegos de aplicaciones web (\textit{Web Application Firewall}): dispositivo o servicio perimetral que filtra el tráfico HTTP antes de que llegue a la aplicación.}
}

\newglossaryentry{modsecurity}
{
    name={ModSecurity},
    description={Motor de WAF de código abierto que se usa con reglas externas; en este trabajo se configuró con OWASP CRS.}
}

\newglossaryentry{crs}
{
    name={OWASP CRS},
    description={Conjunto de reglas genéricas de OWASP (\textit{Core Rule Set}) para ModSecurity, con niveles de paranoia que controlan qué tan estrictas son.}
}

\newglossaryentry{owasp}
{
    name={OWASP Top 10},
    description={Lista de las diez categorías de riesgo más críticas para aplicaciones web, publicada por OWASP; este trabajo usa la edición 2021.}
}

\newglossaryentry{mitre}
{
    name={MITRE ATT\&CK},
    sort={MITRE ATT&CK},
    description={Base de conocimiento de tácticas y técnicas de atacantes, usada aquí para asociar cada vector de ataque a una técnica.}
}

\newglossaryentry{cwe}
{
    name={CWE},
    description={Enumeración común de debilidades de software (\textit{Common Weakness Enumeration}): catálogo de tipos de defectos que originan vulnerabilidades.}
}

\newglossaryentry{sqli}
{
    name={SQL Injection},
    description={Ataque que inserta sentencias SQL en una entrada de la aplicación para alterar las consultas que ejecuta la base de datos.}
}

\newglossaryentry{xss}
{
    name={XSS},
    description={\textit{Cross-Site Scripting}: ataque que inyecta código ejecutable en una página que otros usuarios verán.}
}

\newglossaryentry{pathtraversal}
{
    name={Path Traversal},
    description={Ataque que manipula rutas de archivo (por ejemplo con \texttt{../}) para leer archivos fuera del directorio previsto; incluye la inclusión local de archivos (LFI).}
}

\newglossaryentry{cmdi}
{
    name={Command Injection},
    description={Ataque que logra ejecutar comandos del sistema operativo a través de una entrada de la aplicación.}
}

\newglossaryentry{middleware}
{
    name={\textit{middleware}},
    sort={middleware},
    description={Función que se ejecuta entre la recepción de una petición y la lógica de negocio de la aplicación; en Express se registra con \texttt{app.use}.}
}

\newglossaryentry{rf}
{
    name={Random Forest},
    description={Algoritmo supervisado que combina muchos árboles de decisión; en la librería es el único modelo con autoridad para bloquear una petición.}
}

\newglossaryentry{iforest}
{
    name={Isolation Forest},
    description={Algoritmo no supervisado de detección de anomalías que aísla las observaciones atípicas con particiones aleatorias; en la librería solo registra anomalías y nunca bloquea.}
}

\newglossaryentry{onnx}
{
    name={ONNX},
    description={Formato abierto para serializar modelos de aprendizaje automático, que permite ejecutarlos fuera de Python.}
}

\newglossaryentry{workerthreads}
{
    name={\texttt{worker\_threads}},
    sort={worker_threads},
    description={Módulo de Node.js que ejecuta código en hilos independientes del hilo principal; la inferencia de los modelos se ejecuta ahí.}
}

\newglossaryentry{eventloop}
{
    name={\textit{Event Loop}},
    sort={Event Loop},
    description={Bucle de eventos de Node.js que atiende las peticiones en un único hilo; bloquearlo degrada a toda la aplicación.}
}

\newglossaryentry{failopen}
{
    name={\textit{fail-open}},
    sort={fail-open},
    description={Política por la cual, ante un error o un tiempo de espera del módulo de detección, la petición se deja pasar en lugar de bloquearse.}
}

\newglossaryentry{monitor}
{
    name={Modo \texttt{monitor}},
    sort={Modo monitor},
    description={Modo de operación en el que la librería registra las peticiones que bloquearía, pero no las bloquea.}
}

\newglossaryentry{block}
{
    name={Modo \texttt{block}},
    sort={Modo block},
    description={Modo de operación en el que la librería responde HTTP 403 a las peticiones que clasifica como ataque.}
}

\newglossaryentry{pesp}
{
    name={p95},
    description={Percentil 95 del tiempo de respuesta: el valor por debajo del cual queda el 95~\% de las peticiones.}
}

\newglossaryentry{deltap95}
{
    name={$\Delta p95$},
    sort={Delta p95},
    description={Diferencia entre el p95 con el middleware activo y el p95 de la línea base sin él; en este trabajo se mide en milisegundos.}
}

\newglossaryentry{funo}
{
    name={F1},
    description={Media armónica de la precisión y el recall; es la métrica principal de clasificación del trabajo.}
}

\newglossaryentry{auc}
{
    name={AUC-ROC},
    description={Área bajo la curva ROC; mide qué tan bien separa un modelo dos clases sin depender de un umbral.}
}

\newglossaryentry{wilson}
{
    name={Intervalo de Wilson},
    description={Intervalo de confianza para una proporción, adecuado para muestras pequeñas o proporciones cercanas a 0 o 1.}
}

\newglossaryentry{casidup}
{
    name={Casi-duplicado},
    description={Registro muy parecido a otro pero no idéntico; si queda repartido entre entrenamiento y prueba, puede inflar las métricas.}
}

\newglossaryentry{minhash}
{
    name={MinHash},
    description={Técnica para estimar rápidamente la similitud entre conjuntos, usada para detectar casi-duplicados.}
}

\newglossaryentry{fuga}
{
    name={Fuga de datos},
    description={Presencia en el conjunto de evaluación de información disponible durante el entrenamiento, lo que sobreestima el desempeño.}
}

\newglossaryentry{corpuslimpio}
{
    name={Corpus limpio},
    description={Conjunto de payloads de evaluación filtrado para excluir los que coinciden con el entrenamiento; en este trabajo es el de la Ronda~5.}
}

\newglossaryentry{useragent}
{
    name={\texttt{User-Agent}},
    sort={User-Agent},
    description={Cabecera HTTP con la que el cliente se identifica (navegador, cURL, entre otros); su representación en el corpus sesgó al modelo.}
}

\newglossaryentry{falsopositivo}
{
    name={Falso positivo},
    description={Petición legítima que el sistema clasifica o bloquea como ataque.}
}

\newglossaryentry{parmodelos}
{
    name={Par de modelos},
    description={Combinación de un archivo \texttt{rf.onnx} y uno \texttt{if.onnx} identificada por su suma SHA-256 (pares A, B y C).}
}

\newglossaryentry{tarball}
{
    name={\textit{tarball}},
    sort={tarball},
    description={Archivo comprimido \texttt{.tgz} de un paquete de npm, publicado o construido localmente.}
}

\newglossaryentry{kanban}
{
    name={Kanban},
    description={Metodología de gestión basada en un tablero de columnas que representan el estado de las tareas.}
}

\newglossaryentry{artillery}
{
    name={Artillery},
    description={Herramienta de pruebas de carga usada para medir la latencia con y sin el middleware.}
}

\newglossaryentry{sqlmap}
{
    name={sqlmap},
    description={Herramienta de código abierto que automatiza la detección y la explotación de inyecciones SQL; generó los payloads de SQL Injection del corpus limpio.}
}

\newglossaryentry{zap}
{
    name={OWASP ZAP},
    description={Herramienta de OWASP para pruebas de seguridad de aplicaciones web; se usó para confirmar el payload de XSS.}
}

\newglossaryentry{mcp}
{
    name={MCP},
    description={Protocolo de Contexto de Modelo (\textit{Model Context Protocol}): estándar para que un asistente de inteligencia artificial use herramientas externas.}
}

\fi

% ==============================================================================
% MÁRGENES Y FORMATO GENERALES
% ==============================================================================
\usepackage[top=1in, left=1.5in, right=1in, bottom=1in]{geometry}
%Options: Sonny, Lenny, Glenn, Conny, Rejne, Bjarne, Bjornstrup
\usepackage[Sonny]{fncychap}

% ==============================================================================
% DEFINICIONES DE LA PLANTILLA
% ==============================================================================
\graphicspath{ {figuras/} }
\definecolor{uvg-green}{RGB}{17,71,52}
\newcommand{\defaultparformat}[1]{
	{\setlength{\parskip}{2ex}
     \input{#1}}
}
\ifdefined\capsecuvg
	\renewcommand\thechapter{\Roman{chapter}}
    \renewcommand\thesection{\Alph{section}}
	\renewcommand\thesubsection{\arabic{subsection}}
    \renewcommand\thesubsubsection{\alph{subsubection}}
\fi
\counterwithout{figure}{chapter}
\counterwithout{table}{chapter}
\counterwithout{equation}{chapter}

\newcommand{\blankpage}{
\newpage
\thispagestyle{empty}
\mbox{}
\newpage
}
% ==============================================================================

% Comandos definidos por el usuario en el archivo comandos_usuario.tex
\input{2-paquetes_y_comandos_usuario}

% ==============================================================================
% CUERPO DEL TRABAJO
% ==============================================================================
\pagestyle{headings}
\begin{document}

% Los archivos de capítulo ahora traen su propio \chapter{} y ya no se
% envuelven con \defaultparformat (ese mecanismo asumía que main.tex ponía el
% \chapter{} y el archivo era solo cuerpo). Se aplica aquí una sola vez, para
% todo el documento, el mismo \parskip que \defaultparformat aplicaba por
% sección, para no perder el formato de párrafo por defecto de la plantilla.
\ifdefined\parpordefecto
	\setlength{\parskip}{2ex}
\fi

% ==============================================================================
% PORTADA
% ==============================================================================
\ifdefined\printver
    \let\CAPportada\undefined
\fi 

\ifdefined\CAPportada
    \input{portada}
\fi

% ==============================================================================
% PRIMERAS PÁGINAS (Carátulas más hojas de guarda)
% ==============================================================================
\ifdefined\CAPcaratula
	\input{caratula}
\fi

% ==============================================================================
% HOJA DE FIRMAS
% ==============================================================================
\ifdefined\CAPfirmas
	\newpage
	\cleardoublepage\phantomsection
	\thispagestyle{empty}
	\vspace*{0.5in}
	\large Vo.Bo.:\\[1cm]
	\begin{center}
		(f) \rule[1pt]{4 in}{1pt}\\
		\nombreasesor
	\end{center}
	\vspace{1in}

	Tribunal Examinador:\\[1cm]
	\begin{center}
		(f) \rule[1pt]{4 in}{1pt}\\
		\nombreasesor \\[1in]
		(f) \rule[1pt]{4 in}{1pt}\\
		\nombreprimerex
	\end{center}
	\vspace{1in}

	Fecha de aprobación: Guatemala, \rule[1pt]{0.5 in}{1pt} de \rule[1pt]{1 in}{1pt} de \rule[1pt]{1 in}{1pt}.
	\normalsize
\fi

% Comentar para formato estilo libro en la numeración de páginas (NO 
% compatible con la guía UVG 2019)
\pagestyle{plain}
% ==============================================================================
% CONTENIDO DEL TRABAJO
% ==============================================================================
% PREFACIO
% ------------------------------------------------------------------------------
\ifdefined\CAPprefacio
	\newpage
	\cleardoublepage\phantomsection
    \chapter*{Prefacio}
    \ifdefined\parpordefecto
    	\defaultparformat{a-prefacio}
    \else
    	% ==============================================================================
% PREFACIO: DECLARACIÓN SOBRE EL USO DE INTELIGENCIA ARTIFICIAL
% ==============================================================================
% Adaptado (2026-10-12) de la estructura de la declaración del otro integrante:
% rol del autor, herramientas empleadas, y cierre de responsabilidad.
%
% Este tomo todavía no tiene agradecimientos. Cuando existan, van antes de la
% declaración (esta es el último bloque del prefacio).
%
% En main.tex, el \chapter*{Prefacio} lo pone el documento principal; en modo
% independiente se agrega aquí.
\ifdefined\modoensamble
\else
	\input{0-preambulo_standalone}
	\chapter*{Prefacio}
\fi

\subsection*{Declaración sobre el uso de inteligencia artificial}

El autor definió el enfoque técnico de su componente, seleccionó los vectores de ataque y su
mapeo a OWASP Top 10 y MITRE ATT\&CK, auditó y etiquetó los datos, construyó el entorno de
evaluación, diseñó la capa de captura y respuesta (respuesta HTTP 403, registro de eventos,
\textit{webhooks} e interfaz de línea de comandos), calibró el umbral de bloqueo, elaboró el
corpus de evaluación, ejecutó las pruebas de penetración, las mediciones de latencia, memoria y CPU
y la validación en DVNA, node-api-goat y dvws-node, e implementó el extractor inicial de 72
características y el paso a \textit{worker\_threads}. Decidió además el alcance y los
resultados que se reportan.

Para la búsqueda bibliográfica empleó Gemini (Google) como herramienta de apoyo para identificar
fuentes y temas potencialmente relevantes; leyó las fuentes originales y conservó únicamente las que
respaldaban lo que se afirma en el texto. Para la documentación técnica, la redacción y la
verificación de coherencia de cifras, criterios y atribución de roles entre este trabajo y el del
componente de Inteligencia Artificial empleó a Claude, de Anthropic, mediante conversación, y a
Claude Code, mediante un agente de código, que además se usó para refactorizar código y como
revisor adicional de \textit{pull requests}; esa revisión complementó, y no sustituyó, la revisión
cruzada obligatoria entre los integrantes del proyecto descrita en la
Sección~\ref{subsec:met-flujo}. Ninguna de estas herramientas se utilizó para generar datos,
mediciones ni resultados: las cifras provienen de los archivos y las ejecuciones de los
repositorios del proyecto, y las derivadas de ellos (medianas, intervalos de confianza) se
calcularon a partir de esos archivos. El uso de estas herramientas se declara conforme a la \textit{Guía para el uso ético y responsable de la inteligencia artificial} de la Universidad del Valle de Guatemala. El autor dirigió ese trabajo, contrastó las cifras con sus
fuentes, revisó las referencias bibliográficas contra las fuentes originales, salvo las señaladas
en la bibliografía como no accesibles, puede explicar el contenido y asume la responsabilidad del
documento.

\ifdefined\modoensamble
\else
	\end{document}
\fi

    \fi
    \addcontentsline{toc}{chapter}{Prefacio}
\fi

% ÍNDICE GENERAL
% ------------------------------------------------------------------------------
\ifdefined\CAPindice
	\newpage
    \cleardoublepage\phantomsection
	\renewcommand{\contentsname}{Índice}
    %\phantomsection
    \pdfbookmark{\contentsname}{toc}
    %\pdfbookmark{Índice}{toc}
	\tableofcontents
\fi

% LISTADO DE FIGURAS / LISTADO DE CUADROS
% ------------------------------------------------------------------------------
\ifdefined\CAPfiguras
	\newpage
    \cleardoublepage\phantomsection
	\renewcommand{\listfigurename}{Lista de figuras}
	\listoffigures
	\addcontentsline{toc}{chapter}{Lista de figuras}
\fi

\ifdefined\CAPcuadros
	\newpage
    \cleardoublepage\phantomsection
	\renewcommand{\listtablename}{Lista de cuadros}
	\listoftables
	\addcontentsline{toc}{chapter}{Lista de cuadros}
\fi

% RESUMEN, ABSTRACT, INTRODUCCIÓN, ANTECEDENTES, JUSTIFICACIÓN, OBJETIVOS,
% ALCANCE, MARCO TEÓRICO, CAPÍTULOS, CONCLUSIONES, RECOMENDACIONES
% ------------------------------------------------------------------------------
% PENDIENTES: ninguno redactado todavía. Se están construyendo uno por uno.
% CAPresumen, CAPabstract, CAPintroduccion, CAPjustificacion, CAPobjetivos,
% CAPmarcoteorico, CAPconclusiones y CAPrecomendaciones están definidos de
% forma incondicional en 1-opciones_adicionales.tex (no se debe modificar ese
% archivo), así que sus bloques se comentan aquí en vez de tocar esas
% banderas. El bloque de "CAPÍTULOS" (j-capitulos) tampoco tiene bandera
% \ifdefined en la plantilla original -- se ejecuta siempre -- así que
% también se comenta directamente.
%
\ifdefined\CAPresumen
	\newpage
	\cleardoublepage\phantomsection
	\chapter*{Resumen}
	\ifdefined\modoensamble
\else
	\input{0-preambulo_standalone}
	\chapter*{Resumen}
\fi

Las startups que desarrollan servicios sobre Node.js rara vez cuentan con un equipo de seguridad dedicado. Este trabajo presenta logSguarDian, una librería npm que implementa un middleware de autoprotección en tiempo de ejecución (RASP) para Express y detecta inyección SQL, Cross-Site Scripting, Path Traversal e inyección de comandos. El middleware intercepta cada petición, extrae un vector de 76 características y un Random Forest decide el bloqueo (umbral de 0.35); un Isolation Forest registra las anomalías y nunca bloquea. Ambos modelos se ejecutan en formato ONNX (Open Neural Network Exchange) dentro de hilos de trabajo (\texttt{worker\_threads}). Este documento corresponde al componente de ciberseguridad: la selección y el mapeo de los vectores de ataque a OWASP y MITRE ATT\&CK, un corpus de 383{,}202 peticiones etiquetadas, el entorno de evaluación y las pruebas de penetración. Sobre el conjunto de prueba, el F1 fue de 0.9239 a 0.9956 por categoría. Sobre payloads no vistos, la versión 0.1.1 bloqueó 99.5~\% en inyección SQL, 100.0~\% en Path Traversal y 92.5~\% en inyección de comandos; en Cross-Site Scripting solo se confirmó un caso. El sobrecosto de latencia fue de 12.08~ms en carga normal y de 12.85~ms en carga de ataque, dentro del límite de 15~ms, pero el consumo de memoria (+332~MiB) y de CPU superó los criterios. Con tráfico legítimo de aplicaciones de terceros la librería generó bloqueos indebidos, por lo que no es apta para modo de bloqueo sin calibración previa. El objetivo general se cumplió de forma parcial.

\medskip
\noindent\textbf{Palabras clave:} RASP, middleware, Node.js, detección de intrusiones, Random Forest, Isolation Forest.

\ifdefined\modoensamble
\else
	\end{document}
\fi

	\addcontentsline{toc}{chapter}{Resumen}
\fi

\ifdefined\CAPabstract
	\newpage
	\cleardoublepage\phantomsection
	\chapter*{Abstract}
	\ifdefined\modoensamble
\else
	\input{0-preambulo_standalone}
	\chapter*{Abstract}
\fi

Startups that build services on Node.js rarely have a dedicated security team. This work presents logSguarDian, an npm library that implements a runtime application self-protection (RASP) middleware for Express and detects SQL injection, Cross-Site Scripting, Path Traversal and command injection. The middleware intercepts each request and extracts a 76-feature vector; a Random Forest decides whether to block (threshold 0.35), while an Isolation Forest records anomalies and never blocks. Both models run in ONNX (Open Neural Network Exchange) format inside worker threads. This document covers the cybersecurity component: the selection and mapping of attack vectors to OWASP and MITRE ATT\&CK, a corpus of 383,202 labeled requests, the evaluation environment and the penetration tests. On the held-out test set, the F1 score ranged from 0.9239 to 0.9956 per category. On unseen payloads, version 0.1.1 blocked 99.5\% of SQL injection, 100.0\% of Path Traversal and 92.5\% of command injection; only one cross-site scripting case was confirmed. The latency overhead was 12.08~ms under normal load and 12.85~ms under attack load, within the 15~ms limit, but memory (+332~MiB) and CPU consumption exceeded the criteria. With legitimate traffic from third-party applications the library produced false blocks, so it is not suitable for blocking mode without prior calibration. The general objective was partially met.

\medskip
\noindent\textbf{Keywords:} RASP, middleware, Node.js, intrusion detection, Random Forest, Isolation Forest.

\ifdefined\modoensamble
\else
	\end{document}
\fi

	\addcontentsline{toc}{chapter}{Abstract}
\fi

\ifdefined\CAPintroduccion
	\newpage
	\cleardoublepage
	% ==============================================================================
% INTRODUCCIÓN
% ==============================================================================
% Adaptado de protocolo_sebastian_v1 (3).pdf, sección 1 "Descripción General
% del Trabajo" y subsección 1.1 "Entorno tecnológico". Se conserva el
% contenido y las citas del protocolo; se reescribió el registro (de
% "protocolo" a capítulo de tesis) y se dejó fuera, a propósito, el detalle
% exhaustivo de los cinco entregables y de la Definición del Problema
% (protocolo, secciones 1 y 2): esos dos bloques se reservan para los
% capítulos de Alcance/Metodología y Antecedentes/Justificación
% respectivamente, para no duplicar contenido entre capítulos.
%
% El último párrafo (organización del documento) es una convención de
% redacción estándar de introducciones de tesis, no contenido tomado del
% protocolo.
\ifdefined\modoensamble
\else
	\input{0-preambulo_standalone}
\fi

\pagenumbering{arabic}
\setcounter{page}{1}
\chapter{Introducción}

Las aplicaciones web modernas exponen interfaces de programación cuya superficie de ataque incluye la inyección SQL, el \textit{Cross-Site Scripting}, el \textit{Path Traversal} y la inyección de comandos del sistema operativo. Estas categorías figuran entre los riesgos más críticos del estándar OWASP Top 10 \cite{owasp2021} y comparten una propiedad: la evidencia del ataque reside en una sola petición HTTP. Las \textit{startups} y empresas tecnológicas que desarrollan servicios sobre Node.js rara vez cuentan con un equipo de seguridad dedicado. Este trabajo parte de la premisa de que, para esas organizaciones, las soluciones perimetrales resultan costosas de operar y mantener. De ahí surge la necesidad de una herramienta de protección que opere dentro de la propia aplicación.

El presente trabajo de graduación corresponde a la contribución individual de Sebastián
Huertas Gómez al proyecto integrador \emph{logSguarDian}, desarrollado en paralelo con el
estudiante Diego Valenzuela (carné 22309), quien presenta un trabajo independiente
correspondiente al componente de Inteligencia Artificial del mismo proyecto. Ambos
trabajos comparten el contexto general de \emph{logSguarDian}, pero los entregables,
objetivos específicos, metodología y evaluación documentados en este informe corresponden
exclusivamente a la contribución del autor en el área de ciberseguridad. Los puntos de
integración técnica con el componente de Inteligencia Artificial se identifican
explícitamente a lo largo del documento como interfaces (entrada/salida) y no como
entregables propios.

El trabajo consiste en el diseño, desarrollo e implementación de una librería de software
distribuida a través de npm, orientada a la detección de amenazas de seguridad en
endpoints de aplicaciones desarrolladas en Node.js. Esta integración tecnológica busca
democratizar el acceso a capacidades avanzadas de Inteligencia Artificial, transformando la
defensa tradicional y reactiva en un análisis probabilístico y predictivo en tiempo real. El
proyecto tiene como finalidad proporcionar una solución accesible, automatizada y
técnicamente fundamentada que permita a startups y empresas sin departamento formal de
ciberseguridad fortalecer la protección de sus servicios web frente a ataques comunes y
avanzados.

La librería se integra como \textit{middleware} dentro de aplicaciones Node.js y analiza cada
petición HTTP dirigida a los \textit{endpoints} antes de que alcance la lógica de negocio. Para limitar su impacto sobre el rendimiento del servidor, implementa un \textit{pipeline} de ingeniería de
características que transforma cada petición en un vector numérico de 76 características y ejecuta
la inferencia en hilos de trabajo (\texttt{worker\_threads}), fuera del \textit{event loop} de la
aplicación. El análisis es petición a petición: el sistema no correlaciona peticiones entre sí, y
las características temporales que calcula el extractor no forman parte de la entrada de los
modelos. A partir de ese vector, el sistema identifica patrones asociados a inyección SQL (SQLi),
\textit{Cross-Site Scripting} (XSS), \textit{Path Traversal}/\textit{Local File Inclusion} (LFI) e
inyección de comandos del sistema operativo (\textit{Command Injection}, CMDi), y registra como anómalas
las peticiones que se desvían de la distribución del tráfico legítimo.

La detección se apoya en modelos de aprendizaje automático desarrollados por el componente de
Inteligencia Artificial del proyecto, sobre criterios técnicos de ciberseguridad definidos en este
trabajo: la selección de vectores de ataque y su alineación con OWASP y MITRE ATT\&CK. El componente
de Inteligencia Artificial aporta un modelo híbrido: un clasificador supervisado (Random Forest) que
clasifica las cuatro categorías de ataque y es el único con autoridad para bloquear una petición, y
un detector de anomalías no supervisado (Isolation Forest) que modela el tráfico legítimo y marca en
el registro las peticiones atípicas, sin bloquearlas. Ambos modelos se ejecutan dentro del proceso
de la aplicación, sin servicios externos, y se entrenaron con un corpus etiquetado construido a
partir de fuentes públicas, complementado con conjuntos sintéticos propios.

El proyecto se sitúa en el entorno de aplicaciones web empresariales desarrolladas con
Node.js, específicamente en interfaces de programación de aplicaciones (API) REST utilizadas por startups y empresas tecnológicas en
Guatemala que carecen de un equipo dedicado de seguridad informática.

\section*{Entorno tecnológico}

Las startups y empresas tecnológicas guatemaltecas que constituyen el contexto de este
proyecto comparten un conjunto de características de infraestructura bien definidas. En
cuanto a proveedores de nube, la adopción regional se concentra principalmente en Amazon
Web Services (AWS), Microsoft Azure y Google Cloud Platform (GCP), que en conjunto dominan
el mercado global de infraestructura cloud con cuotas del 30~\%, 21~\% y 12~\% respectivamente en el cuarto trimestre de 2024 \cite{synergy2025cloud}. Para el contexto de este proyecto se asume que las startups guatemaltecas despliegan sobre alguno de
estos tres proveedores o sobre proveedores de infraestructura más simples y de menor costo, como
DigitalOcean o Linode. El diseño de la librería no depende de ningún proveedor, porque se ejecuta
dentro del proceso de la aplicación.

En términos de arquitectura, las startups con servicios basados en Node.js tienden a
adoptar uno de tres patrones predominantes. El primero es la arquitectura monolítica
modular con Express.js o Fastify, típica de equipos pequeños y proyectos en etapa inicial,
donde todos los endpoints se sirven desde un único proceso Node.js desplegado en una
instancia virtual (EC2, Droplet o equivalente). El segundo es la arquitectura de
microservicios, en la que cada dominio funcional (autenticación, pagos, notificaciones) se
implementa como un servicio independiente que expone su propia API REST o GraphQL, con
comunicación entre servicios a través de HTTP o colas de mensajes como RabbitMQ. El tercer patrón es la arquitectura serverless, donde funciones individuales se despliegan en servicios como AWS Lambda o
Google Cloud Functions, eliminando la gestión directa de servidores, pero introduciendo
nuevos vectores de ataque a nivel de endpoint. Express.js, NestJS y Fastify figuran entre los frameworks web más usados según la encuesta de Stack
Overflow de 2025 (19.9~\%, 6.7~\% y 2.9~\% del total de encuestados, respectivamente)
\cite{stackoverflow2025survey}; a ellos se suma Koa, que esa encuesta no cuantifica.

En todos estos patrones, los endpoints HTTP representan la superficie de ataque principal.
Las aplicaciones suelen operar detrás de un proxy inverso como Nginx o un API Gateway
administrado, pero la lógica de negocio y la validación de entradas recaen directamente
sobre el código Node.js. La ausencia de departamentos de seguridad dedicados en la mayoría
de estas organizaciones implica que estas capas de lógica raramente cuentan con mecanismos
automatizados de detección de patrones anómalos, lo que constituye la brecha que este
proyecto pretende cerrar.

La solución propuesta se integra directamente en el ciclo de ejecución del servidor,
actuando como una capa adicional de protección que complementa mecanismos tradicionales
como validaciones de entrada y controles de autenticación. De esta manera, la librería no
reemplaza la necesidad de buenas prácticas de desarrollo seguro, sino que funciona como un
sistema de detección y apoyo a la gestión de riesgos en entornos con recursos limitados.

\section*{Organización del documento}

El resto del documento se organiza de la siguiente forma: el Capítulo~2 (Antecedentes) y el
Capítulo~3 (Justificación) desarrollan el problema y las deficiencias de las soluciones existentes;
el Capítulo~4 establece los objetivos y el Capítulo~5 delimita el alcance; el Capítulo~6 (Marco
teórico) cubre los fundamentos de Node.js, de la seguridad de aplicaciones web y de las técnicas de
detección empleadas; el Capítulo~7 (Metodología) describe el enfoque, los procedimientos y las
métricas; los Capítulos~8 y~9 presentan y analizan los resultados; y los Capítulos~10 y~11 exponen las conclusiones y las recomendaciones. Los anexos reúnen la
procedencia de los modelos y de los paquetes publicados, el historial de bloqueo del conjunto de
prueba, la auditoría de fuentes y la evidencia de la corrida de latencia, y el glosario define los
términos técnicos empleados.

\ifdefined\modoensamble
\else
	\section*{Referencias}
	\printbibliography[heading=none]
	\end{document}
\fi

\fi

\ifdefined\CAPantecedentes
	\newpage
	% ==============================================================================
% ANTECEDENTES
% ==============================================================================
% Contenido redactado y provisto directamente por el estudiante (no adaptado
% del protocolo). Se aplicaron dos correcciones editoriales antes de
% incorporarlo:
%   1. La segunda tabla (comparativa de soluciones RASP) usaba columnas "l"
%      sin ancho fijo con encabezados largos ("Detección de variantes
%      desconocidas"), lo que se sale del margen de la página. Se cambió a
%      columnas p{} con \small, igual que la primera tabla.
%   2. La entrada bibliográfica mereani2019 traía \type{@article} con campos
%      de conferencia (booktitle, organization); se corrigió a
%      \type{@inproceedings}, que es el tipo correcto en biblatex para esos
%      campos.
\ifdefined\modoensamble
\else
	\input{0-preambulo_standalone}
\fi

\chapter{Antecedentes}

\section{Estado del arte del problema}
La detección automatizada de ataques a aplicaciones web mediante técnicas de aprendizaje
automático ha sido un área de investigación activa desde al menos la década pasada, pero su
consolidación como línea aplicable a arquitecturas de producción reales es reciente. La literatura
puede organizarse en tres corrientes que convergen directamente en el problema que aborda este
trabajo: (i) modelos de aprendizaje automático especializados en la detección de un vector de
ataque específico (inyección SQL, XSS, path traversal o inyección de comandos) evaluados sobre
conjuntos de datos aislados; (ii) soluciones comerciales de protección en tiempo de ejecución
(RASP) orientadas a entornos empresariales con presupuesto de seguridad dedicado; y (iii) el
contexto regulatorio y de madurez tecnológica de Guatemala, donde tanto la oferta de herramientas
de seguridad accesibles como el marco legal que las incentivaría se encuentran en un estado
incipiente. Ninguna de estas tres corrientes, por separado, resuelve el problema identificado en la
sección 2 del protocolo de este trabajo: la ausencia de un mecanismo de detección multivector,
integrado nativamente en el runtime de Node.js, distribuible sin costo y operable por equipos sin
un analista de seguridad dedicado. Este capítulo revisa los antecedentes más relevantes de cada
corriente, identifica sus aportes y limitaciones, y sintetiza el vacío de investigación que justifica la
propuesta de \textit{logSguarDian}.

\section{Contexto de ciberseguridad y marco regulatorio en Guatemala}
El ecosistema tecnológico guatemalteco presenta una combinación de crecimiento acelerado en la
adopción de aplicaciones web y ausencia casi total de controles de seguridad estructurados. El
Índice de Delincuencia Organizada reporta un aumento de los delitos cibernéticos en el país que atribuye a la falta de regulación y
control en internet \cite{lahora2025ley6347}. Esta carencia regulatoria no es un dato
anecdótico: constituye en sí misma un antecedente directo del problema, porque en ausencia de
un marco legal que obligue a las organizaciones a invertir en controles mínimos, la adopción de
herramientas de seguridad depende exclusivamente de su accesibilidad económica y operativa, lo
que es precisamente el vacío que este proyecto busca cerrar.

Guatemala cuenta desde 2024 con la Iniciativa de Ley 6347, "Ley de Ciberseguridad", que propone
la tipificación de nuevos delitos informáticos y la creación de un Centro de Respuesta a Incidentes
de Seguridad Informática (CSIRT-GT) como órgano técnico nacional \cite{congreso2025dictamen}.
La iniciativa recibió dictamen favorable de la Comisión de Asuntos de Seguridad Nacional en
agosto de 2025 y fue aprobada en segunda lectura por el pleno del Congreso el 23 de septiembre
de 2025; sin embargo, a la fecha de redacción de este documento permanece sin aprobación final,
con su tercer debate reactivado hasta abril de 2026 tras más de seis meses de estancamiento
\cite{lahora2026estancada}. Un análisis técnico-jurídico de la Alianza por un Congreso Eficiente sobre la iniciativa 6347 señaló
además ambigüedades en varios tipos penales que comprometen el principio de legalidad y recomendó
una revisión legislativa antes de su aprobación \cite{cien2025boletin}. Este antecedente
regulatorio es relevante para el posicionamiento del proyecto por dos razones. Primero, confirma
que, incluso si la ley se aprueba, su enfoque es punitivo y de respuesta institucional a incidentes,
no preventivo a nivel de aplicación; no existe en la iniciativa ningún mecanismo que reduzca la
brecha de herramientas de defensa activa para las startups que constituyen la mayoría del tejido
empresarial guatemalteco. Segundo, mientras el marco legal no entre en vigor, la adopción de
controles de seguridad en el país seguirá siendo enteramente voluntaria, lo que refuerza la
necesidad de que dichos controles sean gratuitos y de instalación trivial para que exista algún
incentivo real de adopción.

\section{Trabajos académicos en detección de ataques web mediante aprendizaje automático}
La aplicación de aprendizaje automático a la detección de ataques específicos de aplicaciones web
ha producido resultados cuantitativos sólidos, pero con limitaciones estructurales consistentes en
su diseño experimental. El Cuadro \ref{cuadro:trabajosprevios} resume los trabajos más
relevantes para el contexto de este proyecto.

\begin{table}[h]
\small
\begin{tabular}{|p{2.1cm}|p{1.6cm}|p{2.2cm}|p{1.6cm}|p{1.6cm}|p{3.3cm}|}
\hline
\textbf{Trabajo} & \textbf{Vector} & \textbf{Técnica} & \textbf{Dataset} & \textbf{Métrica} & \textbf{Limitación principal} \\ \hline
Peralta-Garcia et al. \cite{peralta2024sql} & SQL Injection & Random Forest, árbol de decisión, SVM & Dataset público de 22{,}764 registros & Accuracy = 99~\%; F1 = 98~\% (Random Forest) & Solo cubre SQLi; el modelo se despliega en una arquitectura de microservicios que valida peticiones, es decir, como servicio separado de la aplicación \\ \hline
Mereani \& Howe \cite{mereani2019} & XSS & SVM, k-NN, Random Forest & Código JavaScript etiquetado & No reportado & Clasifica código del lado del cliente, no parámetros de peticiones entrantes al servidor \\ \hline
Wang et al. \cite{seyyar2024} & Command Injection & CCBA: CNN de doble canal, BiLSTM y atención & Dataset propio de inyección de comandos; evaluación cruzada en CSIC 2010 \cite{torrano2010csic} y HTTPParam & Accuracy = 99.3~\%; recall = 98.2~\% (dataset propio) & Modelo de aprendizaje profundo evaluado en una GPU; no reporta latencia de inferencia ni despliegue como middleware \\ \hline
SAP node-rasp \cite{saparchive2019} & SQLi, NoSQLi, Path Traversal & Análisis léxico con reglas (no ML) & No aplica & No reportado & Requiere modificar el runtime de Node.js; abandonado desde LTS 10; no distribuible como paquete npm \\ \hline
Durmuşkaya y Bayraklı \cite{durmuskaya2025waf} & SQLi, XSS, CMDi, LFI & Árbol de decisión (comparado con KNN, regresión logística, Naive Bayes y SVM) & CSIC 2010 \cite{torrano2010csic}, HTTPParams 2015 y conjunto propio & Accuracy = 93.27~\% (WAF en vivo) & Opera como proxy inverso externo (mitmproxy, Python); evaluado contra DVWA; sin cifra de latencia reportada \\ \hline
\end{tabular}
\caption[Comparativa de trabajos previos en detección de ataques web con ML]{Comparativa de trabajos previos en detección de ataques a aplicaciones web mediante aprendizaje automático}
\label{cuadro:trabajosprevios}
\end{table}
{\footnotesize Nota. Elaboración propia con base en los trabajos citados en cada fila. Las métricas son las que reporta cada autor sobre sus propios datos y no son directamente comparables entre sí. WAF: \textit{Web Application Firewall}; SQLi: \textit{SQL Injection}; XSS: \textit{Cross-Site Scripting}; CMDi: \textit{Command Injection}; LFI: \textit{Local File Inclusion}.}

El análisis de estos cinco trabajos revela un patrón consistente antes que hallazgos aislados. Con la excepción parcial de Durmuşkaya y Bayraklı \cite{durmuskaya2025waf}, quienes despliegan en su
WAF un único árbol de decisión binario aunque evalúan el desempeño por vector con modelos separados,
la totalidad de los trabajos revisados entrena y evalúa un clasificador por vector de ataque de forma
independiente, sin considerar que en un escenario de producción real un mismo endpoint puede
recibir simultáneamente intentos de inyección SQL, XSS y path traversal. Incluso el trabajo más cercano en arquitectura a la propuesta de este proyecto, el WAF de Durmuşkaya
y Bayraklı \cite{durmuskaya2025waf}, resuelve el problema mediante un proxy inverso externo
(mitmproxy, implementado en Python), lo que introduce un
salto de red adicional y no captura el contexto interno del proceso Node.js (memoria del \texttt{event
loop}, tiempos de procesamiento internos, encadenamiento de middlewares) que sí está disponible
para una librería integrada directamente como middleware. Ninguno de los trabajos revisados
reporta el impacto de su solución sobre la latencia del servidor protegido bajo carga sostenida, dato
que es un criterio de aceptación explícito de este proyecto.

Es pertinente señalar también una limitación metodológica: las métricas reportadas no son homogéneas entre los trabajos revisados (\textit{accuracy} y F1 en
\cite{peralta2024sql}; \textit{accuracy} y \textit{recall} en \cite{seyyar2024}; \textit{accuracy} en \cite{durmuskaya2025waf}), lo que impide una comparación directa, y confiar en la
\textit{accuracy} sola es una práctica problemática cuando los datasets de ataques suelen
estar desbalanceados a favor del tráfico legítimo, ya que una accuracy alta puede ocultar una tasa
de falsos negativos inaceptable para un sistema de seguridad en producción.

\section{Soluciones RASP existentes: antecedentes tecnológicos}
\label{sec:rasp-antecedentes}
\label{sec:rasp-existentes}
Además de los antecedentes académicos, es necesario revisar los desarrollos tecnológicos
disponibles comercialmente bajo la categoría de Runtime Application Self-Protection (RASP), ya
que representan el estado práctico de la industria frente al mismo problema. El Cuadro
\ref{cuadro:rasp} compara las soluciones más relevantes.

\begin{table}[h]
\small
\begin{tabular}{|p{2.4cm}|p{1.6cm}|p{2cm}|p{1.8cm}|p{2.6cm}|p{1.8cm}|}
\hline
\textbf{Solución} & \textbf{Soporte Node.js} & \textbf{Costo} & \textbf{Plataforma externa} & \textbf{Detección de variantes desconocidas} & \textbf{Operación autónoma} \\ \hline
Contrast Protect \cite{contrastsecurity} & Sí & Licencia empresarial & Sí & Sí (IAST + RASP) & No \\ \hline
Datadog App and API Protection (antes ASM; incorpora Sqreen) \cite{datadog2026nodejs,datadog2021sqreen} & Sí & Suscripción & Sí & Parcial & No \\ \hline
OpenRASP (Baidu) \cite{baidu2026openrasp} & No & Gratuito & No & Limitado & Sí \\ \hline
\textit{logSguarDian} (propuesto) & Sí & Gratuito (npm) & No & Sí (medida sobre corpus limpio) & Sí \\ \hline
\end{tabular}
\caption[Comparativa de soluciones RASP existentes]{Comparativa entre soluciones RASP comerciales existentes y la propuesta de este trabajo}
\label{cuadro:rasp}
\end{table}
{\footnotesize Nota. Elaboración propia con base en la documentación de cada producto
\cite{contrastsecurity,datadog2026nodejs,baidu2026openrasp}. Las columnas de costo, plataforma
externa y detección de variantes desconocidas reflejan el criterio del autor y no una tarifa
publicada.}

Esta comparación establece un antecedente tecnológico claro: existe tecnología RASP madura en
el mercado, pero ninguna de las soluciones revisadas combina soporte nativo para Node.js,
gratuidad, y operación sin dependencia de una plataforma externa de análisis. Contrast Protect y Datadog, las opciones comerciales revisadas con soporte para Node.js, dependen, según el criterio del autor, de una plataforma del proveedor, lo que introduce costo recurrente y consideraciones de
privacidad de datos que resultan prohibitivas para una startup guatemalteca sin presupuesto de
seguridad. OpenRASP, la única alternativa gratuita entre las soluciones revisadas, ofrece agentes solo para
Java y PHP \cite{baidu2026openrasp}, y su enfoque
de detección basado en reglas es más limitado frente a variantes desconocidas que un clasificador
entrenado sobre la distribución estadística de los ataques.

\section{Análisis comparativo y vacíos identificados}
La revisión conjunta de los antecedentes académicos y tecnológicos permite identificar tres brechas
consistentes que ningún trabajo previo resuelve de forma simultánea.

La primera es la \textbf{brecha de integración en runtime}. Los trabajos académicos de detección de
SQLi y XSS revisados operan sobre datasets de cadenas de texto o payloads aislados, sin
integrarse en el flujo de ejecución del servidor; incluso el WAF de Durmuşkaya y Bayraklı \cite{durmuskaya2025waf}, la propuesta arquitectónicamente
más cercana entre las revisadas, opera como intermediario de red externo y no como componente interno del proceso Node.js. Esto significa que ninguno de estos
sistemas ha sido validado bajo restricciones de latencia de producción con acceso al contexto
interno completo de una petición HTTP.

La segunda es la \textbf{brecha de cobertura multivector con enfoque unificado dentro del propio
proceso de la aplicación}. Aunque Durmuşkaya y Bayraklı \cite{durmuskaya2025waf} sí cubren varios vectores de ataque (SQLi,
XSS, inyección de comandos y LFI) en un mismo sistema, lo hacen desde un proxy inverso externo,
entrenado con CSIC 2010 \cite{torrano2010csic}, HTTPParams 2015 y un conjunto propio, y evaluado contra DVWA, una
aplicación de formularios, y no contra APIs REST modernas que transmiten payloads en cuerpos JSON en
lugar de parámetros de URL. El resto de
los trabajos académicos revisados aborda un único vector por sistema.

La tercera es la \textbf{brecha de viabilidad operacional para organizaciones sin equipo de
seguridad}, que es simultáneamente un vacío académico y comercial. Los trabajos académicos
revisados reportan métricas de rendimiento del modelo pero no reportan impacto en latencia del
servidor anfitrión ni consumo de memoria del proceso; las soluciones comerciales que sí resuelven
estos aspectos operacionalmente lo hacen a un costo de licencia y de dependencia de
infraestructura externa incompatible con el perfil de una startup guatemalteca. SAP node-rasp
\cite{saparchive2019} es el único antecedente que intentó específicamente la protección en
runtime de Node.js sin dependencias externas, pero su enfoque de modificar el binario del runtime
lo hizo incompatible con entornos de producción estándar y fue abandonado sin soporte desde
versiones posteriores a Node.js LTS 10.

\section{Síntesis y relación con la propuesta de investigación}
Los antecedentes revisados confirman que el problema abordado por este trabajo no es la ausencia
de investigación en detección de ataques web mediante aprendizaje automático, tema que cuenta
con una literatura considerable y en expansión activa, como lo demuestra la publicación de Durmuşkaya y Bayraklı \cite{durmuskaya2025waf} en 2025. El vacío real
es la ausencia de una solución que combine simultáneamente cuatro condiciones que hasta ahora
solo han sido resueltas de forma parcial y nunca en conjunto: integración nativa en el runtime de
la aplicación en lugar de un componente de red externo, cobertura de múltiples vectores de ataque
en un único pipeline de detección, gratuidad y distribución mediante un mecanismo trivial de
instalación (\texttt{npm install}), y evaluación explícita del impacto en latencia y consumo de recursos
como criterio de aceptación del sistema. A esto se suma el vacío de contexto local: en ausencia de
un marco regulatorio vigente en Guatemala \cite{lahora2026estancada} y de herramientas
adaptadas a la realidad operativa de las startups tecnológicas del país, la adopción de controles de
seguridad depende enteramente de que existan opciones accesibles, lo cual no ha sido el caso
hasta ahora. \textit{logSguarDian} se posiciona precisamente en la intersección de estos vacíos: no
compite con las soluciones RASP empresariales ni pretende sustituir a los WAF perimetrales
existentes, sino que ocupa el espacio que ninguno de los antecedentes revisados cubre
completamente, proporcionando un mecanismo de detección multivector, integrado, gratuito y
operable sin equipo de seguridad dedicado.

\ifdefined\modoensamble
\else
	\section*{Referencias}
	\printbibliography[heading=none]
	\end{document}
\fi

\fi

\ifdefined\CAPjustificacion
	\newpage
	% ==============================================================================
% JUSTIFICACIÓN
% ==============================================================================
% Contenido redactado y provisto directamente por el estudiante (no adaptado
% del protocolo; reestructura las subsecciones 3.6-3.9 del protocolo
% "Utilidad para..." en una organización más estándar de tesis: pertinencia,
% relevancia técnica, relevancia social/económica, relevancia institucional,
% beneficiarios). No se modificó el contenido, solo se revisó.
%
% NOTA: el protocolo incluía una subsección explícita "3.7 Utilidad para el
% estudiante" que este capítulo no reproduce como sección propia. Puede ser
% una omisión involuntaria o una decisión editorial (es una categoría de
% justificación más débil/subjetiva que las demás); no se asume ninguna de
% las dos, se deja señalado.
%
% Verificación de fuentes (2026-08-08): se confirmó por WebFetch que
% soy502fortinet2026 respalda literalmente la cifra de 726 millones de
% intentos de ataque citada (dato de 2025). No se pudo verificar
% mastercard2025 (la URL devuelve 403 Forbidden); la cifra de 200% de
% incremento y el ranking "top 10 en Latinoamérica" quedan sin confirmar por
% este medio y deben revisarse manualmente antes de la entrega final.
% Actualización 2026-10-09: zeng2024 se reemplazó por zhen2025anomaly (mismo
% documento IEEE, autores y año corregidos). Se retiró la cita a verma2020 de la
% afirmación sobre el costo de los SIEM: ese artículo es una revisión de IDS
% para IoT y no respalda la afirmación. Resuelto el mismo día: la oración se
% reformuló sobre enisa2021smes (verificado en el PDF: págs. impresas 16, 17 y
% 18); se retiró la afirmación sobre "Machine Learning integrado en SIEM" e
% "infraestructura de análisis masivo", que no tiene fuente verificada.
% Actualización 2026-10-10: se acotó la promesa de entrega del dataset a scripts y
% vectores derivados (dos fuentes, XSS_dataset.csv y command injection.csv, no
% tienen licencia especificada por su origen). Estado verificado ese día: los
% archivos de vectores (training/data_clean/*.parquet|csv, training/splits/*.parquet)
% están en .gitignore y NO están publicados en el repositorio; ver la
% recomendación correspondiente en n-recomendaciones.tex.
\ifdefined\modoensamble
\else
	\input{0-preambulo_standalone}
\fi

\chapter{Justificación}

\section{Importancia y pertinencia del problema}
Como se estableció en el capítulo de Antecedentes, la carencia identificada no es la ausencia de
investigación sobre detección de ataques web mediante aprendizaje automático, sino la ausencia
de una solución que integre esa detección de forma nativa en el runtime de aplicaciones Node.js,
cubra múltiples vectores de ataque en un único sistema y sea operable por organizaciones sin
personal de seguridad dedicado. Justificar este trabajo exige, entonces, explicar por qué esa
carencia específica importa y para quién.

Las aplicaciones web construidas sobre Node.js constituyen una superficie de ataque
particularmente expuesta: el ecosistema npm se caracteriza por cadenas extensas de dependencias transitivas: un estudio
empírico publicado en 2022 reportó que el 42.1~\% de los paquetes analizados tenía al menos una
dependencia directa o indirecta vulnerable en su última versión \cite{zerouali2022impact}, y otro de
2019 concluyó que paquetes individuales pueden afectar a gran parte del ecosistema y que un número
muy reducido de cuentas de mantenedores podría usarse para inyectar código malicioso en la mayoría
de los paquetes \cite{zimmermann2019small}. Además, la ausencia de controles a nivel de aplicación
deja expuesta la lógica de negocio incluso cuando existen defensas perimetrales. Los mecanismos de defensa tradicionales, sea un Web Application Firewall (WAF)
o un Sistema de Detección de Intrusiones (IDS), fueron diseñados para operar en el borde de la red
y dependen de firmas estáticas; esta dependencia produce una tasa de falsos negativos
estructuralmente alta frente a variantes ligeramente modificadas de un ataque conocido o frente a
abuso de la lógica de negocio de la aplicación, categorías que no son detectables por coincidencia de patrones. Esta limitación es coherente con la
comparación experimental de Sureda Riera et al., en la que las herramientas RASP evaluadas superaron
a los WAF \cite{bermejohiguera2022}. Este no es un problema exclusivamente técnico: es un problema de
acceso. Las organizaciones que sí pueden mitigar esta limitación suelen hacerlo mediante plataformas de
gestión de eventos e información de seguridad (SIEM). Para las pequeñas y medianas empresas, sin
embargo, estos sistemas representan inversiones considerables, y muchas soluciones de
ciberseguridad requieren conocimiento especializado para implementarse y administrarse; en el
estudio de la Agencia de la Unión Europea para la Ciberseguridad sobre pymes europeas, solo el
28~\% de los participantes de la encuesta declaró haber asignado el rol de responsable de seguridad
de la información \cite{enisa2021smes}. Aunque ese estudio no se refiere a Guatemala, describe las
mismas dos barreras, presupuesto y personal especializado, que enfrentan las organizaciones que
constituyen la población objetivo de este trabajo (Sección~\ref{sec:just-contexto}). El problema, por tanto, no se limita a la existencia de ataques contra
aplicaciones Node.js, sino a la distribución desigual de las herramientas capaces de detectarlos.

\section{Relevancia técnica y científica}
Desde el punto de vista técnico, este trabajo es pertinente porque aborda una limitación
metodológica recurrente en la literatura de detección de ataques web mediante aprendizaje
automático: la evaluación de modelos sobre payloads o cadenas de texto aisladas, sin considerar
las restricciones de un entorno de ejecución real. Los sistemas de detección basados
exclusivamente en reglas colapsan frente a cargas útiles polimórficas o sintaxis de ataque
ligeramente ofuscada que evade patrones conocidos, mientras que los enfoques de aprendizaje
automático permiten generalizar el comportamiento base del servidor y detectar desviaciones
estadísticas sin depender de actualizaciones manuales de firmas \cite{zhen2025anomaly}. La combinación
de un componente supervisado, orientado a clasificar con precisión los vectores de ataque
conocidos, con un componente no supervisado orientado a identificar anomalías estadísticas
asociadas a variantes no observadas durante el entrenamiento, se propone responder a esta limitación sin comprometer la capacidad de generalización del sistema frente a amenazas nuevas; en qué medida lo logra se evalúa en los capítulos de Resultados y de Análisis de resultados.

La relevancia científica del proyecto radica en que el problema de integrar inferencia de
aprendizaje automático dentro de un proceso Node.js sin degradar su rendimiento no está resuelto
de forma satisfactoria en la literatura revisada, precisamente porque la mayoría de los sistemas
comparables opera como componente externo (proxy inverso o plataforma SIEM) y no como
extensión del propio proceso de la aplicación. Ejecutar inferencia dentro del ciclo de vida de una
petición HTTP exige resolver un reto de ingeniería específico de las arquitecturas orientadas a
eventos: evitar que el cómputo del modelo bloquee el bucle de eventos (\textit{event loop}) del cual
depende toda la capacidad de manejar peticiones concurrentes de Node.js. La construcción de un
pipeline de extracción de características e inferencia que opere de forma asíncrona sobre datos de
naturaleza multidimensional, sin bloquear dicho bucle, constituye un aporte técnico verificable
independientemente del desempeño final del modelo de clasificación, y es generalizable a
cualquier otro problema de seguridad en tiempo de ejecución sobre arquitecturas orientadas a
eventos.

\section{Relevancia social y económica en el contexto guatemalteco}
\label{sec:just-contexto}
La pertinencia de este trabajo no puede argumentarse en abstracto: depende directamente del
contexto en que se desarrollaría la solución. En Guatemala, los ataques cibernéticos registrados
durante 2024 aumentaron un 200~\% respecto del año anterior, cifra que posicionó al país entre los
diez de mayor exposición cibernética en Latinoamérica según el Cyber Insights Report de
Mastercard \cite{mastercard2025}. Reportes posteriores muestran que los incidentes continuaron: durante 2026 se reportaron ataques
cibernéticos contra instituciones públicas guatemaltecas, entre ellos el ataque a la Dirección
General de Control de Armas y Municiones (Digecam) ocurrido el 7 de abril \cite{emisoras2026digecam}
y el ataque cibernético confirmado por la Universidad de San Carlos de Guatemala
\cite{emisoras2026usac}, en un país que, además, no cuenta todavía con un marco legal de
ciberseguridad vigente que module la respuesta institucional a estos incidentes
\cite{lahora2026estancada}.

Este panorama de riesgo creciente contrasta con la estructura del tejido empresarial
guatemalteco. Según cifras del Ministerio de Economía, las micro, pequeñas y medianas empresas
representaron el 99.57~\% de las empresas activas con ventas en el año base 2015
\cite{mineco2015mipyme}, y una nota de la Agencia Guatemalteca de Noticias de julio de 2025 afirma
que representan el 96.3~\% de las empresas guatemaltecas \cite{agn2025mipymes}, y
es razonable asumir que una proporción sustancial de las empresas de base tecnológica que
desarrollan servicios sobre Node.js pertenece a este mismo segmento, es decir, opera sin un
Centro de Operaciones de Seguridad ni presupuesto destinado a licencias de plataformas
empresariales de detección. Nótese que esta es una inferencia razonable a partir de la estructura
general del tejido empresarial del país, y no una cifra específica y verificada sobre el segmento de
startups tecnológicas, dato que no está disponible públicamente de forma desagregada; el
argumento de relevancia social se sostiene sobre la composición general del sector empresarial
guatemalteco, no sobre una estadística exclusiva del sector de software que actualmente no existe.
En ese contexto, una herramienta que se distribuye sin costo mediante un gestor de paquetes de
uso estándar y que no exige modificar la infraestructura de red reduce directamente las dos
barreras de adopción más frecuentemente documentadas para estas organizaciones: el presupuesto y la complejidad operativa \cite{enisa2021smes}.

\section{Relevancia institucional}
Para la Universidad del Valle de Guatemala, este trabajo se ubica en la intersección entre Ciencia
de Datos, Ciberseguridad e Inteligencia Artificial, tres áreas de demanda creciente en el mercado
tecnológico regional. La metodología documentada de construcción del corpus (auditoría de fuentes, etiquetado
determinista y partición verificada por hash) y los vectores de características derivados
constituyen un recurso académico que puede emplearse en cursos de seguridad informática y
aprendizaje automático dentro de la Facultad de Ingeniería,
independientemente de si el sistema resultante alcanza los criterios de aceptación definidos para
el proyecto. La alineación explícita del diseño con marcos de referencia reconocidos
internacionalmente, como OWASP Top 10 y MITRE ATT\&CK, permite además que la metodología
documentada sea replicada, contrastada o extendida por trabajos de graduación posteriores dentro
del mismo programa.

\section{Beneficiarios directos e indirectos}
Es posible distinguir con claridad entre los beneficiarios directos de este trabajo y sus beneficiarios
indirectos, distinción necesaria para no sobreestimar el alcance real de la investigación en esta
etapa.

Los beneficiarios directos son las organizaciones guatemaltecas que desarrollan aplicaciones
sobre Node.js sin equipo de seguridad dedicado y que podrían instalar la librería resultante como
un componente adicional de su arquitectura existente sin incurrir en costo de licenciamiento ni en
modificaciones estructurales de su infraestructura de red. Dentro de este grupo se incluyen
también iniciativas comunitarias, organizaciones no gubernamentales con presencia digital y
proyectos de gobierno abierto que comparten la misma restricción de recursos técnicos y
financieros que las startups tecnológicas.

Los beneficiarios indirectos son de tres tipos. Primero, la comunidad académica de la
Universidad del Valle de Guatemala, a través de la metodología documentada de construcción del corpus, de los vectores derivados y del precedente
metodológico que este trabajo establece para investigaciones futuras en la intersección de
ciberseguridad e inteligencia artificial aplicada. Segundo, otros investigadores y equipos de desarrollo de la región centroamericana que, a partir del
material publicado del proyecto, podrían reproducir la construcción del corpus desde las mismas
fuentes públicas y reutilizar los vectores de características derivados para entrenar o validar
sistemas de detección propios, en un contexto donde la investigación aplicada en ciberseguridad con
datos representativos de la región es escasa. El proyecto no redistribuye los datos crudos de las
fuentes cuya licencia no está especificada por su origen: pone a disposición los analizadores y
scripts de normalización, el archivo de mapeo de etiquetas, la auditoría de fuentes y los
vectores de características derivados, que contienen valores numéricos y no el texto de las
peticiones. La reproducción del corpus no es completa: los conjuntos sintéticos de Command
Injection, que aportan el 17.8~\% de las filas de esa clase en el conjunto de prueba, no cuentan
con un generador versionado. Tercero, y de forma
más amplia, el ecosistema tecnológico guatemalteco en su conjunto, en la medida en que la
adopción voluntaria de controles de seguridad accesibles contribuye a elevar el nivel general de
madurez en ciberseguridad del país mientras el marco regulatorio correspondiente continúa en
trámite legislativo \cite{lahora2026estancada}.

Es importante precisar que la magnitud real de estos beneficios, tanto en términos de tasa de
detección como de impacto en el rendimiento del servidor anfitrión, es precisamente lo que este
trabajo se propone medir y no puede darse por establecida antes de la fase de evaluación descrita
en la metodología. Esta justificación argumenta la pertinencia de investigar el problema, no
anticipa el desempeño del sistema que se construirá para abordarlo.

\ifdefined\modoensamble
\else
	\section*{Referencias}
	\printbibliography[heading=none]
	\end{document}
\fi

\fi

\ifdefined\CAPobjetivos
	\newpage
	% ==============================================================================
% OBJETIVOS
% ==============================================================================
% Contenido provisto directamente por el estudiante. El Objetivo específico
% 3 coincide verbatim con protocolo_sebastian_v1 (3).pdf, sección 4.1.
%
% Citas owasp2021 y mitreattack: el estudiante no adjuntó bibliografía esta
% vez. Se verificaron directamente contra las páginas oficiales por WebFetch
% antes de agregarlas al .bib (no se citaron de memoria):
%   - owasp2021: confirmado el título "A03:2021 -- Injection" y la URL
%     canónica en owasp.org.
%   - mitreattack: confirmado el nombre del framework y la organización
%     (MITRE); el sitio no expone un número de versión o fecha de
%     publicación único, se cita como framework de actualización continua.
%
% Corrección 2026-10-02: el Objetivo específico 1 afirmaba que los cuatro
% vectores caen bajo "OWASP Top 10 2021 (A03:2021)". Es incorrecto para
% Path Traversal/LFI, que está mapeado en A01:2021 -- Broken Access Control,
% no en A03:2021 -- Injection (verificado por WebFetch contra
% top10.owasp.org; ver corrección completa documentada en
% i-marco_teorico.tex). Se corrigió para distinguir ambas categorías en vez
% de agrupar los cuatro vectores bajo una sola.
\ifdefined\modoensamble
\else
	\input{0-preambulo_standalone}
\fi

\chapter{Objetivos}

\section{Objetivo general}
Diseñar e implementar una librería npm que integre un pipeline de ingeniería de características y
un modelo híbrido de aprendizaje automático capaz de detectar y clasificar amenazas conocidas y
anomalías estadísticas correspondientes a variantes desconocidas de los vectores cubiertos en
endpoints Node.js en tiempo de ejecución, con el propósito de proveer a startups tecnológicas sin
equipo de seguridad dedicado una herramienta accesible de protección a nivel de capa de
aplicación que no comprometa el rendimiento del servidor anfitrión.

\section{Objetivos específicos}
\begin{enumerate}
    \item Evaluar cuatro vectores de ataque intra-request dirigidos a endpoints Node.js — SQL
    Injection, XSS, Path Traversal/LFI y Command Injection — conforme a su prevalencia en el
    OWASP Top 10 2021 \cite{owasp2021} (SQL Injection, XSS y Command Injection bajo A03:2021 --
    Injection; Path Traversal/LFI bajo A01:2021 -- Broken Access Control) y su cobertura en la
    taxonomía MITRE ATT\&CK \cite{mitreattack}, para establecer las categorías de amenaza que
    estructurarán el dataset de entrenamiento del OE2 y los criterios de evaluación del OE3.

    \item Construir un entorno de prueba controlado y reproducible que genere un dataset
    estructurado con un mínimo de 100{,}000 muestras de tráfico HTTP etiquetadas; balanceando
    tráfico legítimo y ataques de las categorías seleccionadas, documentando los scripts de
    generación y la metodología de etiquetado, a fin de disponer de una base de datos de
    entrenamiento y validación que garantice la representatividad estadística necesaria para el
    entrenamiento de los modelos de detección.

    \item Evaluar la efectividad del sistema de detección implementado mediante precisión, recall
    y F1-score por categoría de ataque, verificando como criterio mínimo de aceptación un
    F1-score $\geq 0.80$ en al menos tres de las cuatro categorías seleccionadas bajo condiciones
    de tráfico simulado, umbral establecido en función del balance operacionalmente tolerable
    entre falsos negativos y falsos positivos en entornos de producción sin analista de seguridad
    dedicado.
\end{enumerate}

\ifdefined\modoensamble
\else
	\section*{Referencias}
	\printbibliography[heading=none]
	\end{document}
\fi

\fi

\ifdefined\CAPalcance
	\newpage
	% ==============================================================================
% ALCANCE
% ==============================================================================
% Contenido provisto directamente por el estudiante. Dos correcciones
% aplicadas antes de incorporarlo:
%   1. "la sección 1.3" (referencia manual, en Alcance operativo) se
%      reemplazó por \label/\ref sobre la subsección "Alcance operativo":
%      con Alcance como capítulo propio, esa subsección será algo como
%      "5.1.3" según dónde quede el capítulo, no "1.3" (número que parece
%      arrastrado de una numeración distinta, probablemente la del
%      protocolo original).
%   2. "la sección de factibilidad de la metodología de este protocolo" se
%      corrigió a "de este documento": este texto ya es parte de la tesis,
%      no del protocolo.
%
% NOTA ESTRUCTURAL (sin resolver): "Alcance" no aparece en la lista de 10
% capítulos de .claude/CLAUDE.md (Introducción...Recomendaciones). Sí existe
% como capítulo opcional de la plantilla UVG (\CAPalcance). Se incluye por
% instrucción explícita del estudiante; el índice final tendrá 11 capítulos
% en vez de 10 si esto se mantiene así.
%
% Corrección 2026-10-02 (criterio de latencia del middleware completo, en
% "Alcance técnico"): el criterio original fijaba un incremento relativo de
% p95 del 5% bajo carga normal y 10% bajo carga de ataque simulado. Diego
% Valenzuela (componente de IA del mismo proyecto integrador) señaló por
% chat (2026-10-02) que un criterio relativo es débil cuando la línea base
% de latencia del servidor es de pocos milisegundos: un 5% de esa base es
% una fracción de milisegundo, difícil de medir de forma confiable y poco
% significativa como criterio de aceptación. Reportó que a él le funcionó
% un límite absoluto de 15~ms en su propio componente, y recomendó usar un
% número fijo en ms en vez de un porcentaje salvo que la línea base ya sea
% grande (> 100~ms). El estudiante confirmó adoptar un único umbral
% absoluto de 15~ms (sin distinguir carga normal de carga de ataque
% simulado) para el middleware completo, reemplazando el 5%/10% relativo.
% El criterio de 5~ms del pipeline de extracción de características
% (aislado, dentro del worker thread) no cambia: ya era un valor absoluto.
\ifdefined\modoensamble
\else
	\input{0-preambulo_standalone}
\fi

\chapter{Alcance}

\section{Alcance del proyecto}

\subsection{Alcance funcional}
El proyecto comprende el diseño, implementación y evaluación de un middleware distribuible
como paquete npm para aplicaciones Node.js basadas en Express, capaz de detectar y responder,
en tiempo de ejecución, a peticiones HTTP correspondientes a cuatro vectores de ataque
intra-request: SQL Injection, Cross-Site Scripting (XSS), Path Traversal/LFI e inyección de
comandos. La respuesta del sistema ante una petición clasificada como maliciosa se limita a dos
mecanismos verificables: el bloqueo de la petición mediante código de respuesta HTTP 403 y el
registro de una alerta. El proyecto entrega, como producto funcional, el paquete npm publicado, los
scripts de construcción del corpus etiquetado y los vectores de características derivados, y la
documentación técnica correspondiente. No se entregan los datos crudos de las fuentes públicas
empleadas: cada fuente se obtiene de su origen conforme a su licencia.

\subsection{Alcance técnico}
\label{subsec:alcance-tecnico}
El sistema se integra exclusivamente como middleware dentro del pipeline de peticiones de
Express.js, sin requerir modificaciones al runtime de Node.js ni infraestructura de red externa
(proxy inverso o plataforma SaaS de terceros). El componente de inferencia consume un modelo
serializado en formato ONNX mediante \texttt{onnxruntime-node}, ejecutado dentro de un
\texttt{worker\_thread} independiente para evitar el bloqueo del Event Loop principal. El alcance
técnico está delimitado por dos restricciones cuantitativas, definidas como criterios de aceptación
verificables: el pipeline de extracción de características debe ejecutarse en un tiempo igual o
menor a 5~ms por petición, medido de forma aislada dentro de su worker thread; y el middleware
completo no debe incrementar la latencia p95 del servidor anfitrión en más de 15~ms absolutos
respecto de la línea base sin middleware activo, tanto bajo carga normal como bajo carga de
ataque simulado. Se adopta un umbral absoluto en lugar de uno relativo porque, en un servidor
cuya línea base de latencia es de pocos milisegundos, un porcentaje fijo representaría una
fracción de milisegundo: un criterio estadísticamente frágil y difícil de verificar de forma
confiable bajo las condiciones de medición de este proyecto.

\subsection{Alcance operativo}
\label{subsec:alcance-operativo}
La validación principal del sistema se realiza en un entorno de simulación controlado y
reproducible, construido mediante Docker Compose, con generación de tráfico legítimo mediante
Artillery y de tráfico malicioso mediante herramientas especializadas por vector de ataque
(sqlmap, OWASP ZAP y scripts propios para Path Traversal/LFI y Command Injection). La
validación extendida con entornos productivos de startups guatemaltecas, contemplada como
escenario complementario en la Fase 4 del cronograma, queda condicionada a la disponibilidad de
un acuerdo de colaboración con dichas organizaciones y no constituye el criterio de aceptación
principal del proyecto.

\section{Delimitación}

\subsection{Lo que comprende este documento}
Este documento delimita específicamente el componente de ciberseguridad del sistema
\textit{logSguarDian}: el análisis y selección de vectores de ataque, la construcción del entorno de
simulación y el dataset etiquetado, la implementación inicial del extractor de características y el paso a \texttt{worker\_threads}, la capa de captura de peticiones HTTP y el sistema de
respuesta en tiempo de ejecución, y la evaluación de cobertura de detección mediante pruebas de
penetración controlada. La integración de estos componentes con el modelo de aprendizaje
automático, así como las pruebas de integración extremo a extremo y la documentación final
conjunta, son responsabilidad compartida de ambos integrantes del proyecto y se documentan en
la sección de metodología correspondiente.

\subsection{Lo que queda fuera del alcance}
\label{subsec:fuera-alcance}
Explícitamente quedan fuera del alcance de este trabajo:

\begin{itemize}
    \item El diseño, entrenamiento, selección de arquitectura y serialización del modelo de
    aprendizaje automático, documentado de forma independiente en el trabajo de graduación de
    Diego Valenzuela, correspondiente al componente de Inteligencia Artificial.
    \item La detección de vectores de ataque distintos a los cuatro delimitados en el Objetivo
    Específico 1 (SQL Injection, XSS, Path Traversal/LFI e inyección de comandos). Categorías
    como Server-Side Request Forgery (SSRF), Cross-Site Request Forgery (CSRF) o ataques de
    denegación de servicio a nivel de aplicación no forman parte del sistema de detección
    propuesto.
    \item El soporte para runtimes distintos a Node.js o frameworks distintos a Express.js. El
    sistema no está diseñado para operar sobre aplicaciones construidas en Django, Spring,
    Laravel u otros stacks tecnológicos.
    \item El reemplazo de mecanismos perimetrales de seguridad existentes. El middleware se
    plantea como una capa complementaria a nivel de aplicación, no como sustituto de un
    Web Application Firewall, un IDS de red o un SIEM empresarial.
    \item El cumplimiento normativo o certificación bajo cualquier marco legal de ciberseguridad
    vigente o en trámite en Guatemala. El proyecto no garantiza ni evalúa conformidad legal;
    su relevancia frente al panorama regulatorio se documenta como antecedente contextual,
    no como criterio de aceptación.
    \item El soporte, mantenimiento o actualización del paquete npm posterior a la fecha de entrega del informe final en noviembre de 2026.
    \item El despliegue en infraestructura de producción de terceros fuera del entorno de
    validación controlado y del escenario complementario descrito en la Sección~\ref{subsec:alcance-operativo},
    cuya realización depende de un acuerdo de colaboración externo al proyecto.
\end{itemize}

\section{Restricciones y supuestos}
El desarrollo del proyecto está sujeto a las siguientes restricciones y supuestos, verificables al
momento de redacción de este documento:

\begin{itemize}
    \item \textbf{Restricción temporal.} El proyecto se ejecuta entre mayo y octubre de 2026,
    organizado en seis fases con dependencias secuenciales explícitas; la Fase 3, correspondiente
    al desarrollo del middleware, depende de la finalización de la Fase 1 (dataset) y la Fase 2
    (entrenamiento del modelo, a cargo del componente de IA).
    \item \textbf{Restricción presupuestaria.} El proyecto opera con un presupuesto de USD 0 en
    condiciones normales, con un techo de USD 50 en el escenario de despliegue cloud temporal
    para validación externa. Esta restricción excluye explícitamente el uso de plataformas de
    seguridad comerciales con licenciamiento de pago como parte del sistema desarrollado.
    \item \textbf{Restricción de recursos humanos.} El componente de ciberseguridad es
    responsabilidad de un único integrante, con dedicación de 20 horas semanales durante las
    Fases 1 a 3 y 15 horas semanales durante las Fases 4 a 6, sin personal externo adicional
    disponible.
    \item \textbf{Supuesto de disponibilidad de dependencias.} Se asume la disponibilidad
    continua de las herramientas de código abierto empleadas (Node.js, Express, Docker,
    sqlmap, OWASP ZAP, Jest) y del dataset público de referencia (OWASP ModSecurity Dataset,
    Zenodo) durante toda la duración del proyecto; el riesgo de interrupción del acceso al
    dataset se mitiga mediante copia local descargada al inicio de la Fase 1.
    \item \textbf{Supuesto de estabilidad del entorno de simulación.} Se asume que el entorno
    Docker Compose reproduce condiciones de tráfico suficientemente representativas del
    comportamiento de una API REST en producción para que los resultados obtenidos en el
    entorno controlado sean generalizables, supuesto que la validación complementaria con
    startups guatemaltecas busca contrastar sin condicionar a ella el cumplimiento del proyecto.
\end{itemize}

\section{Viabilidad del alcance propuesto}
La viabilidad del alcance definido descansa en tres condiciones verificables. Primero, la totalidad
del stack tecnológico requerido (Node.js, Express, onnxruntime-node, Docker) está disponible,
mantenido activamente y documentado para los casos de uso específicos del proyecto, sin
dependencia de tecnología en desarrollo o no probada. Segundo, el dataset base de referencia está
publicado bajo licencia abierta y ha sido empleado previamente en investigaciones de detección de
intrusiones en aplicaciones web, lo que reduce el riesgo de que la construcción del entorno de
simulación se convierta en un cuello de botella del cronograma. Tercero, el alcance funcional se
limita deliberadamente a cuatro vectores de ataque y a dos mecanismos de respuesta, en lugar de
plantear una cobertura exhaustiva de OWASP Top 10, precisamente para mantener el proyecto
ejecutable dentro de las 20 semanas efectivas de desarrollo disponibles y de los recursos humanos
descritos en la sección anterior. Los riesgos identificados que podrían afectar el cumplimiento del
alcance definido, junto con sus planes de mitigación, se documentan de forma exhaustiva en la
sección de factibilidad del Capítulo de Metodología de este documento.

\ifdefined\modoensamble
\else
	\end{document}
\fi

\fi

\ifdefined\CAPmarcoteorico
	\newpage
	% ==============================================================================
% MARCO TEÓRICO
% ==============================================================================
% Contenido provisto directamente por el estudiante (dos veces: la primera
% entrega se cortó a los 50,000 caracteres de por medio de la bibliografía;
% el reenvío completo llegó con la codificación corrupta -- "ó" en vez de
% "ó", etc. -- así que este archivo se reconstruyó a partir del texto de la
% primera entrega, que sí tenía acentos correctos, y se completó con las
% entradas bibliográficas del reenvío).
%
% Corrección aplicada: "la restricción de disponibilidad descrita en la
% sección 4.1" (en Inferencia mediante ONNX) era una referencia numérica
% hardcodeada a la subsección "Alcance técnico" de h-alcance.tex. Con
% Alcance como capítulo propio, esa subsección no será "4.1" (mismo tipo de
% problema que "la sección 1.3" en Alcance). Se agregó
% \label{subsec:alcance-tecnico} en h-alcance.tex y aquí se usa \ref.
%
% Verificación de fuentes (2026-08-15): CVE-2021-27185 confirmado por Snyk
% (command injection, paquete npm samba-client) -- coincide con el texto;
% no se pudo verificar independientemente el puntaje CVSS 9.8 exacto.
% republicamastercard2025 (200% de incremento; tecnológico 7.9%, tercer
% sector más atacado) confirmado literal. gem2024 (segunda tasa de
% actividad emprendedora temprana del mundo, 32.4%) NO se pudo verificar
% (la URL del PDF responde 403); revisar manualmente antes de la entrega.
%
% NOTA: la entrada ecosistemastartup2025 llegó en el .bib del estudiante
% pero no se cita en ningún punto de este capítulo; se agregó al .bib de
% todas formas (no genera error, biblatex solo imprime lo citado) pero
% queda sin usar hasta que se referencie en algún capítulo.
%
% Actualización 2026-09-04: el estudiante volvió a pegar el envío original
% completo (misma codificación corrupta "ó" documentada arriba). Se revisó
% contra el archivo ya reconstruido y NO se reincorporó la mayoría del
% contenido: la cifra "Node.js >40% de uso entre desarrolladores" no tiene
% cita ni en la versión original; las fuentes de blog personal (w3schools,
% Medium: cocca2025/ejaz2025/rusin2026/kalita2024, browserstack2025,
% codeb2024, socket2024, mdnwebdocs) se descartaron por calidad frente a las
% fuentes ya usadas (OpenJS Foundation, expressjs); checkpoint2023 no tiene
% URL verificable; la sección "Entorno de Prueba y Reproducibilidad"
% (Docker/Artillery/DotDotPwn/Commix) es diseño experimental y se dejó
% pendiente para el capítulo de Metodología en vez de duplicarla aquí.
%
% Se intentó agregar una sección "Trabajos previos y posicionamiento" con
% cuadro comparativo de trabajos ML previos, pero se detectó que
% e-antecedentes.tex YA tiene una sección equivalente ("Trabajos académicos
% en detección de ataques web mediante aprendizaje automático",
% \label{cuadro:trabajosprevios}) con los mismos cinco trabajos. Por
% instrucción explícita del estudiante, ese contenido se queda únicamente
% en Antecedentes y no se repite aquí.
%
% Por la misma lógica, se eliminó la subsección "RASP comerciales
% existentes" (Cuadro cuadro:rasp-comercial) de este archivo: duplicaba
% casi por completo la sección "Soluciones RASP existentes: antecedentes
% tecnológicos" (Cuadro cuadro:rasp, ahora \label{sec:rasp-antecedentes})
% de Antecedentes. Se reemplazó por una oración que remite a esa
% comparación en vez de repetirla, y se agregó el label faltante en
% e-antecedentes.tex para poder referenciarla. También se dejó fuera de
% este archivo la sección "Contexto Tecnológico y de Ciberseguridad en
% Guatemala" completa del envío del estudiante (GEM, Mastercard, MINECO,
% BID-OEA): esas mismas cifras y fuentes ya se citaron en el capítulo de
% Justificación; este archivo mantiene solo la versión condensada que ya
% tenía, que remite a Antecedentes/Justificación para el detalle.
%
% Contenido nuevo sí incorporado del envío del estudiante (2026-09-04),
% verificado antes de escribirse: cita faltante de microsoft2009stride en
% STRIDE; mención de Helmet/CORS/RateLimit/JWT/bcrypt en la subsección de
% middleware; subsección nueva "Defensa en profundidad" con figura TikZ
% (fuente cloudflaredefenseindepth NO verificable por WebFetch, 403,
% marcado en el .bib); cuarta limitación de RASP (complejidad de
% actualización del modelo embebido); indicadores de log adicionales para
% SQLi/XSS/Path Traversal/Command Injection (mismas fuentes ya verificadas,
% sin citas nuevas); referencia cruzada OWASP API3:2023 -> Path
% Traversal/LFI.
%
% Explícitamente NO incorporado por falta de fuente o por baja calidad de
% fuente: la cifra "Node.js >40% de uso entre desarrolladores" (sin cita ni
% en el envío original); las estadísticas de escala de npm de socket2024
% (2.5M paquetes, 184 mil millones de descargas/mes, 165,486 mantenedores:
% no se pudo confirmar el contenido específico por WebFetch); las cifras de
% Akamai (311/150 mil millones de ataques, +33%) vía encryptionconsulting2025
% (URL devuelve 404) y del Enterprise Strategy Group (62%) vía equixly2024
% (la portada del blog no confirma la cifra); checkpoint2023 (sin URL) y
% checkpointresearch2024 (la portada no muestra el contenido citado); y las
% fuentes de blog personal (w3schools, Medium: cocca2025/ejaz2025/
% rusin2026/kalita2024, browserstack2025, codeb2024, mdnwebdocs).
%
% Corrección 2026-10-02 (clasificación OWASP/MITRE, verificada por WebFetch
% contra las páginas oficiales antes de escribirse):
%  (1) Path Traversal (CWE-22/CWE-23) está mapeado en A01:2021 Broken Access
%      Control, NO en A03:2021 Injection -- confirmado contra
%      top10.owasp.org/2021/A01_2021-Broken_Access_Control. SQLi (CWE-89),
%      XSS (CWE-79) y OS Command Injection (CWE-78) sí están en A03:2021 --
%      confirmado contra top10.owasp.org/2021/A03_2021-Injection, que
%      además confirma que CWE-22 NO aparece en esa página. Se corrigieron
%      las subsecciones de cada vector, la sección "Marcos de referencia", y
%      se agregó el Cuadro cuadro:owasp-mitre-cwe.
%  (2) Se eliminó la afirmación (agregada por error en la sesión anterior,
%      2026-09-04) de que API3:2023 (Broken Object Property Level
%      Authorization) se corresponde con Path Traversal/LFI: verificado
%      contra api-security.owasp.org que API3:2023 cubre exposición
%      excesiva de propiedades y Mass Assignment, sin relación con
%      manipulación de rutas de archivo.
%  (3) MITRE: verificado contra attack.mitre.org que T1189 (Drive-by
%      Compromise) se centra explícitamente en comprometer el endpoint
%      cliente, con XSS como mecanismo de entrega -- se mantiene solo para
%      XSS. T1190 (Exploit Public-Facing Application) menciona
%      explícitamente SQL injection, path traversal y el OWASP Top 10/CWE
%      Top 25 -- se usa para SQLi, Path Traversal y Command Injection.
%  (4) Las consecuencias de SQLi se ajustaron a lo que A03:2021 documenta
%      (modificación/borrado de datos, invocación de procedimientos
%      almacenados) en vez de la lista genérica anterior.
%  Entradas .bib nuevas, cada una verificada por WebFetch antes de
%  escribirse: owaspa01_2021, owaspa03_2021, owaspapi3_2023, owasptraversal,
%  mitret1189, mitret1190.
%
% Propagación: el mismo error (los cuatro vectores bajo "A03:2021") estaba
% también en g-objetivos.tex (Objetivo específico 1); se corrigió ahí
% también por ser la misma afirmación factual incorrecta, no una decisión
% de alcance distinta.
% Actualización 2026-10-04: las figuras fig:defensa-profundidad y
% fig:rasp-vs-perimetral (antes TikZ) se reemplazaron por las imágenes
% provistas por el estudiante (figuras/figura2.png y figuras/figura3.png).
% El párrafo de defensa en profundidad se ajustó a las cinco capas de la
% figura nueva, y la leyenda de la figura 3 pasó de "arriba/abajo" a
% "izquierda/derecha".
\ifdefined\modoensamble
\else
	\input{0-preambulo_standalone}
\fi

\chapter{Marco teórico}

\section{Node.js y arquitectura orientada a eventos}

\subsection{Modelo de ejecución \textit{event-driven}}
Node.js es un entorno de ejecución de JavaScript, asíncrono y orientado a eventos, que ejecuta el
motor V8 de Google Chrome fuera del navegador \cite{openjsfoundation2024,nodejs2026intro}. A diferencia de
los servidores tradicionales que asignan un hilo por conexión entrante, Node.js opera con un único hilo de ejecución de JavaScript organizado en torno a un \textit{Event
Loop}, que permite realizar operaciones de entrada y salida sin bloqueo al delegarlas al kernel del
sistema siempre que es posible \cite{nodejs2026eventloop}. Según el modelo de ejecución de
JavaScript, cuando la pila de llamadas (\textit{Call Stack}) queda vacía se toma el siguiente
trabajo pendiente de la cola para ejecutarlo \cite{mdn2026execution}.

Este modelo es la razón técnica por la que Node.js es la plataforma elegida para
\textit{logSguarDian}: su capacidad de atender miles de conexiones concurrentes sin bloquear el
hilo principal permite instrumentar peticiones HTTP en tiempo real sin degradar la capacidad de
respuesta del servidor, siempre que el propio mecanismo de detección respete la naturaleza no
bloqueante del entorno. Esta restricción de diseño es la que obliga a que la inferencia del modelo
de detección se ejecute fuera del hilo principal, mediante un \texttt{worker\_thread} independiente,
tal como se detalla en la sección \ref{sec:onnx-inferencia}.

\subsection{Middleware y el pipeline de peticiones en Express}
El middleware es el mecanismo mediante el cual \textit{logSguarDian} se integra en aplicaciones
Node.js existentes. Una función middleware en Express tiene acceso al objeto de petición
(\texttt{req}), al objeto de respuesta (\texttt{res}) y a la siguiente función en la cadena; puede
modificar ambos objetos, terminar el ciclo petición-respuesta, o ceder el control a la siguiente
función mediante \texttt{next()} \cite{expressjsguide}. Cada función middleware se ejecuta en el
orden en que fue registrada, formando un pipeline en el que la petición y la respuesta atraviesan
cada capa antes de llegar al manejador de ruta final, que cierra el ciclo enviando la respuesta al
cliente.

Este mecanismo es lo que hace posible que \textit{logSguarDian} se instale como una dependencia
adicional del proyecto sin requerir modificaciones a la infraestructura existente: al registrarse
como middleware, el sistema obtiene acceso al contenido completo de cada petición antes de que
esta llegue a la lógica de negocio de la aplicación, la posición óptima para interceptar payloads
maliciosos. La Figura \ref{fig:pipeline-middleware} ilustra este pipeline y la posición que ocupa
\textit{logSguarDian} dentro de él.

\begin{figure}[H]
\centering
\begin{tikzpicture}[
  node distance=0.5cm,
  box/.style={draw, rectangle, rounded corners, minimum height=1cm, minimum width=2.6cm, align=center, font=\small},
  detect/.style={draw, rectangle, rounded corners, minimum height=1cm, minimum width=2.6cm, align=center, font=\small, fill=gray!15, thick},
  arr/.style={-{Stealth[length=2mm]}, thick}
]
\node[box] (req) {Petición HTTP\\entrante};
\node[box, right=of req] (helmet) {Middleware\\Helmet / CORS};
\node[detect, right=of helmet] (lsd) {\textit{logSguarDian}\\(captura + inferencia)};
\node[box, right=of lsd] (auth) {Middleware\\de autenticación};
\node[box, right=of auth] (route) {Route\\Handler};
\node[box, below=1.3cm of lsd, xshift=1.5cm] (res) {Respuesta al cliente};

\draw[arr] (req) -- (helmet);
\draw[arr] (helmet) -- (lsd);
\draw[arr] (lsd) -- (auth);
\draw[arr] (auth) -- (route);
\draw[arr] (route.south) -- ++(0,-0.4) -| (res.north);
\draw[arr, dashed] (lsd.south) -- ++(0,-0.6) -| node[below, font=\scriptsize, pos=0.25] {Bloqueo HTTP 403\\si se detecta amenaza} (res.north);
\end{tikzpicture}
\caption[Pipeline de middleware en Express.js]{Pipeline de middleware en Express.js y posición de \textit{logSguarDian} dentro del ciclo de vida de una petición}
\label{fig:pipeline-middleware}
\end{figure}
\vspace{-0.3cm}
{\footnotesize Nota. Elaboración propia, con base en la especificación de middleware de Express \cite{expressjsguide}.}

El ecosistema npm ofrece middleware especializado que se integra en este mismo pipeline para
reforzar la seguridad a nivel de aplicación: Helmet.js configura cabeceras HTTP de seguridad,
\texttt{cors} gestiona las políticas de origen cruzado, y \texttt{express-rate-limit} \cite{expressratelimit2026} controla la frecuencia de peticiones para
mitigar el abuso de la API, mientras que \texttt{jsonwebtoken} \cite{auth02026jsonwebtoken} y
\texttt{bcrypt} \cite{kelektiv2026bcrypt} cubren la gestión de tokens de autenticación y el hashing
de contraseñas, respectivamente. \textit{logSguarDian} se posiciona como una capa adicional
dentro de este mismo pipeline, complementaria a estos controles y no sustituta de ellos: mientras
Helmet y CORS operan sobre cabeceras y políticas de origen, y RateLimit sobre el volumen de
peticiones, \textit{logSguarDian} analiza el contenido semántico del cuerpo, los parámetros y el
comportamiento de cada petición para detectar los cuatro vectores de ataque intra-request
resumidos en el Cuadro~\ref{cuadro:vectores}.

\section{Ecosistema npm}
npm (Node Package Manager) es el gestor de paquetes oficial de Node.js: una herramienta de línea
de comandos que resuelve dependencias contra un registro centralizado de paquetes de código
abierto \cite{npmregistry2024}. La escala de este ecosistema es directamente relevante para el
proyecto en dos sentidos. Primero, define el canal de distribución de \textit{logSguarDian}: al
publicarse como paquete npm, la herramienta hereda un mecanismo de instalación estándar
(\texttt{npm install}) ya familiar para cualquier equipo que desarrolle sobre Node.js, eliminando la
fricción de adopción que sí presentan las soluciones RASP comerciales descritas en la
Sección~\ref{sec:rasp-existentes} de Antecedentes. Segundo, la resolución de dependencias mediante \texttt{package.json}
garantiza que las versiones de las herramientas empleadas en el proyecto (Express, Jest,
onnxruntime-node) queden fijadas y sean reproducibles entre entornos, condición necesaria para la
factibilidad técnica documentada en la metodología de este trabajo.

\section{Contexto de ciberseguridad en Guatemala}
\label{sec:contexto-guatemala}
El capítulo de Antecedentes de este informe documenta en detalle el panorama de amenazas
cibernéticas en Guatemala y el estado del marco regulatorio pendiente (Iniciativa de Ley 6347).
Esta sección retoma únicamente los elementos de ese contexto que son insumo directo para
decisiones de diseño de \textit{logSguarDian}.

La estructura del tejido empresarial guatemalteco es el dato que delimita el perfil de usuario
objetivo de la herramienta. Según datos del Ministerio de Economía, las micro, pequeñas y
medianas empresas representan entre el 96~\% y el 99~\% de las empresas del país, según la fuente y el año consultados
\cite{mineco2015mipyme,agn2025mipymes}, un segmento caracterizado por equipos de desarrollo pequeños y
presupuestos que rara vez contemplan una partida específica para ciberseguridad. El Reporte
Nacional GEM Guatemala 2023-2024 documenta además que el país posee la segunda tasa de
actividad emprendedora temprana más alta del mundo (32.4~\%) \cite{gem2024}, y que iniciativas
como EmprendamosGuate, coordinada por el Ministerio de Economía y PRONACOM, buscan identificar y acompañar a emprendedores con potencial de inversión \cite{pronacom}, un
ecosistema de emprendimiento por oportunidad caracterizado por el BID \cite{iadb2024}; este trabajo
asume, como supuesto de diseño y no como dato verificado, que una parte sustancial de las
aplicaciones de ese ecosistema se construye como APIs REST sobre Node.js.

Frente a este crecimiento, la exposición a amenazas ha aumentado de forma paralela. El Cyber
Insights Report Guatemala de Mastercard documenta un incremento del 200~\% en ciberataques
durante 2024 respecto al año anterior, con el sector tecnológico en tercer lugar (7.9~\%) entre los sectores más afectados, después del
financiero (19.3~\%) y el industrial (16.8~\%) \cite{mastercard2025,republicamastercard2025}. Esta combinación, alta tasa de creación de startups sobre Node.js
y crecimiento sostenido de la exposición a ataques, sin que exista aún un marco legal vigente que
incentive la inversión en controles mínimos \cite{lahora2026estancada}, es la que ancla la
justificación técnica de este proyecto en una realidad operativa concreta y no en una abstracción.

En cuanto a la madurez institucional en ciberseguridad, el informe conjunto del Banco
Interamericano de Desarrollo y la Organización de los Estados Americanos, elaborado con el
Global Cyber Security Capacity Centre de la Universidad de Oxford, sitúa a la región en un proceso
de mejora sostenida entre 2020 y 2025 en las cinco dimensiones del Modelo de Madurez de
Capacidad en Ciberseguridad (CMM), aunque persisten brechas estructurales en recursos, talento y
coordinación intersectorial \cite{bidoea2025}. Estas brechas de talento especializado son
precisamente la condición estructural que impide a las startups guatemaltecas contratar personal
dedicado a la seguridad informática, y son la que fundamenta el requisito de diseño de
\textit{logSguarDian} de operar de forma autónoma, sin necesidad de un analista que interprete
alertas ni ajuste reglas de forma continua.

\section{Fundamentos de seguridad en aplicaciones web}

\subsection{La tríada CIA}
La seguridad de la información se articula convencionalmente en torno a tres propiedades:
confidencialidad, integridad y disponibilidad, preservar las restricciones autorizadas sobre el acceso y la divulgación de la información, mantener
su exactitud y prevenir su modificación no autorizada, y
garantizar el acceso oportuno a los recursos por parte de usuarios legítimos, respectivamente \cite{nistinfosec2026}. En el contexto de una API REST construida sobre Node.js, cada uno de los cuatro
vectores de ataque cubiertos por \textit{logSguarDian} compromete predominantemente una de
estas propiedades: la inyección SQL y el Path Traversal/LFI comprometen la confidencialidad al
exponer datos no autorizados; el Cross-Site Scripting compromete la integridad al permitir la
manipulación de la sesión del usuario; y la inyección de comandos puede comprometer las tres
simultáneamente al otorgar ejecución arbitraria con los privilegios del proceso del servidor. La
disponibilidad, en particular, es una propiedad estructuralmente frágil en Node.js: al operar sobre
un único hilo de ejecución, una sobrecarga deliberada del Event Loop puede degradar o
interrumpir el servicio para la totalidad de los clientes concurrentes, razón por la cual el propio
diseño de \textit{logSguarDian} debe evitar convertirse él mismo en un vector de degradación de
disponibilidad, de ahí los criterios de latencia máxima definidos en el Alcance de este trabajo.

\subsection{Superficie de ataque en APIs REST}
La superficie de ataque de una aplicación web es el conjunto de puntos de entrada por los que un
agente malicioso puede introducir datos o extraer información del sistema. En una API REST, esta
superficie corresponde a cada endpoint HTTP expuesto junto con los métodos, parámetros,
cabeceras y cuerpos de solicitud que acepta. Las APIs presentan una superficie de ataque
estructuralmente mayor que las aplicaciones web tradicionales orientadas a la renderización de
páginas, porque cada endpoint es un punto de entrada independiente con su propio contrato de
datos, y porque la comunicación sin estado de REST implica que cada petición debe transportar
toda la información necesaria para su procesamiento, incluyendo credenciales y parámetros que un
atacante puede manipular directamente sin necesidad de interactuar con una interfaz gráfica
intermedia. Esta característica es la que justifica que el punto de instrumentación de
\textit{logSguarDian} sea el objeto \texttt{req} de cada petición entrante, y no el tráfico de red
agregado, dado que es en ese objeto donde se materializa la totalidad de la superficie de ataque
relevante para los cuatro vectores cubiertos.

\subsection{Defensa en profundidad}
La defensa en profundidad es una estrategia de ciberseguridad que implementa múltiples capas de
controles a lo largo de la infraestructura de un sistema, bajo la premisa de que ningún control
individual es infalible: ante la eventual evasión de una capa, las capas subsiguientes deben ser
capaces de detectar, contener o mitigar el avance del ataque \cite{cloudflaredefenseindepth}.
Aplicada al contexto de una API REST en Node.js, la estrategia organiza los controles en cinco
capas funcionales: en la capa de red operan los firewalls de capas 3 y 4 y los IDS de red; en la
capa perimetral, los WAF basados en firmas; en la capa de aplicación, los middlewares de
validación de entrada (Helmet, CORS) y los mecanismos de autenticación; en la capa de detección en
tiempo de ejecución, las herramientas RASP que inspeccionan la petición dentro del propio proceso;
y en la capa de datos, el cifrado en tránsito y en reposo y el control de acceso a la base de
datos. La Figura \ref{fig:defensa-profundidad} ubica a \textit{logSguarDian} dentro de este
modelo, en la capa de detección en tiempo de ejecución.

\begin{figure}[H]
\centering
\includegraphics[width=0.85\textwidth]{figura2.png}
\caption[Capas del modelo de defensa en profundidad]{Capas del modelo de defensa en profundidad aplicadas a una API REST en Node.js, con \textit{logSguarDian} en la capa de detección en tiempo de ejecución (RASP)}
\label{fig:defensa-profundidad}
\end{figure}
\vspace{-0.3cm}
{\footnotesize Nota. Elaboración propia, con base en Cloudflare \cite{cloudflaredefenseindepth}.}

Este modelo posiciona con precisión el aporte de \textit{logSguarDian}: no reemplaza los controles
de red ni los de aplicación ya existentes, sino que complementa la capa de detección y respuesta,
que en el perfil de startup guatemalteca sin equipo de seguridad dedicado descrito en el
capítulo de Antecedentes suele estar ausente por completo.

\subsection{Modelo de amenazas STRIDE aplicado al proyecto}
STRIDE es un modelo de clasificación de amenazas que organiza los riesgos de seguridad en seis
categorías: \textit{Spoofing} (suplantación de identidad), \textit{Tampering} (manipulación de
datos), \textit{Repudiation} (negación de una acción realizada), \textit{Information Disclosure}
(divulgación de información), \textit{Denial of Service} (denegación de servicio) y
\textit{Elevation of Privilege} (escalada de privilegios) \cite{microsoft2025stride}. El Cuadro \ref{cuadro:stride} mapea los
cuatro vectores de ataque cubiertos por \textit{logSguarDian} contra las categorías STRIDE que
comprometen directamente.

\begin{table}[H]
\centering
\small
\begin{tabular}{|L{3.4cm}|L{5.0cm}|L{4.6cm}|}
\hline
\textbf{Vector de ataque} & \textbf{Categoría STRIDE comprometida} & \textbf{Propiedad de la tríada CIA} \\ \hline
SQL Injection & Information Disclosure, Tampering & Confidencialidad, Integridad \\ \hline
Cross-Site Scripting & Spoofing, Information Disclosure & Confidencialidad, Integridad \\ \hline
Path Traversal / LFI & Information Disclosure & Confidencialidad \\ \hline
Command Injection & Tampering, Elevation of Privilege & Confidencialidad, Integridad, Disponibilidad \\ \hline
\end{tabular}
\caption[Mapeo de vectores de ataque a categorías STRIDE]{Mapeo de los cuatro vectores de ataque cubiertos por \textit{logSguarDian} a las categorías del modelo STRIDE}
\label{cuadro:stride}
\end{table}
{\footnotesize Nota. Elaboración propia con base en las categorías de STRIDE \cite{microsoft2025stride} y en la tríada de confidencialidad, integridad y disponibilidad \cite{nistinfosec2026}; la asignación de cada vector es criterio del autor.}

Dos categorías de STRIDE quedan explícitamente fuera de la cobertura de la primera versión de
\textit{logSguarDian}: \textit{Repudiation}, que requeriría un sistema de auditoría e integridad de
logs independiente del mecanismo de detección, y la cobertura completa de \textit{Elevation of
Privilege} más allá de la vía de Command Injection, que abarcaría vulnerabilidades de control de
acceso como BOLA (\textit{Broken Object Level Authorization}). Esta delimitación es consistente
con el Alcance definido en este informe y se documenta aquí como una decisión metodológica
explícita, no como una omisión no reconocida.

\section{Runtime Application Self-Protection (RASP)}
\label{sec:rasp}
RASP designa una tecnología de seguridad que se integra directamente dentro de la aplicación o su
entorno de ejecución, en lugar de operar como un control perimetral externo, habilitando la
detección y prevención de ataques en tiempo real desde el interior del propio proceso \cite{contrastsecurity,bermejohiguera2022}. Esta distinción arquitectónica es la que define la categoría tecnológica en la
que se posiciona \textit{logSguarDian}: a diferencia de un WAF o un IDS de red, que inspeccionan
el tráfico desde fuera del proceso y pierden visibilidad sobre el contenido cifrado o el contexto de
ejecución de la aplicación, un sistema RASP intercepta la petición dentro del propio proceso
Node.js, con acceso directo al objeto \texttt{req} antes de que cualquier capa de cifrado externa
sea relevante. La Figura \ref{fig:rasp-vs-perimetral} ilustra esta diferencia.

\begin{figure}[H]
\centering
\includegraphics[width=0.95\textwidth]{figura3.png}
\caption[Comparación arquitectónica: control perimetral vs. RASP]{Comparación arquitectónica entre un control perimetral (WAF/IDS, izquierda) y un sistema RASP operando dentro del proceso (\textit{logSguarDian}, derecha)}
\label{fig:rasp-vs-perimetral}
\end{figure}
\vspace{-0.3cm}
{\footnotesize Nota. Elaboración propia, con base en la definición de RASP de Contrast Security \cite{contrastsecurity} y el análisis
comparativo de herramientas de protección en tiempo de ejecución de Sureda Riera et al.
\cite{bermejohiguera2022}.}

La arquitectura RASP no está exenta de limitaciones que deben reconocerse en el diseño de
\textit{logSguarDian}. La primera es el \textit{overhead} de rendimiento introducido por la
instrumentación: al interceptar cada petición para analizarla antes de que llegue al route handler,
el middleware añade latencia medible. Este riesgo se mitiga mediante el procesamiento asíncrono
no bloqueante descrito en la sección \ref{sec:onnx-inferencia}, y su magnitud es precisamente uno
de los criterios de aceptación cuantitativos del proyecto. La segunda es el riesgo de falsos
positivos: un modelo que clasifica tráfico legítimo como malicioso degrada el servicio desde la
perspectiva del negocio. En \textit{logSguarDian} este riesgo depende del umbral de bloqueo del
Random Forest (\texttt{RF\_THRESHOLD}), único componente con autoridad de bloqueo; el Isolation
Forest solo marca anomalías en el registro y no bloquea peticiones
(Sección~\ref{subsec:met-f2}). La tercera es la
dependencia de cobertura: un sistema RASP únicamente detecta ataques que atraviesan el
middleware instrumentado, por lo que ataques dirigidos a la infraestructura subyacente (la
instancia de base de datos, el proveedor de nube, otras dependencias npm) quedan fuera de su
visibilidad, lo que posiciona a \textit{logSguarDian} como una capa complementaria dentro de un
modelo de defensa en profundidad (Figura~\ref{fig:defensa-profundidad}) y no como sustituto de
los controles perimetrales descritos en la Sección~\ref{sec:soluciones-existentes}. La cuarta es la
complejidad de actualización del
modelo embebido: una vez instalada la librería en un proyecto Node.js, el modelo de ML que opera
dentro de ella corresponde a la versión con la que fue entrenado y publicado, de modo que, si
emergen nuevos vectores de ataque para los que el modelo no fue entrenado, el desarrollador debe
actualizar la librería para incorporar un modelo reentrenado. A diferencia de los WAF administrados
en la nube, que actualizan sus reglas de forma centralizada y transparente, \textit{logSguarDian}
depende del ciclo estándar de actualización de paquetes npm, lo que introduce una ventana de
exposición ante amenazas emergentes entre versiones publicadas.

\section{Vectores de ataque cubiertos por logSguarDian}
El Cuadro \ref{cuadro:vectores} resume los cuatro vectores de ataque intra-request que
\textit{logSguarDian} está diseñado para detectar, su punto de materialización específico en el
ecosistema Node.js/Express y el tipo de detección que aplica el sistema.

\begin{table}[H]
\centering
\small
\begin{tabular}{|L{2.1cm}|L{3.4cm}|L{3.8cm}|L{2.1cm}|}
\hline
\textbf{Vector} & \textbf{Punto de materialización en Node.js} & \textbf{Indicador principal en logs} & \textbf{Tipo de detección} \\ \hline
SQL Injection & Parámetros de URL, cuerpo POST, query strings concatenados en consultas & Tokens SQL en parámetros; respuestas HTTP 500; longitud anómala de campo & Supervisado (clasificación) \\ \hline
Cross-Site Scripting & Campos de texto, cabeceras y parámetros renderizados sin sanitizar & Etiquetas HTML/\texttt{<script>}; codificaciones URL/Base64 en campos de texto & Supervisado (clasificación) \\ \hline
Path Traversal / LFI & Parámetros de ruta o nombre de archivo usados en el módulo \texttt{fs} & Secuencias \texttt{../}; rutas absolutas a archivos sensibles del sistema & Supervisado (clasificación) \\ \hline
Command Injection & Módulo \texttt{child\_process} (\texttt{exec}, \texttt{execSync}) con entrada no sanitizada & Metacaracteres de shell (\texttt{;}, \texttt{\textbar}, \texttt{\&\&}); nombres de comandos del SO & Supervisado (clasificación) \\ \hline
\end{tabular}
\caption[Vectores de ataque cubiertos por logSguarDian]{Vectores de ataque cubiertos por \textit{logSguarDian}: punto de materialización, indicadores en logs y tipo de detección}
\label{cuadro:vectores}
\end{table}
{\footnotesize Nota. Elaboración propia con base en OWASP \cite{owaspa01_2021,owaspa03_2021}. El Random Forest es el único componente que bloquea; el Isolation Forest marca anomalías en el registro sobre cualquier petición, sin bloquearla.}

Los indicadores del Cuadro~\ref{cuadro:vectores} que dependen de la respuesta del servidor, como
los códigos HTTP 500 o 200 o la concentración de errores en un \textit{endpoint}, son útiles para
el análisis posterior de los registros, pero no forman parte del vector de características.
\textit{logSguarDian} decide antes de que exista una respuesta, por lo que el código de estado se
excluye de la entrada de ambos modelos (Sección~\ref{subsec:met-f2}); la detección se basa
únicamente en el contenido y la estructura de la petición.

\subsection{SQL Injection}
La inyección SQL ocurre cuando datos de entrada no sanitizados se concatenan directamente en
una consulta a base de datos, permitiendo a un atacante alterar la lógica de la consulta original.
OWASP clasifica este vector (\textit{Common Weakness Enumeration}, CWE-89) dentro de la categoría A03:2021 -- Injection, y documenta
que una consulta alterada puede devolver la totalidad de los registros de una tabla y, en
escenarios más graves, modificar o eliminar datos e incluso invocar procedimientos almacenados
\cite{owaspa03_2021}. En el contexto de Node.js, el vector se materializa en cualquier endpoint
que construya consultas mediante concatenación de cadenas en lugar de consultas parametrizadas;
un atacante que envía \texttt{' OR 1=1 --} en un campo de autenticación puede lograr que la
consulta resultante devuelva la totalidad de una tabla, evadiendo el control de acceso.

Los intentos de inyección SQL dejan una huella característica en los logs de Node.js que permite
su detección mediante análisis de comportamiento: la presencia de tokens SQL en los parámetros
(comillas simples, doble guion, punto y coma, y palabras reservadas como UNION, SELECT, DROP
u OR); un patrón de respuestas HTTP 500 concentradas en un mismo endpoint en un período corto,
generado cuando el payload malformado llega a la capa de consulta sin ser filtrado; y una longitud
de campo anómalamente mayor a la esperada, por ejemplo un campo de usuario que normalmente
recibe entre 5 y 30 caracteres y recibe un payload de 200 o más durante un ataque. La combinación
de estos tres indicadores es la firma comportamental que \textit{logSguarDian} extrae para
clasificar este vector mediante el modelo supervisado.

\subsection{Cross-Site Scripting (XSS)}
El XSS consiste en la inserción de código JavaScript malicioso en contenido que posteriormente es
renderizado por otros usuarios, en sus variantes reflejada, almacenada y basada en DOM. OWASP
lo clasifica (CWE-79) dentro de A03:2021 -- Injection \cite{owaspa03_2021}. Dentro de la
taxonomía MITRE ATT\&CK, este vector se corresponde con la técnica T1189
(\textit{Drive-by Compromise}), cuya documentación oficial identifica explícitamente la inserción
de scripts maliciosos mediante interfaces de aplicaciones web que permiten contenido controlado
por el usuario como uno de los mecanismos de entrega de esta técnica \cite{mitret1189}. Esta
técnica se centra en comprometer el equipo cliente del usuario que visita el sitio, no el servidor,
lo que es coherente con el hecho de que el XSS afecta primariamente a los usuarios de la
aplicación y no a la infraestructura que la aloja. En una API REST, el riesgo se materializa
cuando la API sirve contenido dinámico renderizado directamente en el navegador del cliente; un atacante que inyecta un script
almacenado puede robar tokens de sesión o \textit{JSON Web Tokens} (JWT) y ejecutar acciones autenticadas en nombre del
usuario comprometido.

En los logs de Node.js, los intentos de XSS presentan tres indicadores observables. El primero es
la presencia de etiquetas o atributos HTML en campos que por diseño de la API solo deberían
recibir texto plano, como \texttt{<script>}, \texttt{onerror=}, \texttt{onload=} o \texttt{javascript:}.
El segundo es el uso de codificaciones para evadir filtros superficiales, como codificación URL
(\texttt{\%3Cscript\%3E}) o de entidades HTML (\texttt{\&lt;script\&gt;}). El tercero es un patrón
temporal doble característico del XSS almacenado: una petición POST o PUT con el payload hacia
un endpoint de escritura, seguida de peticiones GET al recurso donde ese contenido queda
almacenado.

\subsection{Path Traversal / Local File Inclusion (LFI)}
El Path Traversal consiste en la manipulación de rutas de archivo mediante secuencias como
\texttt{../} para acceder a archivos fuera del directorio raíz autorizado por la aplicación \cite{owasptraversal}. LFI es un vector relacionado que ocurre cuando la aplicación
no solo lee, sino que incluye y evalúa el archivo referenciado en su propio contexto de
ejecución, lo que puede escalar a ejecución remota de código. Ambos vectores comparten la misma
causa raíz: la ausencia de validación y canonicalización de la entrada de usuario antes de
construir rutas del sistema de archivos.

A diferencia de los otros tres vectores, el OWASP Top 10:2021 no clasifica el Path Traversal
como una categoría de inyección: los identificadores CWE-22 (\textit{Path Traversal}) y CWE-23
(\textit{Relative Path Traversal}) están mapeados en A01:2021 -- Broken Access Control
\cite{owaspa01_2021}, mientras que SQL Injection (CWE-89), XSS (CWE-79) y la inyección de comandos
del sistema operativo (CWE-78) están mapeados en A03:2021 -- Injection \cite{owaspa03_2021}. La
inclusión del Path Traversal/LFI entre los vectores cubiertos por \textit{logSguarDian} no se
justifica, por lo tanto, por pertenecer a la misma categoría OWASP que los demás, sino por
compartir con ellos una causa raíz operativa (entrada controlada por el cliente que llega sin
validación a un recurso sensible del servidor) y una huella observable en el objeto \texttt{req}
que un mismo extractor de características puede capturar. En Node.js, el módulo \texttt{fs} y los
servidores de archivos estáticos son los puntos de integración más comunes donde este vector se
materializa cuando la entrada no es sanitizada.

Las técnicas de evasión documentadas añaden complejidad a la detección mediante reglas
estáticas: codificación URL simple (\texttt{\%2e\%2e\%2f} como equivalente de \texttt{../}),
codificación URL doble (\texttt{\%252e\%252e\%252f}), codificación Unicode (\texttt{..\%c0\%af}) e
inserción de bytes nulos para truncar extensiones verificadas por el servidor \cite{owasptraversal}. Esta variabilidad sintáctica es la que hace insuficiente una lista de denegación
basada en firmas fijas, y la que justifica un clasificador supervisado entrenado sobre la
distribución estadística de estos patrones. Los indicadores observables en los logs de Node.js son
tres: la secuencia \texttt{../} y sus variantes codificadas en parámetros que por diseño solo deberían
recibir nombres de archivo; la aparición de rutas absolutas o nombres de archivos sensibles del
sistema operativo (\texttt{/etc/passwd}, \texttt{.env}, \texttt{config.json}); y una respuesta HTTP 200
con contenido no esperado, o 500 cuando la ruta construida no existe.

\subsection{Command Injection}
La inyección de comandos del sistema operativo ocurre cuando una aplicación transmite entrada de
usuario no sanitizada a un intérprete de comandos del sistema; OWASP la clasifica (CWE-78) dentro
de A03:2021 -- Injection \cite{owaspa03_2021}. En Node.js, la
superficie de ataque se concentra en el módulo \texttt{child\_process}, cuyas funciones
\texttt{exec} y \texttt{execSync} invocan un shell del sistema operativo que interpreta
metacaracteres como \texttt{;}, \texttt{\textbar}, \texttt{\&\&} y \texttt{\$()}. Esta clase de
vulnerabilidad está documentada en incidentes reales del ecosistema npm: CVE-2021-27185 afectó
al paquete \texttt{samba-client}, donde el uso de \texttt{process.exec} sin sanitización de entrada
permitía inyección de comandos con un puntaje CVSS crítico de 9.8 \cite{gitlabadvisory2021}.

La variabilidad sintáctica de los payloads entre plataformas aumenta la complejidad de su
detección mediante reglas estáticas: en Linux, los metacaracteres más frecuentes son \texttt{;},
\texttt{\textbar}, \texttt{\&\&}, \texttt{\textbar\textbar} y la sustitución de comandos mediante
\texttt{`cmd`} o \texttt{\$(cmd)}; en Windows, los equivalentes son \texttt{\&}, \texttt{\textbar},
\texttt{\&\&} y \texttt{\textbar\textbar} aplicados sobre \texttt{cmd.exe}. Los indicadores observables
en los logs de Node.js son la presencia de estos metacaracteres en campos que por diseño deberían
recibir valores alfanuméricos, y la aparición de nombres de comandos del sistema operativo
(\texttt{cat}, \texttt{whoami}, \texttt{wget}, \texttt{curl}) dentro de parámetros de la API.

\section{Marcos de referencia para la selección de vectores}
La selección de los cuatro vectores cubiertos por \textit{logSguarDian} no responde a un criterio
arbitrario, sino a su correspondencia con marcos de referencia reconocidos internacionalmente y
basados en evidencia empírica de incidentes reales. Esa correspondencia no es uniforme: los
cuatro vectores no pertenecen a una única categoría de OWASP, y el Cuadro~\ref{cuadro:owasp-mitre-cwe}
documenta con precisión dónde se ubica cada uno.

El OWASP Top 10 es una lista de las diez categorías de riesgo más críticas en aplicaciones web,
construida a partir del análisis de vulnerabilidades reportadas en la industria \cite{owasp2021}.
Este trabajo emplea la edición 2021, sobre la que se definieron el protocolo y la taxonomía del
dataset; la edición 2025 ya está disponible \cite{owasp2025top10}. En ella Injection pasa de A03 a A05 y
conserva los CWE de inyección de comandos, XSS y SQL (CWE-78, CWE-79 y CWE-89), y Path Traversal
(CWE-22) sigue dentro de A01, por lo que la asignación de los cuatro vectores no cambia de fondo; la
edición 2025 no se adoptó para mantener la trazabilidad con las decisiones ya tomadas. En la edición 2021, tres de los cuatro vectores se agrupan bajo A03:2021 --
Injection: la inyección SQL (CWE-89), el XSS (CWE-79) y la inyección de comandos del sistema
operativo (CWE-78) \cite{owaspa03_2021}. El Path Traversal (CWE-22) no pertenece a esa categoría:
está mapeado en A01:2021 -- Broken Access Control \cite{owaspa01_2021}. La primera versión de
\textit{logSguarDian} cubre, por lo tanto, vectores de dos categorías OWASP distintas, y la
justificación de agruparlos en una misma herramienta es operativa (todos materializan su payload
en el objeto \texttt{req} y son detectables a partir de las mismas características extraídas), no
taxonómica.

\begin{table}[H]
\centering
\small
\begin{tabular}{|L{3.4cm}|L{2.2cm}|L{4.4cm}|L{2.4cm}|}
\hline
\textbf{Vector} & \textbf{CWE principal} & \textbf{OWASP Top 10:2021} & \textbf{MITRE ATT\&CK} \\ \hline
SQL Injection & CWE-89 & A03 -- Injection & T1190 \\ \hline
Cross-Site Scripting & CWE-79 & A03 -- Injection & T1189 \\ \hline
Path Traversal / LFI & CWE-22 & A01 -- Broken Access Control & T1190 \\ \hline
Command Injection & CWE-78 & A03 -- Injection & T1190 \\ \hline
\end{tabular}
\caption[Mapeo de vectores a CWE, OWASP Top 10 y MITRE ATT\&CK]{Mapeo de los cuatro vectores cubiertos por \textit{logSguarDian} a CWE, OWASP Top 10:2021 y MITRE ATT\&CK}
\label{cuadro:owasp-mitre-cwe}
\end{table}
\vspace{-0.3cm}
{\footnotesize Nota. Elaboración propia, con base en OWASP \cite{owaspa01_2021, owaspa03_2021} y MITRE \cite{mitret1189, mitret1190}. Dentro de las categorías del Top 10:2021 usadas aquí no existe un CWE específico de inclusión local
de archivos. A01 mapea CWE-22 \cite{owaspa01_2021} y A03 mapea CWE-98 \cite{owaspa03_2021}, que
MITRE define para la inclusión de archivos en PHP y que documenta ``local file inclusion'' solo como
término alterno \cite{mitre2026cwe98}. LFI se asocia aquí a CWE-22 \cite{mitre2026cwe22} como
variante de la manipulación de rutas.}

El OWASP API Security Top 10, en su edición 2023, reconoce que la superficie de ataque de las
APIs modernas difiere estructuralmente de las aplicaciones web tradicionales, y posiciona como
riesgo principal la Autorización Rota a Nivel de Objeto (BOLA, API1:2023) \cite{owaspapi2023};
esta categoría no está cubierta por \textit{logSguarDian} en su primera versión, consistente con
la delimitación de alcance ya señalada para el modelo STRIDE. La categoría API3:2023 (Broken
Object Property Level Authorization), por su parte, abarca la exposición de propiedades sensibles
de un objeto y la asignación masiva de propiedades que el cliente no debería poder modificar
\cite{owaspapi3_2023}; es un riesgo de autorización a nivel de propiedad y no guarda relación con
la manipulación de rutas de archivo, por lo que ninguno de los cuatro vectores cubiertos tiene un
equivalente directo en el API Security Top 10 2023.

El framework MITRE ATT\&CK es una base de conocimiento pública que documenta tácticas,
técnicas y procedimientos utilizados por adversarios reales en incidentes observados
\cite{mitreattack}. Para el dominio Enterprise, relevante para este proyecto, la táctica de Acceso
Inicial es la que agrupa los vectores cubiertos. La técnica T1190 (\textit{Exploit Public-Facing
Application}) describe la explotación de debilidades en sistemas expuestos a Internet, incluyendo
sitios web y bases de datos, y remite al OWASP Top 10 como referencia de las vulnerabilidades web
más comunes \cite{mitret1190}; es la técnica que corresponde a SQL Injection, Path Traversal y
Command Injection. La técnica T1189 (\textit{Drive-by Compromise}) se reserva para el XSS, dado
que su documentación oficial lo menciona como mecanismo de entrega y se centra en el compromiso
del cliente y no del servidor \cite{mitret1189}. La alineación de los vectores seleccionados con
ambos marcos garantiza que la cobertura de amenazas de \textit{logSguarDian} responde a evidencia
empírica de incidentes reales, y no a un criterio de selección arbitrario.

\section{Fundamentos de detección basada en aprendizaje automático}
\label{sec:ml-conceptos}
La detección de los cuatro vectores cubiertos por \textit{logSguarDian} combina dos paradigmas de
aprendizaje automático que responden a dos naturalezas distintas de amenaza. Esta sección
introduce ambos paradigmas al nivel conceptual necesario para comprender las decisiones de
arquitectura del sistema; el detalle del entrenamiento, la selección de hiperparámetros y la
evaluación comparativa de modelos candidatos se documenta en el trabajo de graduación de Diego
Valenzuela, correspondiente al componente de Inteligencia Artificial del proyecto.

\subsection{Clasificación supervisada}
El aprendizaje supervisado construye un modelo a partir de un conjunto de datos etiquetado, en el
que cada muestra de tráfico HTTP está marcada como legítima o como perteneciente a una de las
cuatro categorías de ataque cubiertas. Random Forest (RF), el algoritmo de clasificación supervisada considerado para \textit{logSguarDian}, es un método de ensamble que combina las predicciones de
múltiples árboles de decisión entrenados sobre subconjuntos aleatorios de características y
muestras, reduciendo el sobreajuste que afecta a un árbol de decisión individual
\cite{breiman2001random}. Este paradigma es apropiado para los cuatro vectores cubiertos porque cada uno
presenta patrones sintácticos y estructurales relativamente consistentes, discernibles como
características extraíbles del objeto \texttt{req} (presencia de tokens SQL, metacaracteres de
shell, secuencias de \textit{traversal}), sobre los cuales un clasificador supervisado puede
aprender fronteras de decisión efectivas.

\subsection{Detección de anomalías no supervisada}
El aprendizaje no supervisado, a diferencia del anterior, no requiere que las muestras de
entrenamiento estén etiquetadas: aprende la distribución estadística del comportamiento
considerado normal y señala como anómala cualquier observación que se desvíe significativamente de
ella. Isolation Forest (IF), el algoritmo de detección de anomalías empleado en \textit{logSguarDian},
aísla observaciones mediante particiones aleatorias recursivas del espacio de características: las
anomalías, al ser estadísticamente distintas de la mayoría de los datos, requieren en promedio
menos particiones para quedar aisladas, lo que permite asignar una puntuación de anomalía sin
etiquetas \cite{liu2008isolation}. Este paradigma complementa al clasificador supervisado, que
asigna cualquier entrada a la clase conocida más cercana aunque no se parezca a ninguna de las
vistas durante el entrenamiento. Su papel en el sistema es, sin embargo, acotado: el Isolation
Forest no bloquea peticiones, solo las marca como anómalas en el registro (\texttt{pass\_anomaly}),
y su contribución a la detección de variantes no vistas solo se midió, de forma indirecta, sobre el
conjunto de validación del componente de Inteligencia Artificial \cite{valenzuela2026ia}
(Sección~\ref{sec:met-validez}). La correspondencia entre cada paradigma y el tipo de indicador
de seguridad que produce se documenta en la Sección~\ref{sec:ioc-ioa}.

\subsection{Inferencia en tiempo de ejecución mediante ONNX}
\label{sec:onnx-inferencia}
Independientemente del algoritmo de origen, el componente de inferencia de \textit{logSguarDian}
consume un modelo serializado en formato ONNX (\textit{Open Neural Network Exchange}), un
formato estándar abierto para representar modelos de aprendizaje automático de forma
interoperable entre distintos frameworks de entrenamiento \cite{onnx2026}. La librería carga este
archivo mediante \texttt{onnxruntime-node} \cite{onnxruntimedocs} y ejecuta la inferencia dentro de
un \texttt{worker\_thread} de Node.js independiente del hilo principal; los worker threads están
pensados para operaciones de JavaScript intensivas en CPU \cite{nodejs2026worker}, y con esta
decisión de diseño se busca que la carga computacional del modelo no ocupe el Event Loop durante el
ciclo de vida de la petición. Esta decisión de arquitectura es la aplicación directa, a nivel de
implementación, de la restricción de disponibilidad descrita en la Sección~\ref{subsec:alcance-tecnico}:
sin este aislamiento, la propia herramienta de detección se convertiría en la causa de la
degradación de servicio que está diseñada para prevenir.

\section{Logs de servidor y análisis de comportamiento}
Los logs de servidor son la fuente primaria de información sobre el comportamiento de una
aplicación en producción. En Node.js y Express, los logs de acceso registran método, URL, código
de estado, tamaño de la solicitud y respuesta, dirección IP y tiempo de respuesta; los logs de error
capturan excepciones y fallos de validación que pueden revelar intentos de explotación, un volumen
inusualmente alto de respuestas 400 o 500 en un endpoint específico es indicador temprano de
sondeo con payloads malformados.

El análisis de comportamiento parte de establecer una línea base de tráfico normal, contra la cual
se evalúan desviaciones. Estos patrones pueden ser volumétricos (alta frecuencia de solicitudes
desde una misma IP, picos de errores 401/403) o semánticos (caracteres asociados a inyección en
parámetros, longitud anómala de campos, codificaciones inusuales). \textit{logSguarDian} trabaja
solo con los patrones semánticos de cada petición individual: el clasificador supervisado y el
detector de anomalías reciben ambos el vector de características de una sola petición, y las
características temporales del extractor (ventanas de 1, 5 y 60~s) no forman parte de su entrada.
Los patrones volumétricos y de secuencia, que requieren correlacionar varias peticiones, quedan
fuera del alcance de esta versión, junto con la denegación de servicio
(Sección~\ref{subsec:fuera-alcance}).

Morgan es el middleware de registro de peticiones HTTP de referencia en el ecosistema Express,
distribuido bajo la organización oficial \texttt{expressjs} en npm \cite{expressjsmorgan}. Se apoya
en el sistema de \textit{streams} de Node.js para escribir el registro sin bloquear la respuesta, y
permite definir tokens personalizados mediante \texttt{morgan.token()} para capturar campos que sus
formatos predefinidos no incluyen \cite{expressjsmorgan}. Es, por lo tanto, el punto de referencia
de cómo se registra el tráfico en una aplicación Express típica. \textit{logSguarDian} no depende
de Morgan ni extiende su registro: intercepta el objeto \texttt{req} dentro de su propia función
middleware, extrae el vector de características antes de que el manejador de ruta procese la
petición y almacena cada detección en su propio registro de eventos en SQLite. La diferencia es de
propósito: Morgan registra el tráfico para su consulta posterior, mientras que
\textit{logSguarDian} analiza cada petición para decidir, en el momento, si la bloquea.

\section{Indicadores de Compromiso (IoC) e Indicadores de Ataque (IoA)}
\label{sec:ioc-ioa}
Un Indicador de Compromiso (IoC) es un artefacto técnico u observable que sugiere que un ataque es
inminente o está en curso, o que ya pudo haberse producido un compromiso
\cite{nistioc2025,crowdstrikeioc}. Un Indicador de Ataque (IoA), en cambio, se centra en identificar
la intención y el comportamiento de un adversario mientras el ataque está en curso, enfoque que CrowdStrike presentó en 2016 como alternativa a la detección basada en IoC, a los que
critica por equivaler a firmas \cite{crowdstrike2016ioa,crowdstrikeioa}. Esta distinción es operativamente central para
\textit{logSguarDian}, que opera dentro del proceso Node.js e intercepta cada petición antes de que
llegue al manejador de ruta: las alertas que produce son, por construcción, indicadores de ataque,
porque se generan mientras el ataque está en curso y no a partir de evidencia forense posterior.
Ambos componentes producen este tipo de indicador: el Random Forest al clasificar y bloquear la
petición, y el Isolation Forest al marcarla como anómala en el registro. El
Cuadro~\ref{cuadro:ioc-ioa} sintetiza la diferencia entre ambos tipos de indicador.

\begin{table}[H]
\centering
\small
\begin{tabular}{|L{3.8cm}|L{4.4cm}|L{4.6cm}|}
\hline
\textbf{Dimensión} & \textbf{Indicador de Compromiso (IoC)} & \textbf{Indicador de Ataque (IoA)} \\ \hline
Naturaleza & Reactivo & Proactivo \\ \hline
Momento de detección & Típicamente después del incidente & Durante el ataque \\ \hline
Base de conocimiento & Firmas y artefactos conocidos & Comportamiento y patrones de la petición \\ \hline
Efectividad ante variantes desconocidas & Baja & Alta \\ \hline
\end{tabular}
\caption[Comparativa entre IoC e IoA en logSguarDian]{Comparativa entre Indicadores de Compromiso (IoC) e Indicadores de Ataque (IoA)}
\label{cuadro:ioc-ioa}
\end{table}
\vspace{-0.3cm}
{\footnotesize Nota. Elaboración propia, con base en CrowdStrike \cite{crowdstrikeioc, crowdstrikeioa}.}

\section{Soluciones de seguridad existentes}
\label{sec:soluciones-existentes}
El Cuadro \ref{cuadro:soluciones} compara las categorías de solución de seguridad disponibles
frente a \textit{logSguarDian} en las dimensiones más relevantes para el contexto de una startup
guatemalteca sin equipo de seguridad dedicado.

\begin{table}[H]
\centering
\small
\begin{tabular}{|L{2.6cm}|L{2.6cm}|L{1.6cm}|L{1.4cm}|L{1.9cm}|L{1.8cm}|}
\hline
\textbf{Solución} & \textbf{Tipo de análisis} & \textbf{Visibilidad capa 7} & \textbf{Equipo dedicado} & \textbf{Costo} & \textbf{Integración nativa en Node.js} \\ \hline
WAF (Cloudflare, AWS WAF) & Firmas estáticas & Parcial & Sí & Alto & No (perimetral) \\ \hline
IDS (Suricata, Snort) & Tráfico de red & No (ciego a HTTPS) & Sí & Medio & No (nivel de red) \\ \hline
SIEM (Splunk, Wazuh) & Correlación de eventos & Parcial & Sí (centro de operaciones de seguridad, SOC) & Muy alto & No (infraestructura externa) \\ \hline
\textit{logSguarDian} & Clasificación ML por petición & Sí (dentro del proceso) & No & Gratuito & Sí (middleware nativo) \\ \hline
\end{tabular}
\caption[Comparativa de soluciones de seguridad para APIs REST en Node.js]{Comparativa de soluciones de seguridad existentes frente a \textit{logSguarDian}}
\label{cuadro:soluciones}
\end{table}
\vspace{-0.3cm}
{\footnotesize Nota. Elaboración propia. Las calificaciones (Alto, Medio, Muy alto, equipo dedicado, visibilidad de
capa 7) son criterio del autor; el costo relativo del SIEM se apoya en ENISA \cite{enisa2021smes}.}

\subsection{Web Application Firewall (WAF)}
Un WAF opera como intermediario entre clientes y servidor, inspeccionando el tráfico HTTP/HTTPS
contra un conjunto de reglas y firmas predefinidas, principalmente en la capa 7 del modelo OSI.
Los WAF basados en firmas y en el OWASP Core Rule Set (CRS) requieren ajuste continuo y mantenimiento especializado para ser efectivos sin generar un volumen
excesivo de falsos positivos: el propio proyecto CRS advierte que los niveles de paranoia más altos
aumentan la posibilidad de bloquear tráfico legítimo por falsas alarmas y que puede ser necesario
agregar reglas de exclusión para aplicaciones con entradas complejas \cite{crs2026paranoia}. Por
construcción, una regla solo detecta lo que su patrón describe, por lo que estos WAF son poco
eficaces contra vulnerabilidades de lógica de negocio o variantes de ataque que no coinciden con
ninguna regla existente. Estas
dos limitaciones, mantenimiento especializado y ceguera ante variantes desconocidas, son
precisamente las que un modelo de aprendizaje automático entrenado sobre la distribución
estadística de los patrones de ataque, en lugar de un conjunto fijo de reglas, está diseñado para
mitigar.

\subsection{Sistemas de Detección de Intrusiones (IDS)}
Un IDS monitorea tráfico de red o actividad de sistema en busca de comportamientos sospechosos
\cite{nist80094}. Los IDS de red pierden visibilidad completa frente al tráfico cifrado mediante
HTTPS, dado que sin descifrar la conexión TLS (\textit{Transport Layer Security}) solo observan metadatos y patrones del flujo
\cite{papadogiannaki2021survey}, una ceguera crítica para detectar ataques de capa de aplicación como SQLi o XSS, donde el vector
malicioso está en el contenido del cuerpo de la solicitud. Adicionalmente, en la revisión realizada para este trabajo no se identificó ningún IDS convencional,
de red o de host, diseñado para integrarse en el modelo de concurrencia orientado a eventos de
Node.js, lo que deja una brecha arquitectónica que es la que justifica el enfoque de
\textit{logSguarDian} como middleware nativo: en lugar de observar el tráfico desde fuera del
proceso, opera desde dentro, con acceso pleno al contenido descifrado y al contexto de ejecución.

\subsection{Sistemas de Gestión de Eventos e Información de Seguridad (SIEM)}
Un SIEM agrega, correlaciona y analiza eventos de seguridad provenientes de múltiples fuentes de
infraestructura. Su principal limitación para el contexto de startups guatemaltecas es el costo: para las
pequeñas y medianas empresas, los SIEM representan inversiones considerables y, como muchas
soluciones de ciberseguridad, requieren conocimiento especializado para implementarse y
administrarse \cite{enisa2021smes}. Las soluciones SIEM ofrecidas como servicio implican además el
envío de los registros a una infraestructura externa, lo que introduce costo operativo y
consideraciones de privacidad de datos.

El capítulo de Antecedentes documenta en detalle la comparación entre las soluciones RASP
comerciales existentes (Contrast Protect, Datadog, OpenRASP) y \textit{logSguarDian}
(Sección~\ref{sec:rasp-existentes} y su cuadro comparativo); esta sección no la repite y remite a
esa comparación para el detalle de soporte de Node.js, costo y dependencia de infraestructura
externa de cada alternativa.

\section{Herramientas de simulación, pruebas y serialización}
La construcción del entorno de validación de \textit{logSguarDian} se apoya en un conjunto de
herramientas de código abierto seleccionadas por su reproducibilidad y su adopción estándar en el
ecosistema Node.js. El entorno de simulación se implementa mediante Docker Compose, lo que
garantiza que el entorno de generación de tráfico y ataques sea reproducible en cualquier momento
del proyecto y por cualquier integrante del equipo \cite{dockerdocs}; esta práctica de
contenerización como mecanismo de reproducibilidad computacional está documentada en la
literatura de ciencia de datos como una condición necesaria para la validez de resultados
experimentales dependientes de un entorno específico \cite{moreau2023}. La generación de tráfico
legítimo se realiza mediante Artillery, una herramienta de pruebas de carga que define escenarios
de tráfico mediante configuración declarativa \cite{artillerydocs}. Para el tráfico malicioso, el
protocolo contemplaba sqlmap, OWASP ZAP y Commix \cite{sqlmap, owaspzap, stasinopoulos2019}. En la
ejecución, los payloads de SQL Injection del corpus de evaluación limpio se obtuvieron con sqlmap
y los de Path Traversal y Command Injection con generadores propios, deliberadamente no derivados
de listas públicas; OWASP ZAP se dirigió mediante su API REST (Sección~\ref{sec:met-instrumentos}). Commix
interviene de forma indirecta: la lista \texttt{command-injection-commix.txt} de SecLists, derivada
de esa herramienta, es una de las fuentes de la clase \texttt{cmdi} del corpus de entrenamiento y
del corpus de la Ronda~4.

La validación de la librería distribuida como paquete npm emplea Jest como framework de pruebas
unitarias e de integración, combinado con Supertest para ejecutar peticiones HTTP contra la
aplicación sin necesidad de levantar un servidor real \cite{jestdocs, supertestdocs}. Esta
combinación permite verificar de forma determinista que el middleware intercepta peticiones con
payloads maliciosos conocidos y responde con los códigos de estado correctos ante cada categoría de
ataque. La latencia no se verifica con estas pruebas, sino con el benchmark de carga descrito en la
Sección~\ref{subsec:met-f4}.

\ifdefined\modoensamble
\else
	\section*{Referencias}
	\printbibliography[heading=none]
	\end{document}
\fi

\fi

\ifdefined\CAPmetodologia
	\newpage
	% ==============================================================================
% METODOLOGÍA
% ==============================================================================
% Evolución del Marco Metodológico de protocolo_sebastian_v1 (3).pdf
% (sección 7), actualizada al estado real del proyecto al 2026-10-02.
%
% Fuentes de cada cifra y procedimiento (repositorio logSguarDian y
% logSguarDian-vulnerable-project):
%   - Dataset y partición: training/split.py, training/DEDUP_METHODOLOGY.md,
%     docs/dataset-audit.md, training/results/v11_test_results.json.
%   - Umbral RF: docs/decision-policy.md secciones 2.2.1 y 2.2.2.
%   - Corpus limpio (Ronda 5): attack-sim/build_corpus.py,
%     attack-sim/exclusion_filter.py, docs/config3b-results.md.
%   - Fuga: training/leakage_check.py, training/results/leakage_report.json.
%   - Latencia/memoria/CPU: .github/workflows/latency-benchmark.yml,
%     attack-sim/run_latency_suite.sh, artillery-*.yml, monitor-memory.sh.
%   - Apps de terceros: docs/vulnerable-projects-tests/,
%     docs/findings-evidence/juice-shop-ua-bias/.
%
% Decisiones del estudiante incorporadas (conversación 2026-10-02):
%   - Criterio de latencia absoluto de 15 ms (ya fijado en h-alcance.tex),
%     justificado porque una cifra en milisegundos es más representativa
%     que un porcentaje sobre una línea base de pocos milisegundos.
%   - Umbral RF calibrado con rf_v3.
%   - Tablero Kanban de cuatro columnas.
%   - Kanban en Notion + flujo main <- develop <- rama con aprobación
%     restringida a los dos integrantes + GitHub Actions.
%   - La validación con startups no se realizó; se sustituye por la
%     validación en aplicaciones de terceros.
%   - [Actualizado 2026-10-10: ZAP/XSS y Config 1/3a/3b ya se ejecutaron sobre el corpus limpio.]
%     Antes pendientes: (ZAP/XSS, Config 1/3a/3b con corpus
%     limpio, latencia con el modelo sincronizado, memoria en escenario
%     único) no se ejecutarán; se declaran como restricciones.
%   - El paquete mlops se incorpora por recomendación del asesor, del lado
%     del componente de IA; aquí solo se describe como interfaz.
%
% Fechas de la Tabla de fases: reconstruidas del historial de git de ambos
% repositorios (primer y último commit relacionado). Pendientes de
% confirmación por el estudiante.
\ifdefined\modoensamble
\else
	\input{0-preambulo_standalone}
\fi

\chapter{Metodología}
\label{cap:metodologia}

Este capítulo describe cómo se ejecutó el componente de ciberseguridad de \textit{logSguarDian} y
cómo se obtuvieron, procesaron y validaron los datos que sustentan el capítulo de Resultados.
Constituye una evolución del Marco Metodológico del Protocolo de Graduación: conserva el
enfoque, la metodología de gestión, las seis fases y las métricas allí definidas, e incorpora los
cambios que el avance real del proyecto hizo necesarios. Cada cambio se declara de forma
explícita en la Sección~\ref{sec:met-cambios}, junto con su causa; ninguna cifra o criterio del
protocolo se sustituye en silencio.

El proyecto integrador se compone de dos componentes desarrollados en paralelo. Este documento
cubre el componente de ciberseguridad; el diseño, entrenamiento y serialización de los modelos
de aprendizaje automático pertenecen al componente de Inteligencia Artificial y aquí se tratan
únicamente como entradas, con su contrato de interfaz (Sección~\ref{subsec:met-f2}).

% ------------------------------------------------------------------------------
\section{Enfoque y tipo de investigación}
\label{sec:met-enfoque}

El trabajo adopta un enfoque cuantitativo, de tipo aplicado y con diseño experimental. Es
cuantitativo porque todas las preguntas que plantean los objetivos se responden mediante
mediciones numéricas sobre conjuntos de datos y condiciones de carga definidos de antemano,
contrastadas con criterios de aceptación fijados antes de la medición
\cite{hernandez2014}. Es aplicado porque su producto, una librería npm publicada, responde a un
problema operativo concreto: la ausencia de protección a nivel de aplicación en servicios
Node.js de organizaciones sin equipo de seguridad dedicado.

El diseño es experimental en el sentido de que el investigador manipula de forma deliberada la
variable independiente y mide su efecto sobre variables dependientes en un entorno controlado
\cite{hernandez2014}. La variable independiente es la configuración de defensa del servidor
anfitrión, con cuatro niveles: sin protección (Configuración~1, grupo de control), con
\textit{logSguarDian} (Configuración~2), con un WAF ModSecurity con OWASP CRS (Configuración~3a)
y con ambas capas superpuestas (Configuración~3b). Las variables dependientes son la tasa de
detección por categoría de ataque, la tasa de falsos positivos sobre tráfico legítimo, la
latencia p95 de respuesta y el consumo de memoria y CPU. La línea base sin middleware y el
servidor anfitrión idéntico entre condiciones (misma imagen Docker, misma base de datos
sembrada) controlan las variables extrañas que podrían explicar las diferencias observadas.

La unidad de análisis es la petición HTTP individual dirigida a un endpoint de una aplicación
Express, caracterizada por su método, ruta, cadena de consulta, cuerpo y cabeceras. No participan
personas como sujetos de estudio; por esa razón no aplica un muestreo de población humana, y la
Sección~\ref{sec:met-unidad} define las poblaciones y muestras de peticiones utilizadas.

% ------------------------------------------------------------------------------
\section{Metodología de trabajo}
\label{sec:met-trabajo}

\subsection{Kanban como metodología de gestión}
La gestión del trabajo se realizó con Kanban \cite{anderson2010}, sobre un tablero en Notion
compartido por los dos integrantes del proyecto. La selección, justificada en el protocolo, se
mantuvo por tres razones que el desarrollo confirmó en la práctica:

\begin{itemize}
    \item \textbf{Ausencia de un receptor externo de incrementos.} No existe un Product Owner
    ni un cliente que valide entregas periódicas; las ceremonias de Scrum habrían carecido de
    propósito metodológico real.
    \item \textbf{Dependencias técnicas secuenciales.} El corpus es prerrequisito del
    entrenamiento, y el modelo serializado es prerrequisito de la integración en el middleware.
    Kanban representa ese flujo con columnas de estado sin imponer iteraciones de duración fija
    que obliguen a avanzar con una dependencia sin resolver.
    \item \textbf{Respuesta continua a hallazgos experimentales.} Durante el proyecto aparecieron
    hallazgos que obligaron a reabrir trabajo ya cerrado: la fuga de datos entre el corpus de
    evaluación y el de entrenamiento, el sesgo del modelo hacia la cabecera \texttt{User-Agent} y
    el colapso del servidor bajo volumen de datos realista. Kanban permitió repriorizarlos de
    inmediato sin esperar al cierre de un sprint.
\end{itemize}

El tablero se organizó en cuatro columnas que representan los estados posibles de una tarea:
No empezado, En proceso, En revisión y Completado. Respecto al protocolo, el tablero no incluyó
una columna Bloqueado. El criterio de salida de la columna En revisión es la verificación del indicador
cuantitativo asociado a la tarea (Sección~\ref{sec:met-metricas}), no una apreciación subjetiva
de completitud.

\subsection{Flujo de control de versiones e integración continua}
\label{subsec:met-flujo}
El tablero se complementó con un flujo de ramas en Git que convierte la revisión en un control
obligatorio y no opcional:

\begin{itemize}
    \item \textbf{Ramas.} Todo cambio se desarrolla en una rama propia con prefijo por tipo
    (\texttt{feat/}, \texttt{fix/}, \texttt{test/}, \texttt{docs/}, \texttt{perf/},
    \texttt{chore/}) creada desde \texttt{develop}. La integración sigue la secuencia
    \texttt{rama} $\rightarrow$ \texttt{develop} $\rightarrow$ \texttt{main}; un flujo de
    GitHub Actions rechaza cualquier pull request hacia \texttt{main} cuyo origen no sea
    \texttt{develop}.
    \item \textbf{Aprobación restringida.} Las reglas de protección de rama del repositorio solo
    permiten que los dos integrantes del proyecto aprueben y fusionen pull requests, de modo que
    cada cambio de un integrante es revisado por el otro.
    \item \textbf{Integración continua (CI).} Cada pull request ejecuta en GitHub Actions
    (\texttt{ubuntu-latest}) los trabajos de análisis estático (ESLint), compilación
    (\texttt{tsc}), pruebas del extractor de características (incluida la prueba de paridad
    numérica contra la implementación de referencia en Python), pruebas del paquete
    \texttt{core} (middleware, almacenamiento, webhooks y la interfaz de línea de comandos, CLI) y pruebas del paquete
    \texttt{mlops}. La fusión de \texttt{develop} a \texttt{main} publica automáticamente el
    paquete en el registro npm si la versión no ha sido publicada antes.
    \item \textbf{Trazabilidad de decisiones.} Las decisiones técnicas y su justificación se
    registran en una bitácora versionada, con una entrada por tarea cerrada, y cada
    reentrenamiento del modelo va precedido de un commit que fija el hash del conjunto de prueba
    (Sección~\ref{subsec:met-f1}), con una excepción: el rebloqueo del 23 de septiembre se hizo
    después de restaurar los datos de la generación vigente.
\end{itemize}

Este flujo tiene una finalidad metodológica, no solo de organización: cada cifra reportada en
Resultados queda asociada a un commit, a un script versionado y, cuando aplica, a una ejecución
de GitHub Actions cuyo registro puede consultarse posteriormente.

\subsection{Distribución de roles}
El Cuadro~\ref{tab:met-roles} distribuye las responsabilidades. Las responsabilidades del
componente de IA se listan solo para delimitar la frontera entre componentes.

\begin{table}[H]
\centering
\small
\begin{tabular}{|L{2.6cm}|L{4.9cm}|L{5.3cm}|}
\hline
\textbf{Integrante} & \textbf{Rol principal} & \textbf{Responsabilidades clave} \\ \hline
Sebastián Huertas & Componente de ciberseguridad & Extractor inicial de 72 características y el
paso a \texttt{worker\_threads}; selección y mapeo de vectores de ataque;
auditoría de fuentes y etiquetado; entorno Docker de evaluación; capa de captura y respuesta
(HTTP 403, registro, webhooks, CLI); calibración del umbral de bloqueo; corpus de evaluación
limpio; pruebas de penetración controlada; benchmarks de latencia, memoria y CPU; validación en
aplicaciones de terceros. \\ \hline
Diego Valenzuela & Componente de IA & Análisis y justificación de las características, definición de
las no disponibles en inferencia y del contrato de entrada; cuatro características adicionales y
correcciones del extractor; diseño de la arquitectura de inferencia en hilos (hilo dedicado de Random
Forest, grupo de Isolation Forest y lectura asíncrona de su resultado); entrenamiento, selección y
serialización ONNX de los modelos; pipeline de reentrenamiento offline (\texttt{mlops}). \\ \hline
Ambos & Integración y documentación & Arquitectura de la librería; pruebas de integración
extremo a extremo; revisión cruzada de pull requests; publicación en npm; documentación técnica
final. \\ \hline
\end{tabular}
\caption[Distribución de roles y responsabilidades]{Distribución de roles y responsabilidades del equipo}
\label{tab:met-roles}
\end{table}
{\footnotesize Nota. Elaboración propia.}

La frontera entre componentes en el pipeline de características es la siguiente. El extractor
inicial de 72 características y el paso a \texttt{worker\_threads} los
implementó Sebastián Huertas. Diego Valenzuela diseñó la arquitectura de inferencia en hilos (hilo
dedicado de Random Forest, grupo de Isolation Forest y lectura asíncrona de su resultado), analizó
y justificó cada característica, definió cuáles no existen en inferencia y fijó el contrato de
entrada, añadió cuatro características y las correcciones del extractor surgidas del análisis del
modelo, y desarrolló el entrenamiento, la serialización ONNX y el pipeline de reentrenamiento.

% ------------------------------------------------------------------------------
\section{Unidad de estudio, población y muestras}
\label{sec:met-unidad}

La población de interés es el conjunto de peticiones HTTP que puede recibir un endpoint de una
aplicación Node.js/Express, legítimas o pertenecientes a alguna de las cuatro categorías de
ataque del Objetivo Específico~1. Como esa población no es enumerable, el estudio emplea cinco
muestras, cada una construida para responder una pregunta distinta. El Cuadro~\ref{tab:met-muestras}
las resume.

\begin{table}[H]
\centering
\small
\begin{tabular}{|L{2.8cm}|L{3.6cm}|L{2.4cm}|L{3.9cm}|}
\hline
\textbf{Muestra} & \textbf{Composición} & \textbf{Tipo de muestreo} & \textbf{Pregunta que responde} \\ \hline
M1. Corpus etiquetado & 383{,}202 peticiones de siete fuentes públicas auditadas, dos listas públicas de payloads, conjuntos sintéticos y telemetría curada:
268{,}240 de entrenamiento, 57{,}481 de validación y 57{,}481 de prueba & Estratificado por
clase y fuente (70/15/15) & ¿El modelo clasifica correctamente las cuatro categorías? (OE2, OE3) \\ \hline
M2. Corpus limpio de evaluación (Ronda~5) & 200 SQLi (sqlmap), 200 Path Traversal y 200 Command
Injection (generadores propios); 1 payload de XSS confirmado con OWASP ZAP & Intencional, con filtro de exclusión contra el entrenamiento &
¿Qué proporción de payloads nunca vistos detecta el sistema desplegado? (OE3) \\ \hline
M3. Tráfico de carga & Navegación legítima y tráfico mixto generados con Artillery a
$\approx$100 peticiones/s & Generado, cinco repeticiones por condición & ¿Cuánto degrada el
middleware el rendimiento del servidor? (Objetivo general) \\ \hline
M4. Tráfico de navegador & Rutas sin parámetros (\texttt{/login}, \texttt{/posts},
\texttt{/profile}) con cadenas \texttt{User-Agent} de navegadores reales & Intencional & ¿El
modelo bloquea tráfico legítimo de navegadores reales? (OE3) \\ \hline
M5. Aplicaciones de terceros & DVNA \cite{appsecco2026dvna}, node-api-goat \cite{layro012026nodeapigoat}, dvws-node \cite{snoopysecurity2026dvwsnode} y OWASP Juice Shop \cite{kimminich2026juiceshop} & Intencional,
por diversidad de arquitectura & ¿La calibración generaliza fuera de la aplicación de
referencia? (Objetivo general) \\ \hline
\end{tabular}
\caption[Muestras de peticiones utilizadas en el estudio]{Muestras de peticiones HTTP utilizadas en el estudio}
\label{tab:met-muestras}
\end{table}
{\footnotesize Nota. Elaboración propia con datos de los repositorios del proyecto.}

La muestra M1 es un insumo compartido con el componente de IA; su tamaño satisface el mínimo de
100{,}000 muestras del OE2. La muestra M2 reemplaza al corpus de evaluación inicial (Ronda~4),
descartado por fuga de datos (Sección~\ref{subsec:met-fuga}). La categoría XSS de M2 se reduce a un único payload confirmado con OWASP ZAP (n=1); la restricción
correspondiente se declara en la Sección~\ref{sec:met-restricciones}.

% ------------------------------------------------------------------------------
\section{Fases, actividades y procedimientos}
\label{sec:met-fases}

El proyecto conserva las seis fases del protocolo. El Cuadro~\ref{tab:met-fases} compara el
período planificado con el período real, reconstruido a partir del historial de commits de los
repositorios del proyecto. Las fases se solaparon más de lo previsto: la detección de problemas
durante las pruebas obligó a reabrir las fases de construcción del dataset y de entrenamiento,
lo que corresponde al ciclo iterativo descrito en el protocolo.

\begin{table}[H]
\centering
\small
\begin{tabular}{|L{2.6cm}|L{1.9cm}|L{2.1cm}|L{1.2cm}|L{4.6cm}|}
\hline
\textbf{Fase} & \textbf{Protocolo} & \textbf{Real} & \textbf{OE} & \textbf{Entregable verificable} \\ \hline
F1. Construcción del dataset & 1--31 may & 11 may -- 20 jun & OE1, OE2 & Corpus $\geq$100{,}000
muestras etiquetadas; auditoría de fuentes y licencias; taxonomía de cinco clases \\ \hline
F2. Entrenamiento del modelo (IA) & 1--19 jun & 11 jun -- 26 sep & OE3 & Modelos ONNX con paridad
$<$0.1~\%; lectura del conjunto de prueba \\ \hline
F3. Desarrollo de la librería & 20 jun -- 1 ago & 30 jun -- 30 ago & OE3 & Paquete
\texttt{logsguardian@0.1.0} publicado en npm, con cuatro grupos de comandos CLI probados con
Jest \\ \hline
F4. Pruebas & 2 ago -- 30 sep & 23 jul -- 10 oct & OE3 & Resultados de las configuraciones 1, 2,
3a y 3b; corpus limpio; benchmark de rendimiento; validación en aplicaciones de terceros \\ \hline
F5. Aplicación de feedback & 15 ago -- 30 sep & 24 jul -- 11 oct & OE3 & Registro de iteraciones
con métricas antes y después de cada corrección; versión 0.1.1 publicada en npm \\ \hline
F6. Documentación & octubre & 16 ago -- oct & Todos & Informe final; README publicado en npm;
documentación técnica \\ \hline
\end{tabular}
\caption[Fases del proyecto: período planificado y real]{Fases del proyecto: período planificado en el protocolo, período real de ejecución y entregables verificables}
\label{tab:met-fases}
\end{table}
{\footnotesize Nota. Elaboración propia. El período real corresponde al primer y último commit
relacionado con cada fase en los repositorios \texttt{logSguarDian} y
\texttt{logSguarDian-vulnerable-project}; las fechas de publicación en npm son las del registro
(UTC). Todas las fechas son de 2026.}

Las subsecciones siguientes describen los procedimientos de cada fase. Para cada uno se indica
su finalidad y el objetivo al que contribuye.

\subsection{Fase 1: análisis de vectores y construcción del dataset}
\label{subsec:met-f1}

\textbf{Procedimiento 1.1. Mapeo de vectores a marcos de referencia (OE1).} Cada uno de los
cuatro vectores se asoció a su categoría del OWASP Top 10 2021, a su debilidad CWE y a su
técnica MITRE ATT\&CK, verificando cada correspondencia contra la página oficial del marco
\cite{owasp2021, owaspa01_2021, owaspa03_2021, mitret1190, mitret1189}. El mapeo distingue
SQL Injection, XSS y Command Injection (A03:2021, Injection) de Path Traversal/LFI (A01:2021,
Broken Access Control). Finalidad: fijar las cinco clases canónicas del dataset (\texttt{sqli},
\texttt{xss}, \texttt{path\_traversal}, \texttt{cmdi}, \texttt{benign}) y las categorías sobre
las que se calculan las métricas del OE3.

\textbf{Procedimiento 1.2. Auditoría de fuentes (OE2).} Se auditaron nueve fuentes públicas
candidatas, registrando para cada una su licencia (verificada contra la API o la página del
repositorio de origen), su DOI cuando existe, su formato, su número de registros y las
categorías que cubre. Se aprobaron ocho y se excluyó el dataset AIT-LDSv2 \cite{landauer2022ait,landauer2023maintainable} porque sus etiquetas
describen el rol del atacante y no el tipo de payload, lo que habría exigido un reetiquetado manual
con riesgo de error superior a su aporte de 3{,}435 filas. De las ocho aprobadas, siete aportan
filas al corpus; la lista de payloads de Path Traversal de omurugur \cite{ugur2019traversal} se aprobó, pero no tiene
analizador y no aporta filas. El corpus final incluye además dos listas públicas de payloads de Command Injection (SecLists y PayloadsAllTheThings
\cite{miessler2026seclists,swisskyrepo2026payloads}), incorporadas después de la auditoría y cuya
licencia (MIT, consultada en sus repositorios el 10 de octubre de 2026) no consta en ella, siete
conjuntos sintéticos propios y 119 filas de telemetría curada del pipeline de reentrenamiento
(Cuadro~\ref{tab:met-origenes}). El dataset de peticiones bloqueadas por OWASP ModSecurity
publicado en Zenodo \cite{lucz2025,lucz2025data}, que el protocolo designaba como base, aporta el 1.0~\% de las
filas; el volumen proviene principalmente de la colección CAPEC
\cite{sureda2022srbh,sureda2022newmultilabel} (55.7~\%) y de modsec-learn
\cite{pralab2024modseclearn,scano2024modseclearn} (26.0~\%). Las fuentes de licencia no especificada por su origen (XSS\_dataset.csv y command
injection.csv, dos de las siete auditadas presentes en el corpus, 3.4~\% de las filas) no se
redistribuyen: sus archivos de trabajo no se versionan en el repositorio y el proyecto entrega
únicamente los scripts de procesamiento y los vectores de características derivados.

\begin{table}[H]
\centering
\small
\begin{tabular}{|L{4.3cm}|L{1.8cm}|L{1.3cm}|L{5.3cm}|}
\hline
\textbf{Origen} & \textbf{Filas} & \textbf{\%} & \textbf{Licencia} \\ \hline
Colección CAPEC (SR-BH 2020) \cite{sureda2022srbh,sureda2022newmultilabel} & 213{,}514 & 55.72 & CC0-1.0, auditada \\ \hline
modsec-learn \cite{pralab2024modseclearn,scano2024modseclearn} & 99{,}645 & 26.00 & MIT, auditada \\ \hline
payload\_full.csv \cite{thakur2025payloads} & 29{,}719 & 7.76 & MIT, auditada \\ \hline
Payloads.csv \cite{shahane2021xss} & 18{,}411 & 4.80 & CC BY 4.0, auditada \\ \hline
XSS\_dataset.csv \cite{shah2020xss} & 10{,}854 & 2.83 & No especificada por el origen, auditada \\ \hline
OWASP ModSecurity (Zenodo) \cite{lucz2025} & 3{,}860 & 1.01 & CC BY 4.0, auditada \\ \hline
command injection.csv \cite{pawase2024cmdi} & 2{,}053 & 0.54 & No especificada por el origen, auditada \\ \hline
SecLists y PayloadsAllTheThings \cite{miessler2026seclists,swisskyrepo2026payloads} (Command Injection) & 2{,}470 & 0.64 & MIT; no consta en la auditoría \\ \hline
Conjuntos sintéticos propios (siete) & 2{,}557 & 0.67 & Propios \\ \hline
Telemetría curada del pipeline de reentrenamiento & 119 & 0.03 & Propia \\ \hline
\textbf{Total} & \textbf{383{,}202} & \textbf{100.00} & \\ \hline
\end{tabular}
\caption[Orígenes de las filas del corpus]{Orígenes de las filas del corpus etiquetado y estado de su licencia}
\label{tab:met-origenes}
\end{table}
{\footnotesize Nota. Elaboración propia con el conteo de la columna \texttt{\_source} del archivo de
vectores de características (repositorio \texttt{logSguarDian}, 10 de octubre de 2026) y con
\texttt{docs/dataset-audit.md}.}

\textbf{Procedimiento 1.3. Normalización y etiquetado determinista (OE2).} Un analizador por
fuente transforma cada registro crudo a un esquema canónico común (\texttt{CanonicalRequest}:
método, ruta, consulta, cuerpo y cabeceras), y un archivo de mapeo versionado
(\texttt{label\_map.yaml}) asigna la clase canónica a partir de la etiqueta original de cada
fuente. El etiquetado no involucra juicio manual por muestra: la misma entrada produce siempre la
misma etiqueta. Las decisiones de mapeo que no son equivalencias conceptuales se declaran como
tales; por ejemplo, la asignación de CAPEC-242 (Code Injection) a la clase \texttt{xss}.

\textbf{Procedimiento 1.4. Deduplicación.} Se eliminan los registros idénticos mediante un hash
SHA-256 de la concatenación \texttt{ruta|consulta|cuerpo|etiqueta}. Los casi-duplicados
se detectan de dos formas y se reportan, pero no se eliminan: por distancia de Levenshtein \cite{levenshtein1966binary} menor
que 3 sobre la consulta, dentro de la misma clase, y por MinHash \cite{broder1997resemblance} con LSH \cite{indyk1998approximate} (128 permutaciones y
bandas de 8 por 16), que se aplicó a las cinco clases. La decisión de marcar sin eliminar se
declara como amenaza a la validez (Sección~\ref{sec:met-validez}).

\textbf{Procedimiento 1.5. Partición y bloqueo del conjunto de prueba (OE2, OE3).} El corpus
unificado se divide 70/15/15 en entrenamiento, validación y prueba, con estratificación conjunta
por clase y fuente y semilla fija. Antes de escribir las particiones, el script verifica tres
condiciones: cero solapamiento de hashes entre particiones, distribución de clases por partición
dentro de $\pm$2~\% de la global y distribución de fuentes dentro de $\pm$5~\%. Si se cumplen,
calcula un hash SHA-256 sobre los hashes ordenados de las filas de prueba y lo versiona
(\texttt{test.lock.sha256}) antes de cualquier entrenamiento. Finalidad: que cualquier lectura
posterior del conjunto de prueba pueda demostrar que la partición leída es la que se bloqueó.

\subsection{Fase 2: entrenamiento del modelo (entrada del componente de IA)}
\label{subsec:met-f2}

El entrenamiento de un Random Forest como clasificador supervisado y de un Isolation Forest como
detector de anomalías, ambos con scikit-learn \cite{pedregosa2011}, corresponde al componente de
IA. Para este componente se fijó el siguiente contrato de interfaz, que es lo que el middleware
consume:

\begin{itemize}
    \item Vector de características calculado por una única implementación en TypeScript
    (\texttt{packages/extractor}), consumida tanto por el pipeline de entrenamiento como por el
    middleware, para eliminar la divergencia entre entrenamiento y producción. El vector tiene 76
    características. En la generación vigente (\texttt{rf\_v11}/\texttt{if\_v10}), el Random Forest
    consume 69 y el Isolation Forest 63; en la generación publicada (\texttt{rf\_v10}/\texttt{if\_v9}),
    67 y 61. Se excluyen por nombre las que no están disponibles en el momento de inspeccionar la
    petición (código de estado y cinco características temporales) y
    \texttt{non\_json\_quote\_count}, agregada al extractor después del entrenamiento de los
    modelos vigentes. El Isolation Forest excluye además seis características con varianza nula o
    casi nula en tráfico legítimo.
    \item Modelos serializados en ONNX y ejecutados con \texttt{onnxruntime-node}
    \cite{onnxruntimedocs}, con una prueba de paridad obligatoria (diferencia máxima entre las
    predicciones de Python y de Node.js menor que 0.1~\%) antes de incorporarlos al paquete.
    \item Cada evaluación se ata a un par concreto de binarios ONNX, identificado por su hash
    SHA-256. Durante el proyecto se usaron tres pares, descritos en el
    Cuadro~\ref{tab:pares-modelos}.
\end{itemize}

\begin{table}[H]
\centering
\small
\begin{tabular}{|L{0.9cm}|L{2.6cm}|L{1.4cm}|L{2.6cm}|L{5.2cm}|}
\hline
\textbf{Par} & \textbf{Versión} & \textbf{Caract. (RF/IF)} & \textbf{Hash (\texttt{rf.onnx} / \texttt{if.onnx})} & \textbf{Evaluaciones en que se usó} \\ \hline
A & \texttt{rf\_v10} / \texttt{if\_v9} (2 de agosto de 2026) & 67 / 61 & \texttt{18083725...} / \texttt{5de374a7...} & Paquete \texttt{logsguardian@0.1.0} (commit \texttt{c3d591a}); Ronda~4 (Configuraciones 2, 3a y 3b sobre los 590 payloads)\textsuperscript{c}; medición principal de la Ronda~5; estudio de \texttt{User-Agent}\textsuperscript{a}; DVNA, node-api-goat y dvws-node (atribución probable)\textsuperscript{b} \\ \hline
B & \texttt{rf\_v11} / \texttt{if\_v10}, primer build (30 de agosto de 2026) & 69 / 63 & \texttt{a7c015f5...} / \texttt{3cf8e45c...} & Repetición de la Configuración~2 sobre los 590 payloads de la Ronda~4 (9 de septiembre de 2026); benchmark de latencia, memoria y CPU en Linux nativo \\ \hline
C & \texttt{rf\_v11} / \texttt{if\_v10}, reentrenamiento (26 de septiembre de 2026) & 69 / 63 & \texttt{25b407e6...} / \texttt{5be28ee8...} & Paquete \texttt{logsguardian@0.1.1} (publicado el 11 de octubre de 2026)\textsuperscript{d}, con el que se verificó la cobertura de Path Traversal y de Command Injection de la Ronda~5 con el ejecutor corregido; lectura del conjunto de prueba, paridad ONNX y análisis SHAP; repetición de la Ronda~5 con el tarball de validación 22550ea; OWASP Juice Shop (copia local) \\ \hline
\end{tabular}
\caption[Pares de modelos ONNX y su uso]{Pares de modelos ONNX utilizados y evaluaciones asociadas a cada uno}
\label{tab:pares-modelos}
\end{table}
{\footnotesize Nota. Elaboración propia con datos de \texttt{docs/model-provenance.md} (repositorio
\texttt{logSguarDian}), verificados por hash contra el historial de git y contra el paquete
descargado del registro de npm. \textsuperscript{a}El archivo de resultados no registra el modelo; la
atribución se apoya en el contenido de \texttt{models/} en el commit del estudio (\texttt{5966712}),
que contiene el par~A (hash verificado). \textsuperscript{b}Los tres proyectos declaran la dependencia
\texttt{\^{}0.1.0} del registro de npm, cuyo único contenido publicado en la fecha de las mediciones
(30 de septiembre de 2026) era el par~A, pero el paquete y el modelo exactos de cada medición no
constan en el repositorio. \textsuperscript{c}En las Rondas 1 a 3,
con un corpus pequeño, el modelo no consta en los resultados. \textsuperscript{d}Contenido verificado
descargando el paquete del registro: modelos con los hashes del par~C, extractor de 76
características, corrección de la recodificación de la consulta y umbral del Isolation Forest leído del archivo \texttt{parity\_report.json} del modelo (0.004206). No se agrega un par~D: el paquete 0.1.1 y el tarball
de validación contienen los mismos modelos, y los pares se identifican por el hash de sus archivos
ONNX; difieren en el middleware, no en el modelo.}

El paquete \texttt{logsguardian@0.1.0} se publicó en npm el 30 de agosto de 2026 desde el commit
\texttt{c3d591a} y contiene el par~A; el primer build de \texttt{rf\_v11}/\texttt{if\_v10} (par~B) se
generó ese mismo día y no se incluyó en él. La versión \texttt{0.1.1}, publicada el 11 de octubre de
2026 (SHA-256 del archivo \texttt{abf8d5c4...}), contiene el par~C, el extractor de 76
características, la corrección del defecto de recodificación de la consulta y el umbral del Isolation Forest leído del archivo \texttt{parity\_report.json} del modelo. En este
documento, \textit{paquete 0.1.0}
designa a la primera y \textit{paquete 0.1.1} a la segunda. Con la 0.1.1 solo se midió la cobertura de la Configuración~2 de la Ronda~5, con el ejecutor
original y con el corregido; las aplicaciones de terceros y el rendimiento no se midieron con ella,
que es posterior a esas pruebas. Las mediciones de la Ronda~5 designadas como del paquete 0.1.0
se hicieron con un tarball construido en el commit \texttt{a3aada0} (SHA-256 \texttt{e09120006dd1...}),
que contiene el par~A pero no es el archivo publicado en el registro (SHA-256
\texttt{c5c1a1da57cf...}); se les llama paquete 0.1.0 por contener el mismo par de modelos.
Para medir la cobertura del modelo vigente antes de que existiera la 0.1.1, se construyó un
paquete a partir del commit \texttt{22550ea} de la rama \texttt{develop}, con el par~C y sin la
corrección, denominado \textit{tarball de validación 22550ea} (SHA-256 \texttt{cf1238de...}). Ese
paquete no se publicó y no debe confundirse con la 0.1.0 ni con la 0.1.1. La evaluación de
cobertura sobre el corpus limpio se ejecutó con el paquete 0.1.0 y con el tarball de validación,
de modo que esa métrica existe para los pares A y C, y para la versión 0.1.1 publicada solo en la Configuración~2.

El pipeline de reentrenamiento offline con compuertas automáticas (\texttt{packages/mlops}) se
incorporó por recomendación del asesor como parte del componente de IA. No modifica el
comportamiento de la librería publicada: los modelos se entrenan fuera de línea y se distribuyen
fijos con cada versión, sin aprendizaje a partir del tráfico del usuario.

\textbf{Procedimiento 2.1. Calibración del umbral de bloqueo (OE3, ciberseguridad).} La
política de decisión otorga al Random Forest la autoridad exclusiva de bloqueo: una petición se
bloquea cuando la probabilidad de la clase de ataque más probable es igual o mayor que el umbral
\texttt{RF\_THRESHOLD}. El Isolation Forest nunca bloquea; solo marca la petición como anómala
en el registro (\texttt{pass\_anomaly}), porque una política disyuntiva (bloquear si cualquiera
de los dos modelos alerta) habría bloqueado cerca del 10~\% del tráfico legítimo. El umbral se
calibró con el modelo \texttt{rf\_v3} mediante un barrido entre 0.10 y 0.70 sobre la partición
de validación, nunca sobre la de prueba, registrando para cada valor la precisión, el recall, los
falsos positivos y la detección por clase. Se seleccionó 0.35, el extremo conservador de la
meseta en que recall y falsos positivos permanecen constantes. El repositorio registra una
confirmación de ese umbral sobre el conjunto de prueba de \texttt{rf\_v3}, pero no conserva el
artefacto de esa lectura, por lo que esa confirmación no es verificable. Para \texttt{rf\_v11} el
umbral se mantuvo sin recalibrar. El umbral del Isolation Forest se seleccionó también sobre
validación y se fija en los metadatos de cada modelo; sin embargo, en dos generaciones previas
(\texttt{if\_v2} e \texttt{if\_v5}) el objetivo interno de calibración (falsos positivos de 8~\% y
luego de 6~\%) se definió después de observar un resultado de prueba que no lo cumplía.

\textbf{Procedimiento 2.2. Lectura del conjunto de prueba (OE3).} La evaluación final de los
modelos publicados ejecuta un script versionado (\texttt{evaluate\_test.py}) que: (1) recalcula el
hash del conjunto de prueba y aborta si no coincide con el bloqueado; (2) carga los binarios ONNX
del repositorio (par C), no los modelos de Python; (3) calcula el reporte de clasificación y el AUC-ROC
por clase en esquema uno contra el resto, asignando a cada petición la clase de mayor
probabilidad (sin aplicar \texttt{RF\_THRESHOLD}); (4) evalúa el Isolation Forest con el umbral de
sus metadatos; y (5) guarda un manifiesto de sumas de verificación del conjunto de prueba y de
ambos binarios junto al resultado. Finalidad: que cada cifra del OE3 quede atada al contenido
exacto, byte a byte, del modelo y de los datos que la produjeron.

Para la generación \texttt{rf\_v11} el conjunto de prueba se leyó tres veces: el 24 de septiembre
de 2026, con un modelo entrenado antes del rebloqueo del 23 de septiembre y un umbral del
Isolation Forest de 0.00807; y el 25 y el 26 de septiembre, con modelos reentrenados sobre la
partición vigente y un umbral de 0.00421. Las cifras reportadas en Resultados corresponden a la
lectura del 26 de septiembre (par C).

\subsection{Fase 3: desarrollo de la librería}
\label{subsec:met-f3}

\textbf{Procedimiento 3.1. Captura y respuesta.} El middleware transforma el objeto
\texttt{req} de Express al esquema canónico, extrae el vector de características y envía la
inferencia a \texttt{worker\_threads}, para no bloquear el Event Loop. La arquitectura de esa
inferencia en hilos (un hilo dedicado al Random Forest, un grupo rotativo de dos hilos para el
Isolation Forest, y una respuesta que se resuelve con el veredicto del Random Forest sin esperar
al Isolation Forest, cuya respuesta tardía actualiza el registro de forma asíncrona) es una
decisión de diseño del componente de Inteligencia Artificial, justificada en el trabajo de
graduación de Diego Valenzuela \cite{valenzuela2026ia}; aquí se describe solo como la interfaz que el middleware consume.
Ante una petición bloqueada, el middleware responde HTTP 403 en
modo \texttt{block}; en modo \texttt{monitor} solo registra. Todo evento se persiste en SQLite y
los eventos accionables (\texttt{block}, \texttt{pass\_anomaly}) se notifican a los webhooks
registrados. Ante un error o un tiempo de espera excedido del hilo de inferencia (50~ms por
defecto), el middleware deja pasar la petición (\textit{fail-open}), para que un fallo interno
no convierta la librería en un vector de denegación de servicio contra la aplicación anfitriona.

\textbf{Procedimiento 3.2. Verificación funcional.} Cada módulo se valida con pruebas
automatizadas en Jest \cite{jestdocs} y Supertest \cite{supertestdocs} que se ejecutan en
integración continua. Las correcciones de defectos se acompañan de una prueba de regresión y se
verifican mediante sabotaje y reversión: se reintroduce el defecto de forma temporal, se
confirma que la prueba nueva falla y se revierte. Antes de la publicación, el paquete empaquetado
(\texttt{pnpm pack}) se instaló en un directorio limpio fuera del monorepo y se probó contra una
aplicación Express real con una petición SQLi (respuesta 403) y una legítima (respuesta 200).

\subsection{Fase 4: pruebas}
\label{subsec:met-f4}

Las pruebas se ejecutan sobre una aplicación de referencia deliberadamente vulnerable
(\texttt{logSguarDian-vulnerable-project}), construida con Express, PostgreSQL y plantillas EJS,
y desplegada con Docker Compose \cite{dockerdocs, moreau2023}. La aplicación expone endpoints
vulnerables a cada vector: \texttt{GET /posts/search?q=} y \texttt{POST /login} (SQLi),
\texttt{POST /posts} y \texttt{POST /profile} (XSS), \texttt{GET /posts/:id/attachment?file=}
(Path Traversal/LFI) y \texttt{POST /admin/ping} (Command Injection).

\textbf{Procedimiento 4.1. Configuraciones experimentales (OE3).} La misma aplicación se evalúa
en cuatro configuraciones que solo difieren en la capa de defensa:

\begin{itemize}
    \item \textbf{Configuración 1 (control):} aplicación sin protección, para confirmar que cada
    payload es explotable y establecer la línea base.
    \item \textbf{Configuración 2:} aplicación con \textit{logSguarDian} en modo \texttt{block}.
    \item \textbf{Configuración 3a:} aplicación detrás de Nginx con ModSecurity v3 \cite{owasp2026modsecurity} y OWASP CRS
    \cite{owaspcrs} en nivel de paranoia 1 y modo de bloqueo, sin \textit{logSguarDian}.
    \item \textbf{Configuración 3b:} ambas capas superpuestas. El registro de auditoría de
    ModSecurity y el registro de eventos de \textit{logSguarDian} se cruzan para atribuir cada
    bloqueo a la capa que lo produjo, dado que ambas responden con un código de bloqueo.
\end{itemize}

Se eligió el nivel de paranoia 1 porque es la configuración realista para la población
objetivo: organizaciones sin equipo de seguridad que pueda ajustar los falsos positivos de los
niveles superiores.

\textbf{Procedimiento 4.2. Auditoría de fuga entre evaluación y entrenamiento.}
\label{subsec:met-fuga}
La fuga de datos, es decir, la presencia en el conjunto de evaluación de información disponible
durante el entrenamiento, es una causa frecuente de sobreestimación del desempeño en estudios
basados en aprendizaje automático \cite{kapoor2023}. Un script versionado
(\texttt{leakage\_check.py}) cruza cada corpus de evaluación contra cada partición en dos
niveles: coincidencia exacta (hash del payload) y coincidencia de plantilla. La plantilla
normaliza la decodificación URL, las mayúsculas y los dígitos; para Command Injection reemplaza
además el marcador aleatorio de SecLists por un símbolo fijo. La auditoría confirmó que el corpus
de la Ronda~4 tenía fuga en Command Injection (los 200 payloads coinciden por plantilla con la
fuente SecLists usada en entrenamiento) y que el conjunto de 500 payloads de la suite E2E se
había muestreado del corpus completo (entre 70 y 75 coincidencias exactas con entrenamiento por
cada 100 payloads de cada clase). En consecuencia, la detección de payloads no vistos se mide
exclusivamente con el corpus limpio M2; el conjunto E2E se usa para verificar el funcionamiento
extremo a extremo de la librería.

\textbf{Procedimiento 4.3. Construcción del corpus limpio (OE3).} El corpus M2 se genera con un
script versionado (\texttt{build\_corpus.py}) y semilla fija (20260923), de modo que dos
ejecuciones sobre las mismas entradas producen el mismo corpus:

\begin{itemize}
    \item \textbf{SQLi:} sqlmap \cite{sqlmap} v1.10.9, con los parámetros \texttt{--batch
    --level=3 --risk=2 --technique=BEUSTQ}, contra los dos endpoints SQLi. Se registra cada
    petición enviada mediante su archivo de tráfico; se usan las formas de payload que sqlmap
    intentó, no su veredicto de explotabilidad, porque los bloqueos 403 del middleware alteran
    sus heurísticas.
    \item \textbf{Path Traversal/LFI:} generador propio que combina profundidades de 2 a 12
    niveles, ocho estilos de codificación (crudo, barra invertida, codificación URL simple, doble
    y de mayúsculas mixtas, UTF-8 sobrelargo, evasión \texttt{....//} y sufijo de byte nulo) y 28
    archivos objetivo de Linux, \texttt{/proc}, la propia aplicación y Windows.
    \item \textbf{Command Injection:} generador propio, deliberadamente no derivado de SecLists,
    que combina 12 separadores de shell y 4 envolturas de sustitución con 20 comandos de
    post-explotación sobre el sumidero \texttt{ping -c 3 \$\{host\}}.
    \item \textbf{XSS:} OWASP ZAP \cite{owaspzap}, con un \textit{Active Scan} dirigido por su API REST
    contra \texttt{POST /posts} (campos \texttt{content} y \texttt{title}) y \texttt{POST /profile}
    (campo \texttt{bio}), con sesión autenticada. ZAP reporta una alerta por cada combinación de
    parámetro y técnica confirmada como vulnerable, no cada payload probado; por eso el corpus de
    XSS contiene un solo payload, el confirmado, y no se completó con variantes escritas a mano.
\end{itemize}

Todo candidato pasa por un filtro de exclusión antes de entrar al corpus: se descarta si su hash
(o el de su forma codificada o decodificada) coincide con algún payload de entrenamiento de la
misma clase o, en Command Injection, si su plantilla coincide con alguna plantilla de la fuente
SecLists o de la clase \texttt{cmdi} de entrenamiento. Los payloads coincidentes se descartan, no
se marcan, y cada descarte queda registrado con su motivo y la fuente coincidente. De cada
categoría de SQLi, Path Traversal y Command Injection se seleccionan 200 payloads finales; el corpus
resultante tiene 601 peticiones.

\textbf{Procedimiento 4.4. Medición de cobertura (OE3).} Un ejecutor (\texttt{run\_large\_corpus.py})
envía cada payload del corpus a su endpoint objetivo con la sesión correspondiente y registra, por
petición, el código de respuesta y, cuando aplica, la clase y la confianza del evento almacenado.
Una petición cuenta como detectada solo si el middleware la bloquea con HTTP 403. Las que el
Isolation Forest marca como anómalas sin bloquearlas (\texttt{pass\_anomaly}), las que terminan por
tiempo de espera y las que devuelven HTTP 500 cuentan como no detectadas. La medición principal usa
el paquete 0.1.0 (par~A), con el ejecutor en la versión anterior a la corrección del 10 de octubre
de 2026; esa corrección evita volver a codificar los payloads de Path Traversal y de Command
Injection que ya llegan codificados, y desde entonces el ejecutor envía sus bytes tal cual. Como
referencia, se repitió sobre el mismo corpus con el tarball de
validación 22550ea (par~C), verificando dentro del contenedor el hash de los modelos cargados. Las
cuatro configuraciones del Procedimiento~4.1 se ejecutaron sobre el corpus limpio; la
Configuración~3a, que no incorpora la librería, se ejecutó en una rama distinta de la aplicación
mediante un árbol de trabajo de git, y de las Configuraciones 3a y 3b no consta el paquete exacto: el directorio de dependencias
de su commit contiene el archivo de la 0.1.0 del registro, pero los resultados no registran cuál
se ejecutó. La Configuración~2 se repitió además con el paquete 0.1.1, instalado desde el registro de npm, y con
el ejecutor corregido.

\textbf{Procedimiento 4.5. Benchmark de rendimiento (objetivo general).} El benchmark se ejecuta
en Linux nativo mediante un flujo de GitHub Actions (\texttt{latency-benchmark.yml}) en un
ejecutor \texttt{ubuntu-latest}, que documenta el hardware del ejecutor en cada corrida. La
aplicación de referencia se levanta con Docker Compose, y la condición se alterna entre línea
base y middleware activo mediante la variable \texttt{LOGSGUARDIAN\_DISABLED}, recreando el
contenedor, sin modificar ningún otro parámetro. Se usa Artillery 2.0.34 \cite{artillerydocs}
con tres escenarios:

\begin{itemize}
    \item \textbf{Normal:} navegación legítima (inicio de sesión, listado, detalle y perfil) a
    25 usuarios virtuales por segundo, que con cuatro peticiones por usuario produce
    $\approx$100 peticiones/s.
    \item \textbf{Ataque:} tráfico mixto concurrente, 70~\% de navegación legítima y 30~\% de
    escenarios con payloads de las cuatro categorías, a 30 usuarios virtuales por segundo
    ($\approx$100 peticiones/s).
    \item \textbf{Volumen:} el escenario normal sobre una base de datos sembrada con 2{,}010
    registros, para representar una aplicación con datos reales. Este escenario no estaba en el
    protocolo; se agregó para contrastar la hipótesis de que una línea base más lenta reduce el
    impacto relativo del middleware, y sus resultados se comparan con el mismo criterio de 15~ms.
\end{itemize}

Cada corrida tiene una fase de calentamiento de 15~s a 5 usuarios/s, descartada del análisis
porque los hilos del Isolation Forest tardan cerca de 5~s en quedar listos, y una fase medida de
60~s. Cada combinación de escenario y condición se repite cinco veces. Durante toda la ejecución,
un proceso paralelo muestrea cada 5~s la memoria y el uso de CPU del contenedor de la aplicación
mediante \texttt{docker stats}. Un script posterior (\texttt{analyze\_resources.py}) asigna cada
muestra a su repetición con las marcas de tiempo de inicio y fin de cada reporte de Artillery,
descarta los primeros 15~s de cada repetición y calcula, para cada par de repeticiones
equivalentes, la diferencia entre la condición con middleware y la línea base. La evidencia
cruda de la corrida (reportes de Artillery, muestras de recursos y hashes de los modelos) está
archivada en el repositorio de la aplicación de referencia. El procedimiento de agregación se
detalla en la Sección~\ref{sec:met-datos}.

\textbf{Procedimiento 4.6. Estudio de falsos positivos con navegadores reales (OE3).} El
procedimiento surgió de un falso positivo reportado al navegar con un navegador real hacia
\texttt{GET /login}. Las corridas de Artillery usan una única cadena \texttt{User-Agent} fija y no
cubren ese caso. Un script (\texttt{ua\_fp\_study.py}) envía peticiones sin consulta ni cuerpo a
\texttt{/login}, \texttt{/posts} y \texttt{/profile} con cadenas \texttt{User-Agent} de Chrome,
Firefox y Safari en macOS y Windows, con y sin sesión, y registra el código de respuesta de cada
combinación.

\textbf{Procedimiento 4.7. Validación de generalización en aplicaciones de terceros (objetivo
general).} Sustituye a la validación con startups guatemaltecas prevista en el protocolo y se
realizó por recomendación del asesor. Se evaluaron cuatro aplicaciones Node.js deliberadamente vulnerables que no participaron en la
calibración ni en el entrenamiento: DVNA, node-api-goat, dvws-node y OWASP Juice Shop. En las tres
primeras se integró el paquete publicado en npm (\texttt{npm install logsguardian}, entonces 0.1.0);
OWASP Juice Shop la evaluó el componente de IA con una copia local del paquete que contiene el
par~C. Las cuatro difieren en arquitectura (renderizado de
plantillas, API pura sin base de datos, SOAP/GraphQL y aplicación de página única). En las tres primeras se confirma primero la explotabilidad sin protección y luego se mide, con el
middleware en modo \texttt{block} y configuración por defecto, la tasa de bloqueo sobre ataques y la tasa de bloqueo
sobre tráfico legítimo representativo de la aplicación. En Juice Shop, en modo \texttt{monitor},
se aisló el efecto de la cabecera \texttt{User-Agent} variando solo esa cabecera (cURL frente a
una cadena de navegador) sobre el mismo conjunto de peticiones.

\subsection{Fase 5: aplicación de feedback}
\label{subsec:met-f5}
Cada hallazgo de la Fase~4 que afecta a una métrica se convierte en una tarea del tablero, con
su propia rama y pull request, y sigue un ciclo fijo: reproducción con evidencia antes de
cualquier cambio; corrección; nueva medición de la métrica afectada; y registro del valor antes
y después en la bitácora de decisiones. Si una corrección requiere reentrenar el modelo, el
conjunto de prueba se bloquea de nuevo antes del reentrenamiento (Procedimiento~1.5). Una
corrección no puede justificarse con la misma lectura de prueba que la motivó.

\subsection{Fase 6: documentación}
\label{subsec:met-f6}
Se consolida un documento de cumplimiento por objetivo, en el que cada afirmación remite al
artefacto que la respalda (script, archivo de resultados, ejecución de CI o captura) y cada
cifra histórica superada se marca como tal. El informe final, el README publicado en npm y la
documentación técnica se derivan de ese documento.

% ------------------------------------------------------------------------------
\section{Técnicas e instrumentos}
\label{sec:met-instrumentos}

El Cuadro~\ref{tab:met-instrumentos} relaciona cada técnica de obtención de datos con el
instrumento que la implementa y con el dato que produce.

\begin{longtable}{|L{2.7cm}|L{3.7cm}|L{3.6cm}|L{2.6cm}|}
\hline
\textbf{Técnica} & \textbf{Instrumento} & \textbf{Dato producido} & \textbf{Procedimiento} \\ \hline
\endfirsthead
\multicolumn{4}{l}{\footnotesize\textit{(continuación)}} \\ \hline
\textbf{Técnica} & \textbf{Instrumento} & \textbf{Dato producido} & \textbf{Procedimiento} \\ \hline
\endhead
\multicolumn{4}{r}{\footnotesize\textit{Continúa en la página siguiente}} \\
\endfoot
\endlastfoot
Análisis documental & Páginas oficiales OWASP, CWE y MITRE ATT\&CK; registros de licencia
de cada fuente & Mapeo de vectores; tabla de auditoría de fuentes & 1.1, 1.2 \\ \hline
Transformación determinista & Analizadores por fuente, \texttt{label\_map.yaml},
\texttt{unify.py}, \texttt{split.py} & Corpus etiquetado, particiones y hash de bloqueo & 1.3--1.5 \\ \hline
Extracción de características & CLI de \texttt{packages/extractor} (TypeScript) & Matriz de
características por petición & 1.3, F2 \\ \hline
Evaluación sobre conjunto bloqueado & \texttt{evaluate\_test.py}, scikit-learn
\cite{pedregosa2011}, \texttt{onnxruntime} & Precisión, recall, F1, AUC-ROC; métricas del
Isolation Forest; manifiesto de sumas & 2.2 \\ \hline
Barrido de umbral & Notebooks de calibración sobre validación & Curvas precisión/recall/FP por
umbral & 2.1 \\ \hline
Pruebas automatizadas & Jest, Supertest, GitHub Actions & Resultado de pruebas por commit & 3.2 \\ \hline
Auditoría de fuga & \texttt{leakage\_check.py} & Coincidencias exactas y por plantilla por
clase y partición & 4.2 \\ \hline
Pruebas de penetración controlada & sqlmap \cite{sqlmap}, generadores propios, filtro de
exclusión y ejecutor del corpus & Corpus limpio; veredicto por
payload y configuración & 4.3, 4.4 \\ \hline
WAF de comparación & ModSecurity v3 + OWASP CRS PL1 \cite{owaspcrs} & Registro de auditoría
del WAF & 4.1 \\ \hline
Prueba de carga & Artillery 2.0.34, script de orquestación de la suite y flujo de CI del
benchmark & Latencia p95 por repetición; errores y tiempos de espera
excedidos & 4.5 \\ \hline
Monitoreo de recursos & \texttt{docker stats} y script de análisis de memoria & Serie temporal de memoria y CPU del contenedor & 4.5 \\ \hline
Experimento de cabecera & Script de estudio de \texttt{User-Agent}; corridas cURL frente a
navegador en Juice Shop & Veredicto por cadena \texttt{User-Agent} y ruta & 4.6, 4.7 \\ \hline
Registro de eventos & SQLite del middleware; CLI \texttt{attacks} y \texttt{endpoints} & Clase,
confianza y veredicto por petición & 4.4, 4.7 \\ \hline
\caption[Técnicas e instrumentos de obtención de datos]{Técnicas e instrumentos de obtención de datos, dato producido y procedimiento asociado}
\label{tab:met-instrumentos}
\end{longtable}
{\footnotesize Nota. Elaboración propia.}

OWASP ZAP \cite{owaspzap}, previsto en el protocolo para generar payloads XSS, está instalado como
aplicación gráfica; no se automatizó mediante el navegador, sino a través de la API REST que expone
su proxy, que permitió crear el contexto, autenticar al usuario y lanzar el \textit{Active Scan}
sin intervención manual. Su resultado es un único payload confirmado, por lo que la cobertura de
XSS tiene n=1 (Sección~\ref{sec:met-restricciones}). La herramienta
clinic.js, prevista para medir CPU, se sustituyó por \texttt{docker stats}: mide el contenedor
completo, puede ejecutarse sin intervención dentro del flujo de CI y produce una serie temporal
alineada con las corridas de Artillery.

% ------------------------------------------------------------------------------
\section{Recursos}
\label{sec:met-recursos}

\subsection{Recursos humanos}
El componente de ciberseguridad estuvo a cargo de un único integrante, con dedicación estimada
de 20 horas semanales en las Fases 1 a 3 y de 15 horas semanales en las Fases 4 a 6. El
componente de IA estuvo a cargo del segundo integrante, que además actuó como revisor obligatorio
de los pull requests. La asesoría académica estuvo a cargo del asesor asignado, cuyas revisiones
del documento de cumplimiento motivaron varias de las correcciones descritas en la
Sección~\ref{sec:met-cambios}. No participó personal externo.

\subsection{Recursos tecnológicos}
El Cuadro~\ref{tab:met-recursos} actualiza los recursos tecnológicos del protocolo con las
herramientas que se usaron en la práctica.

\begin{longtable}{|L{3.0cm}|L{4.5cm}|L{5.3cm}|}
\hline
\textbf{Categoría} & \textbf{Herramienta} & \textbf{Función en el proyecto} \\ \hline
\endfirsthead
\multicolumn{3}{l}{\footnotesize\textit{(continuación)}} \\ \hline
\textbf{Categoría} & \textbf{Herramienta} & \textbf{Función en el proyecto} \\ \hline
\endhead
\multicolumn{3}{r}{\footnotesize\textit{Continúa en la página siguiente}} \\
\endfoot
\endlastfoot
Runtime y framework & Node.js, Express & Plataforma de la librería y de las aplicaciones de prueba \\ \hline
Monorepo & pnpm workspaces, TypeScript & Paquetes \texttt{extractor}, \texttt{core} y \texttt{mlops} \\ \hline
Inferencia & \texttt{onnxruntime-node}, \texttt{worker\_threads} & Ejecución de los modelos sin bloquear el Event Loop \\ \hline
Persistencia & SQLite (\texttt{better-sqlite3}) & Registro de eventos y de webhooks \\ \hline
Entorno de prueba & Docker Compose, PostgreSQL & Aplicación de referencia reproducible \\ \hline
Generación de tráfico & Artillery 2.0.34 & Carga legítima y mixta \\ \hline
Simulación de ataques & sqlmap 1.10.9; generadores propios en Python & Payloads de SQLi, Path Traversal y Command Injection \\ \hline
WAF de comparación & Nginx + ModSecurity v3 + OWASP CRS & Configuraciones 3a y 3b \\ \hline
Análisis de datos y ML & Python, pandas, scikit-learn, Jupyter & Partición, evaluación, calibración y auditoría de fuga \\ \hline
Pruebas & Jest, Supertest & Pruebas unitarias, de integración y E2E \\ \hline
Integración continua & GitHub Actions (\texttt{ubuntu-latest}) & Pruebas por pull request, publicación en npm y benchmark de rendimiento en Linux nativo \\ \hline
Gestión & Notion (tablero Kanban), Git, GitHub & Gestión de tareas, versionado y revisión cruzada \\ \hline
Distribución & Registro público de npm & Publicación de \texttt{logsguardian} \\ \hline
\caption[Recursos tecnológicos utilizados]{Recursos tecnológicos utilizados en el proyecto}
\label{tab:met-recursos}
\end{longtable}
{\footnotesize Nota. Elaboración propia. Todas las herramientas son de código abierto o de uso
gratuito para el proyecto.}

\subsection{Recursos externos y financieros}
Los recursos externos son las fuentes públicas del corpus, las cuatro aplicaciones vulnerables
de terceros, el registro de npm y los ejecutores de GitHub Actions. El componente de
ciberseguridad no requirió licencias de pago. La infraestructura en la nube contemplada en el
protocolo para la validación con startups no se utilizó, porque esa validación no se realizó.

% ------------------------------------------------------------------------------
\section{Recolección, procesamiento y validación de datos}
\label{sec:met-datos}

El Cuadro~\ref{tab:met-datos} describe, para cada flujo de datos, cómo se recolecta, cómo se
procesa y qué verificación debe superar antes de considerarse evidencia válida para el capítulo
de Resultados.

\begin{longtable}{|L{2.3cm}|L{3.4cm}|L{3.4cm}|L{3.7cm}|}
\hline
\textbf{Flujo} & \textbf{Recolección} & \textbf{Procesamiento} & \textbf{Validación} \\ \hline
\endfirsthead
\multicolumn{4}{l}{\footnotesize\textit{(continuación)}} \\ \hline
\textbf{Flujo} & \textbf{Recolección} & \textbf{Procesamiento} & \textbf{Validación} \\ \hline
\endhead
\multicolumn{4}{r}{\footnotesize\textit{Continúa en la página siguiente}} \\
\endfoot
\endlastfoot
Corpus etiquetado (M1) & Descarga de las fuentes públicas (Cuadro~\ref{tab:met-origenes}) & Esquema canónico,
etiquetado por mapeo, deduplicación exacta, extracción de características, partición
estratificada & Cero solapamiento entre particiones; tolerancias de clase ($\pm$2~\%) y fuente
($\pm$5~\%); hash de prueba versionado antes de entrenar \\ \hline
Métricas de clasificación & Lectura del conjunto de prueba con los binarios ONNX publicados
& Reporte de clasificación y AUC-ROC uno contra el resto & Hash de prueba coincidente;
manifiesto de sumas de \texttt{test.parquet}, \texttt{rf.onnx} e \texttt{if.onnx}; umbrales
tomados de los metadatos del modelo \\ \hline
Paridad ONNX & Predicciones de Python y de Node.js sobre el mismo conjunto & Diferencia máxima
absoluta por salida & Prueba automatizada en CI; diferencia $<$0.1~\% \\ \hline
Corpus limpio (M2) & sqlmap y generadores propios con semilla fija & Filtro de exclusión
exacto y por plantilla; selección de 200 por categoría & Reporte de exclusión con cero
coincidencias en el corpus final; reproducibilidad (dos ejecuciones con igual resultado) \\ \hline
Cobertura & Envío por HTTP real del corpus a la aplicación en cada configuración & Conteo de
bloqueos (HTTP 403) por categoría; anomalías sin bloqueo, tiempos de espera y errores 500 cuentan
como no detectados & Explotabilidad confirmada en Configuración~1; atribución por
capa en 3b mediante el registro de auditoría del WAF \\ \hline
Latencia & Reportes JSON de Artillery por repetición & p95 por repetición sobre la fase medida;
mediana de las cinco repeticiones; $\Delta$p95 absoluto y relativo & Exclusión del
calentamiento; tasa real de peticiones $\approx$100/s verificada en el reporte; conteo de
tiempos de espera excedidos a partir del resumen final de cada repetición \\ \hline
Memoria y CPU & Muestreo cada 5~s de \texttt{docker stats} & Asignación de muestras a cada
repetición por marca de tiempo; media por condición (MiB); incremento pareado con middleware menos línea
base & Ventana total $\geq$30~min; se excluyen calentamiento y muestras fuera de una repetición;
contraste entre la ventana medida y la repetición completa \\ \hline
Falsos positivos con navegador (M4) & Script con cadenas \texttt{User-Agent} reales &
Proporción de combinaciones bloqueadas & Sesión y ruta controladas; una sola variable manipulada \\ \hline
Aplicaciones de terceros (M5) & Integración del paquete publicado y scripts de ataque y de
tráfico legítimo por aplicación & Tasa de bloqueo en ataques y en tráfico legítimo, reportadas
siempre juntas & Explotabilidad previa confirmada; evidencia cruda (scripts y JSON) copiada al
repositorio \\ \hline
\caption[Recolección, procesamiento y validación de datos]{Recolección, procesamiento y validación de cada flujo de datos}
\label{tab:met-datos}
\end{longtable}
{\footnotesize Nota. Elaboración propia.}

Tres decisiones de procesamiento requieren justificación explícita, porque afectan la
interpretación de los resultados:

\begin{itemize}
    \item \textbf{Agregación de las repeticiones de latencia.} De las cinco repeticiones de cada condición se reporta la mediana. Se usa la mediana porque es menos sensible que la media a una repetición atípica
    causada por el entorno de ejecución.
    \item \textbf{Cálculo del p95 por repetición.} Artillery reporta percentiles por intervalos
    de aproximadamente 10~s. El analizador calcula el p95 de cada repetición como el promedio,
    ponderado por número de peticiones, de los p95 de cada intervalo posterior al calentamiento,
    no como el percentil exacto sobre todas las muestras. Es una aproximación y se declara como
    tal; su efecto es mayor en el escenario de volumen, donde la latencia crece durante la corrida.
    \item \textbf{Lectura conjunta de detección y falsos positivos.} En las aplicaciones de
    terceros, una tasa de bloqueo alta en ataques no se reporta sin la tasa de bloqueo en tráfico
    legítimo. Un sistema que bloquea casi todo obtiene un recall alto sin discriminar.
\end{itemize}

% ------------------------------------------------------------------------------
\section{Métricas e indicadores}
\label{sec:met-metricas}

Las métricas de clasificación se calculan por clase. Sean $TP_c$, $FP_c$ y $FN_c$ los
verdaderos positivos, falsos positivos y falsos negativos de la clase $c$ (Ecuación~\eqref{eq:met-clasificacion}):

\begin{equation}
P_c = \frac{TP_c}{TP_c + FP_c}, \qquad
R_c = \frac{TP_c}{TP_c + FN_c}, \qquad
F1_c = \frac{2\,P_c\,R_c}{P_c + R_c}
\label{eq:met-clasificacion}
\end{equation}

El F1 es la métrica primaria de aceptación porque el corpus está desbalanceado: la exactitud
global premiaría a un clasificador que favoreciera a la clase mayoritaria, mientras que el F1
penaliza tanto los ataques no detectados como los bloqueos de tráfico legítimo. El AUC-ROC uno
contra el resto se reporta como medida de discriminación independiente del umbral
\cite{fawcett2006}. Para el rendimiento (Ecuación~\eqref{eq:met-latencia}):

\begin{equation}
\Delta p95_{\mathrm{abs}} = \tilde{p}95_{\mathrm{activo}} - \tilde{p}95_{\mathrm{base}}, \qquad
\Delta p95_{\mathrm{rel}} = \frac{\Delta p95_{\mathrm{abs}}}{\tilde{p}95_{\mathrm{base}}} \times 100\,\%
\label{eq:met-latencia}
\end{equation}

\noindent donde $\tilde{p}95$ es la mediana de las cinco repeticiones de una condición.

Para las proporciones observadas (tasas de bloqueo y de detección) se reportan intervalos de confianza de Wilson al 95~\% \cite{wilson1927probable}, preferidos al intervalo normal (Wald) porque mantienen una cobertura cercana a la nominal con muestras pequeñas o proporciones cercanas a 0 o 1 \cite{brown2001interval}.

El criterio de aceptación de latencia es el de la Sección~\ref{subsec:alcance-tecnico}: un
$\Delta p95_{\mathrm{abs}}$ de 15~ms como máximo, tanto bajo carga normal como bajo carga de
ataque. Se expresa en milisegundos y no como porcentaje de la línea base porque un porcentaje
depende por completo de la latencia propia de la aplicación anfitriona. Con una línea base de 2 a
5~ms, como la de la aplicación de referencia, un margen de 5~\% equivale a una fracción de
milisegundo, menor que la variación entre repeticiones del entorno de medición, de modo que el
criterio dejaría de distinguir un middleware costoso de uno ligero. Además, el mismo middleware
obtendría veredictos opuestos según la aplicación en que se instale: en una aplicación más lenta,
el mismo costo se diluye y el porcentaje baja. Una cifra absoluta mide directamente el costo que
el middleware agrega a cada petición, que es la magnitud que percibe el usuario final y la que un
equipo sin especialista en seguridad puede comparar con su propio presupuesto de latencia.
$\Delta p95_{\mathrm{rel}}$ se reporta como indicador complementario.

El Cuadro~\ref{tab:met-metricas} define todas las métricas, su criterio y el objetivo al que se
vinculan.

\begin{longtable}{|L{2.8cm}|L{4.0cm}|L{2.7cm}|L{2.0cm}|L{1.0cm}|}
\hline
\textbf{Indicador} & \textbf{Definición operacional} & \textbf{Criterio} & \textbf{Instrumento} & \textbf{Obj.} \\ \hline
\endfirsthead
\multicolumn{5}{l}{\footnotesize\textit{(continuación)}} \\ \hline
\textbf{Indicador} & \textbf{Definición operacional} & \textbf{Criterio} & \textbf{Instrumento} & \textbf{Obj.} \\ \hline
\endhead
\multicolumn{5}{r}{\footnotesize\textit{Continúa en la página siguiente}} \\
\endfoot
\endlastfoot
Tamaño del corpus & Filas totales en las tres particiones & $\geq$100{,}000 & Partición (1.5) & OE2 \\ \hline
Representación de clases & Clases con muestras en cada partición & 5/5 en cada partición & Partición (1.5) & OE1, OE2 \\ \hline
Integridad de la partición & Hashes compartidos entre particiones & 0 & Partición (1.5) & OE2 \\ \hline
F1 por categoría & $F1_c$ sobre el conjunto de prueba bloqueado & $\geq$0.80 en al menos 3 de 4 categorías & Lectura de prueba (2.2) & OE3 \\ \hline
Precisión y recall por categoría & $P_c$, $R_c$ sobre el conjunto de prueba & Reportados, subordinados al F1 & Lectura de prueba (2.2) & OE3 \\ \hline
AUC-ROC por categoría & Área bajo la curva ROC uno contra el resto & Referencia, sin umbral & Lectura de prueba (2.2) & OE3 \\ \hline
Detección de anomalías del IF & $TP/(TP+FN)$ sobre ataques de prueba, umbral de los metadatos & $\geq$50~\% & Lectura de prueba (2.2) & OE3 \\ \hline
Falsos positivos del IF & $FP/(FP+TN)$ sobre legítimos de prueba & $\leq$10~\% & Lectura de prueba (2.2) & OE3 \\ \hline
Paridad ONNX & Diferencia máxima entre predicciones Python y Node.js & $<$0.1~\% & Prueba de paridad & OG \\ \hline
Cobertura por categoría & Payloads del corpus limpio bloqueados con HTTP 403 / payloads enviados & $\geq$80~\% por categoría & Ejecutor del corpus (4.4) & OE3 \\ \hline
Contribución en defensa en profundidad & Ataques no bloqueados por el WAF que bloquea \textit{logSguarDian} / ataques no bloqueados por el WAF & Descriptivo & Configuraciones 3a/3b & OG \\ \hline
Falsos positivos con navegador & Combinaciones ruta--\texttt{User-Agent} legítimas bloqueadas / combinaciones probadas & Descriptivo; se reporta junto al F1 & Estudio UA (4.6) & OE3 \\ \hline
Generalización a terceros & Tasa de bloqueo en ataques y en tráfico legítimo por aplicación & Descriptivo; siempre en pareja & Scripts por aplicación & OG \\ \hline
Extracción de características & Tiempo por petición, aislado en su hilo & $\leq$5~ms & Benchmark del extractor & OG \\ \hline
Latencia adicional & $\Delta p95_{\mathrm{abs}}$, escenarios normal y de ataque (criterio) y de volumen (referencia) & $\leq$15~ms; $\Delta p95_{\mathrm{rel}}$ complementario & Artillery, CI & OG \\ \hline
Tiempos de espera excedidos & Peticiones con \texttt{ERR\_SOCKET\_TIMEOUT} / peticiones enviadas & Descriptivo (escenario de volumen) & Artillery, CI & OG \\ \hline
Memoria adicional & Incremento pareado de memoria del contenedor con middleware respecto de la línea base (MiB) & $\leq$150~MB (143.1~MiB) & \texttt{docker stats} & OG \\ \hline
Incremento de CPU & Uso de CPU con middleware frente a la línea base & $\leq$15~\% relativo & \texttt{docker stats} & OG \\ \hline
\caption[Métricas e indicadores de evaluación]{Métricas e indicadores: definición operacional, criterio de aceptación, instrumento y objetivo vinculado}
\label{tab:met-metricas}
\end{longtable}
{\footnotesize Nota. Elaboración propia con base en el protocolo de graduación y en
\cite{fawcett2006}. OG: objetivo general.}

% ------------------------------------------------------------------------------
\section{Matriz de trazabilidad}
\label{sec:met-trazabilidad}

El Cuadro~\ref{tab:met-trazabilidad} relaciona cada objetivo con las actividades que lo
abordan, los datos que estas producen, las métricas que se calculan sobre esos datos y el
resultado esperado, que es lo que el capítulo de Resultados contrasta.

\begin{longtable}{|L{1.5cm}|L{2.6cm}|L{2.7cm}|L{2.6cm}|L{3.3cm}|}
\hline
\textbf{Objetivo} & \textbf{Actividades} & \textbf{Datos} & \textbf{Métricas} & \textbf{Resultado esperado} \\ \hline
\endfirsthead
\multicolumn{5}{l}{\footnotesize\textit{(continuación)}} \\ \hline
\textbf{Objetivo} & \textbf{Actividades} & \textbf{Datos} & \textbf{Métricas} & \textbf{Resultado esperado} \\ \hline
\endhead
\multicolumn{5}{r}{\footnotesize\textit{Continúa en la página siguiente}} \\
\endfoot
\endlastfoot
OE1 & 1.1 Mapeo a OWASP, CWE y MITRE & Cuadro de correspondencia verificado & Vectores mapeados con fuente oficial & Cuatro vectores justificados; cinco clases canónicas que estructuran OE2 y OE3 \\ \hline
OE2 & 1.2--1.5 Auditoría, etiquetado, deduplicación y partición & Corpus M1; reporte de auditoría; hash de bloqueo & Tamaño; representación de clases; integridad de partición & $\geq$100{,}000 muestras etiquetadas con etiquetado determinista y documentado \\ \hline
OE3 (modelo) & 2.1 Calibración; 2.2 Lectura de prueba & Conjunto de prueba bloqueado; manifiesto de sumas & F1, precisión, recall y AUC-ROC por clase; detección y FP del IF & F1 $\geq$0.80 en al menos 3 de 4 categorías; IF $\geq$50~\% de detección con $\leq$10~\% de FP \\ \hline
OE3 (sistema) & 4.1--4.4 Configuraciones, auditoría de fuga, corpus limpio y cobertura & Corpus M2; veredictos por payload y configuración & Cobertura por categoría; contribución en defensa en profundidad & Cobertura $\geq$80~\% por categoría sobre payloads no vistos en entrenamiento \\ \hline
OE3 (falsos positivos) & 4.6 Estudio con navegadores & Muestra M4 & Falsos positivos con navegador & Cuantificación del bloqueo de tráfico legítimo no cubierto por Artillery \\ \hline
OG (rendimiento) & 3.1 Arquitectura en hilos; 4.5 Benchmark & Muestra M3; series de memoria y CPU & $\Delta p95$; tiempos de espera; memoria; CPU; extracción & $\Delta p95_{\mathrm{abs}}\leq$15~ms (normal y ataque; volumen como referencia); $\Delta$ memoria $\leq$150~MB (143.1~MiB); CPU $\leq$15~\% \\ \hline
OG (generalización) & 4.7 Integración en terceros & Muestra M5 & Bloqueo en ataques y en legítimos por aplicación & Evidencia de si la calibración transfiere fuera de la aplicación de referencia \\ \hline
OG (librería) & 3.1--3.2 Desarrollo y verificación; publicación & Pruebas en CI; paquete publicado & Pruebas aprobadas; paridad ONNX & Paquete funcional publicado en npm \\ \hline
\caption[Matriz de trazabilidad metodológica]{Matriz de trazabilidad entre objetivos, actividades, datos, métricas y resultados esperados}
\label{tab:met-trazabilidad}
\end{longtable}
{\footnotesize Nota. Elaboración propia. Los números de actividad remiten a los procedimientos
de la Sección~\ref{sec:met-fases}.}

% ------------------------------------------------------------------------------
\section{Cambios respecto al Marco Metodológico del protocolo}
\label{sec:met-cambios}

El Cuadro~\ref{tab:met-cambios} enumera los cambios respecto al protocolo y la causa de cada
uno.

\begin{longtable}{|L{3.4cm}|L{4.3cm}|L{5.1cm}|}
\hline
\textbf{Elemento} & \textbf{Cambio} & \textbf{Causa} \\ \hline
\endfirsthead
\multicolumn{3}{l}{\footnotesize\textit{(continuación)}} \\ \hline
\textbf{Elemento} & \textbf{Cambio} & \textbf{Causa} \\ \hline
\endhead
\multicolumn{3}{r}{\footnotesize\textit{Continúa en la página siguiente}} \\
\endfoot
\endlastfoot
Criterio de latencia & De $\Delta p95$ relativo (5~\%/10~\%) a absoluto ($\leq$15~ms) & Una cifra en milisegundos mide el costo real por petición; un porcentaje sobre una línea base de 2--5~ms es menor que la variación de medición y cambia según la aplicación anfitriona \\ \hline
Agregación de repeticiones & De media de las tres centrales a mediana & Menor sensibilidad a repeticiones atípicas \\ \hline
Entorno de medición & De Docker Desktop en macOS a Linux nativo en GitHub Actions & La capa de virtualización de Docker Desktop alteraba la latencia medida \\ \hline
Escenario de volumen & Agregado & Contrastar el efecto de una línea base realista sobre el impacto del middleware \\ \hline
Origen del dataset (OE2) & De generado en el entorno Docker a construido con fuentes públicas auditadas, listas públicas de payloads y conjuntos propios; el entorno se usa para evaluar & Las fuentes públicas aportan volumen y diversidad de ataques que el entorno controlado no podía generar en el tiempo disponible \\ \hline
Corpus de evaluación & Corpus de SecLists (Ronda~4) reemplazado por corpus limpio (Ronda~5) & Fuga confirmada en Command Injection respecto al entrenamiento \\ \hline
Política de decisión & El Isolation Forest no bloquea; solo registra anomalías & Una política disyuntiva bloqueaba $\approx$10~\% del tráfico legítimo \\ \hline
Validación externa & Startups guatemaltecas sustituidas por cuatro aplicaciones vulnerables de terceros & No se concretó un acuerdo con una startup en el período de pruebas \\ \hline
Estudio de falsos positivos con navegador & Agregado & Falso positivo observado con un navegador real, no cubierto por Artillery \\ \hline
Instrumento de CPU & clinic.js sustituido por \texttt{docker stats} & Ejecución sin intervención en CI y serie temporal alineada con la carga \\ \hline
Gestión & Kanban en Notion con cuatro columnas, sin la columna Bloqueado; complementado con protección de ramas e integración continua & Simplificar el tablero; convertir la revisión cruzada y las pruebas en controles obligatorios \\ \hline
Reentrenamiento offline (\texttt{mlops}) & Agregado, en el componente de IA & Recomendación del asesor \\ \hline
\caption[Cambios respecto al Marco Metodológico del protocolo]{Cambios respecto al Marco Metodológico del protocolo y su causa}
\label{tab:met-cambios}
\end{longtable}
{\footnotesize Nota. Elaboración propia.}

% ------------------------------------------------------------------------------
\section{Amenazas a la validez}
\label{sec:met-validez}

Las condiciones siguientes limitan lo que los datos pueden demostrar. Se declaran aquí para que
el capítulo de Resultados las tenga en cuenta al interpretar cada cifra:

\begin{itemize}
    \item \textbf{Rebloqueos y lecturas repetidas del conjunto de prueba.} El conjunto de prueba
    se volvió a bloquear en cada generación de modelo (el valor del bloqueo cambió nueve veces, entre
    el 11 de junio y el 23 de septiembre de 2026), y en la generación
    \texttt{rf\_v11} se leyó tres veces, con un cambio del umbral del Isolation Forest entre la
    primera y la última lectura. En \texttt{if\_v2} e \texttt{if\_v5}, los objetivos internos de
    calibración del Isolation Forest se fijaron después de observar resultados de prueba. Ver
    resultados de prueba antes de decidir filtra información hacia el diseño. El manifiesto de
    sumas de verificación solo protege las lecturas desde \texttt{rf\_v11}.
    \item \textbf{Procedencia de los datos restaurados.} Los datos de entrenamiento se
    restauraron desde una copia de respaldo el 23 de septiembre de 2026 sin que el repositorio
    registre su origen. La partición vigente suma 383{,}202 filas, 119 más que el archivo
    unificado (383{,}083); la diferencia coincide con las 119 filas de telemetría curada del
    pipeline de reentrenamiento, que el script de unificación excluye y el pipeline incorpora por
    separado (Cuadro~\ref{tab:met-origenes}), por lo que no queda sin explicar. El 17.8~\%
    de las filas de Command Injection del conjunto de prueba (261 de 1{,}468) proviene de
    conjuntos sintéticos cuyo script generador no está en el repositorio; esa porción del corpus no
    es reproducible con los scripts versionados.
    \item \textbf{Métricas sin umbral de bloqueo.} Las métricas del Random Forest sobre el
    conjunto de prueba asignan la clase de mayor probabilidad, mientras que el middleware bloquea
    solo si esa probabilidad alcanza 0.35. Ambas decisiones coinciden en la mayoría de los casos,
    pero no son idénticas.
    \item \textbf{Calibración en una sola aplicación.} El umbral se calibró sobre la partición
    de validación y se confirmó en la aplicación de referencia. Las aplicaciones de terceros son
    la única evidencia de cómo transfiere a otras formas de tráfico.
    \item \textbf{Corpus de ataque de las aplicaciones de terceros.} Las integraciones en DVNA,
    node-api-goat y dvws-node, realizadas por recomendación del asesor para probar la librería
    fuera de la aplicación de referencia, usaron el conjunto de 500 payloads de la suite E2E, que
    comparte payloads con el entrenamiento (Procedimiento~4.2). Su tasa de bloqueo en ataques
    confirma que la librería instalada desde npm funciona en aplicaciones ajenas al proyecto, pero
    no mide la detección de payloads nuevos; esa medición corresponde al corpus limpio (M2). Sus
    tasas de bloqueo sobre tráfico legítimo no dependen de ese corpus y son la evidencia principal
    de generalización.
    \item \textbf{Sesgo de representación de la cabecera \texttt{User-Agent}.} Solo el 0.4~\%
    de las filas legítimas del corpus tiene un \texttt{User-Agent} de más de 40 caracteres,
    frente al 33.3--97.7~\% de las clases de ataque. Las métricas sobre el conjunto de prueba
    heredan ese sesgo y no reflejan el tráfico de navegadores reales.
    \item \textbf{Generaciones de modelo distintas entre evaluaciones.} Las métricas sobre el
    conjunto de prueba corresponden al par~C, que es el contenido del paquete 0.1.1; el paquete 0.1.0, la
    medición principal de cobertura y el estudio de \texttt{User-Agent}, al par~A; las aplicaciones DVNA, node-api-goat y
    dvws-node, con atribución probable al par~A, porque el paquete y el modelo exactos de cada
    medición no constan en el repositorio; el benchmark de rendimiento y la repetición de la
    Configuración~2 del 9 de septiembre, al par~B; y la evaluación en OWASP Juice Shop, al par~C
    (Cuadro~\ref{tab:pares-modelos}). Los pares no tienen la misma entrada: el par~A consume 67 y 61
    características, y los pares B y C, 69 y 63. Comparten, en cambio, la arquitectura de ejecución
    (un hilo de Random Forest, dos de Isolation Forest, comunicación entre hilos y escritura en
    SQLite), por lo que los resultados de rendimiento describen principalmente esa arquitectura y no
    los pesos de un modelo en particular, aunque no se midieron con el paquete 0.1.1. El paquete 0.1.0
    usa el umbral del Isolation Forest calibrado para \texttt{if\_v9} (0.002486), fijado como
    constante en el código; el 0.1.1 lo lee del archivo \texttt{parity\_report.json} del modelo (0.004206).
    \item \textbf{Evaluación no ciega del par~C en la Ronda~5.} Los resultados de la Ronda~5 con el
    par~A se conocieron el 24 de septiembre de 2026 y el par~C se reconstruyó el 26 de septiembre.
    El historial no vincula ambos hechos, pero tampoco permite descartar que esos resultados
    influyeran en el reentrenamiento, por lo que la cobertura del par~C no constituye una
    evaluación estrictamente ciega; el paquete 0.1.1 contiene el mismo par.
    \item \textbf{Dos defectos de recodificación en la cadena de medición.} Hubo dos defectos
    independientes. El primero, en el middleware, que reconstruía la cadena de consulta a partir de
    \texttt{req.query}, ya decodificada una vez por Express, y le añadía una capa de codificación;
    se corrigió el 10 de octubre de 2026 y la corrección está en el paquete 0.1.1. El segundo, en el
    ejecutor de pruebas (\texttt{run\_large\_corpus.py}), que volvía a codificar los payloads de
    Path Traversal y de Command Injection que los generadores ya entregaban codificados (por
    ejemplo \texttt{\%252e\%252e\%252f} o un separador \texttt{\%0a}), de modo que el payload enviado
    no era el que el generador pretendía; se corrigió el 10 de octubre, limitado a esas dos
    categorías. Ninguna corrección bastó por sí sola (Cuadro~\ref{tab:res-pt-causas}). Todas las
    mediciones de Path Traversal y de Command Injection de la Ronda~5, y las de la Ronda~4 que usaron
    el mismo ejecutor, se hicieron con el ejecutor anterior, salvo la Configuración~2 con el paquete
    0.1.1. No se midió el paquete 0.1.0 con el ejecutor corregido, por lo que no puede separarse
    cuánto de la mejora de 85.0~\% a 100.0~\% se debe al modelo, al middleware y al ejecutor; y las
    Configuraciones 3a y 3b, que incluyen un WAF, tampoco se repitieron, así que su comparación en
    esas dos categorías es provisional.
    \item \textbf{La medición con el paquete 0.1.1 no es independiente.} Los defectos se
    encontraron analizando los fallos de este mismo corpus de la Ronda~5, y la medición que los da
    por corregidos (200 de 200 en Path Traversal) se hizo sobre ese corpus. Un resultado obtenido
    sobre los datos que motivaron la corrección es optimista por construcción; confirmarlo exige un
    corpus nuevo, generado con otra semilla y filtrado de la misma manera.
    \item \textbf{Saturación de casi-duplicados.} El corpus no se deduplicó más allá de los
    registros idénticos. La detección por MinHash, descrita en la sección VII.4 del trabajo de
    graduación de Diego Valenzuela \cite{valenzuela2026ia}, encontró saturación en las cinco clases: SQLi tiene
    aproximadamente 395 veces más pares casi duplicados que filas (dominada por la generación por
    plantillas de sqlmap y CAPEC) y la clase legítima, aproximadamente 9.5 veces. SQLi es el
    59.3~\% del conjunto de prueba. Con una partición aleatoria estratificada, variantes casi
    idénticas pueden quedar en entrenamiento y en prueba, lo que puede inflar las métricas, sobre
    todo las de SQLi; este trabajo no cuantificó ese efecto. El análisis por distancia de
    Levenshtein (211 pares, todos de la fuente CAPEC) lo subestimó, porque solo compara consultas de
    1 a 100 caracteres y hasta 10{,}000 comparaciones por clase.
    \item \textbf{Alcance de la evidencia sobre variantes desconocidas.} La detección de
    variantes no vistas de los vectores cubiertos se mide a nivel del sistema con el corpus limpio
    M2. El recall del Isolation Forest se mide sobre las mismas cuatro categorías que ve el Random
    Forest, por lo que no aísla la contribución del detector de anomalías a esa capacidad.
    \item \textbf{Validez externa del entorno simulado.} El entorno controlado no reproduce el
    tráfico de una startup en producción. La validación en aplicaciones de terceros reduce ese
    riesgo, pero no lo elimina.
\end{itemize}

% ------------------------------------------------------------------------------
\section{Restricciones y condiciones de ejecución}
\label{sec:met-restricciones}

El proyecto se ejecutó bajo las restricciones de tiempo (mayo a octubre de 2026), presupuesto y
recursos humanos descritas en el capítulo de Alcance. Además, las siguientes condiciones afectan
el alcance de la evidencia disponible para el capítulo de Resultados y se ejecutarán solo como
trabajo futuro:

\begin{itemize}
    \item \textbf{Cobertura de XSS con n=1.} ZAP confirmó una única vulnerabilidad de XSS en la
    aplicación de referencia (\texttt{POST /profile}, campo \texttt{bio}); la cobertura de esa
    categoría descansa en ese caso y no tiene el peso estadístico de las otras tres (n=200 cada
    una). Que ZAP reporte un solo caso refleja la aplicación evaluada, no un límite de la
    herramienta.
    \item \textbf{Benchmark no repetido con el modelo sincronizado.} Ver la amenaza a la validez
    correspondiente en la Sección~\ref{sec:met-validez}.
    \item \textbf{Sin un escenario sostenido de condición única.} Cada repetición del benchmark
    dura 75~s. Se ejecutó además una corrida de monitoreo continuo de 41.9~minutos (374 muestras,
    con un pico de 494.8~MiB en el contenedor), pero esa serie mezcla condiciones distintas, y el
    indicador heurístico \texttt{possibleLeak}, que compara las dos mitades de la serie, no es
    concluyente por esa razón. No se ejecutó un escenario de 30~minutos o más con una sola condición
    estable, que es el que permitiría descartar una fuga de memoria.
    \item \textbf{Validación con startups no realizada.} No se concretó un acuerdo de
    colaboración en el período de pruebas; la validación en entornos productivos queda fuera de
    la evidencia del proyecto.
\end{itemize}

El Cuadro~\ref{tab:met-riesgos} actualiza los riesgos metodológicos del protocolo con su
estado al cierre de la ejecución.

\begin{longtable}{|L{3.6cm}|L{2.3cm}|L{6.9cm}|}
\hline
\textbf{Riesgo del protocolo} & \textbf{Estado} & \textbf{Mitigación aplicada} \\ \hline
\endfirsthead
\multicolumn{3}{l}{\footnotesize\textit{(continuación)}} \\ \hline
\textbf{Riesgo del protocolo} & \textbf{Estado} & \textbf{Mitigación aplicada} \\ \hline
\endhead
\multicolumn{3}{r}{\footnotesize\textit{Continúa en la página siguiente}} \\
\endfoot
\endlastfoot
Ningún modelo alcanza F1 $\geq$0.80 & No materializado & Datos reales adicionales en lugar de muestras sintéticas para Command Injection; lectura de prueba con verificación de hash \\ \hline
Fuga de datos entre particiones & Materializado en el corpus de evaluación, no en las particiones & Auditoría exacta y por plantilla; corpus limpio con filtro de exclusión \\ \hline
Diferencia de paridad ONNX $>$0.1~\% & No materializado & Prueba de paridad obligatoria en CI \\ \hline
Falsos positivos ante tráfico real & Materializado & El Isolation Forest no bloquea; estudio con navegadores; modo \texttt{monitor} documentado para una primera instalación \\ \hline
Imposibilidad de acuerdo con startups & Materializado & Sustitución por cuatro aplicaciones de terceros \\ \hline
Extracción de características $>$5~ms & No materializado & Extractor único en TypeScript con benchmark propio \\ \hline
\caption[Estado de los riesgos metodológicos]{Estado de los riesgos metodológicos del protocolo al cierre de la ejecución}
\label{tab:met-riesgos}
\end{longtable}
{\footnotesize Nota. Elaboración propia.}

La viabilidad del método descansa en que todos los procedimientos usan herramientas gratuitas
y scripts versionados, se ejecutan en hardware local o en ejecutores públicos de integración
continua y dejan sus datos crudos en los repositorios del proyecto. Esto permite que otro
investigador los repita sin depender del entorno de los autores.

\ifdefined\modoensamble
\else
	\section*{Referencias}
	\printbibliography[heading=none]
	\end{document}
\fi

\fi

\ifdefined\CAPresultados
	\newpage
	% ==============================================================================
% RESULTADOS
% ==============================================================================
% Organizado por objetivo (general, OE1, OE2, OE3), con un cuadro SMART al
% inicio de cada uno. Cada cifra proviene de un artefacto del repositorio;
% la fuente se indica en la nota de cada cuadro o figura.
%
% Fuentes principales (repositorio logSguarDian salvo indicación):
%   - docs/cybersecurity-objectives-compliance.md (estado al 2026-10-02).
%   - training/results/v11_test_results.json (lectura de prueba 2026-09-26).
%   - training/models/parity_report.json (vigente, rf_v11/if_v10).
%   - docs/results.md (A19, A24, F4.4).
%   - benchmarks/results/extractor-2026-10-09T23-34-*.json (extractor, PR #97;
%     reemplaza la cifra F1.8 de docs/results.md, que no tenía artefacto).
%   - docs/vulnerable-projects-tests/ y docs/limitations.md (terceros).
%   - logSguarDian-vulnerable-project: docs/config3b-results.md,
%     attack-sim/ua_fp_study_results.json, exclusion_report.json.
%
% Imágenes: copiadas de docs/cibersecurity-images/ y training/results/ a
% figuras/resultados/. Las capturas que muestran una generación de modelo o
% un estado anterior se rotulan como tales en su leyenda; ninguna se
% presenta como cifra vigente si no lo es.
%
% Se omitieron obj3-b (reporte de validación de una generación anterior,
% redundante con obj3-a y con el cuadro vigente) y obj3-f (misma Ronda 4
% que obj3-g, solo Configuración 2).
\ifdefined\modoensamble
\else
	\input{0-preambulo_standalone}
\fi

\chapter{Resultados}
\label{cap:resultados}

Este capítulo presenta la evidencia obtenida mediante los procedimientos del
Capítulo~\ref{cap:metodologia}, organizada por objetivo: primero el objetivo general y después
los tres objetivos específicos. Cada sección inicia con un cuadro que descompone el objetivo según
los criterios SMART \cite{doran1981smart} (específico, medible, alcanzable, relevante y temporal) y asocia cada criterio
con la evidencia correspondiente. A continuación se presentan los datos, mediciones y capturas que
sustentan esa evidencia. Cuando un indicador tiene un criterio de aceptación, se indica si el valor
medido lo cumple; la interpretación de los resultados se desarrolla en el capítulo de Análisis de resultados.

Cada resultado indica el par de modelos ONNX con que se obtuvo, según el
Cuadro~\ref{tab:pares-modelos}: las métricas sobre el conjunto de prueba y la evaluación en OWASP
Juice Shop corresponden al par~C, que es el contenido del paquete 0.1.1; el paquete 0.1.0, la
medición principal de cobertura sobre el corpus limpio y el estudio de \texttt{User-Agent}, al
par~A; las aplicaciones DVNA, node-api-goat y
dvws-node, con atribución probable al par~A, pues el modelo exacto de cada medición no consta en el
repositorio; y el benchmark de rendimiento y la repetición de la Configuración~2 del 9 de
septiembre, al par~B. Algunas
capturas reflejan generaciones anteriores del modelo o del proyecto y se incluyen como valores
iniciales para la comparación; su leyenda lo indica de forma explícita.

% ==============================================================================
\section{Objetivo general}
\label{sec:res-og}

El objetivo general establece: diseñar e implementar una librería npm que integre un pipeline de
ingeniería de características y un modelo híbrido de aprendizaje automático capaz de detectar y
clasificar amenazas conocidas y anomalías estadísticas en endpoints Node.js en tiempo de
ejecución, sin comprometer el rendimiento del servidor anfitrión. El Cuadro~\ref{tab:res-smart-og}
lo descompone según los criterios SMART.

\begin{table}[H]
\centering
\small
\begin{tabular}{|L{2.2cm}|L{5.2cm}|L{5.6cm}|}
\hline
\textbf{Criterio} & \textbf{Definición en el objetivo} & \textbf{Evidencia obtenida} \\ \hline
Específico & Librería npm para Express con pipeline de características y modelo híbrido (Random
Forest + Isolation Forest) & Paquete \texttt{logsguardian} publicado (0.1.0 y 0.1.1); extractor de 76
características; RF con autoridad de bloqueo e IF de registro de anomalías
(Sección~\ref{subsec:res-og-libreria}) \\ \hline
Medible & Paridad ONNX $<$0.1~\%; extracción $\leq$5~ms; $\Delta p95_{\mathrm{abs}}\leq$15~ms;
$\Delta$ memoria $\leq$150~MB (143.1~MiB); CPU $\leq$15~\% & Valores en los Cuadros~\ref{tab:res-paridad},
\ref{tab:res-extractor}, \ref{tab:res-latencia} y \ref{tab:res-memoria} \\ \hline
Alcanzable & Stack de código abierto, sin costo de licencias & Construido con Node.js,
\texttt{onnxruntime-node}, SQLite, Docker y GitHub Actions, sin licencias de pago \\ \hline
Relevante & Protección a nivel de aplicación para organizaciones sin equipo de seguridad &
Instalación con un solo comando (\texttt{npm i logsguardian}); validación en cuatro aplicaciones de
terceros (Sección~\ref{subsec:res-og-terceros}) \\ \hline
Temporal & Mayo a octubre de 2026 & Publicación en npm el 30 de agosto de 2026 (0.1.0) y el 11 de octubre de 2026 (0.1.1); mediciones finales entre el 23 de septiembre y el 11 de octubre de 2026 \\ \hline
\end{tabular}
\caption[Criterios SMART del objetivo general]{Criterios SMART del objetivo general y evidencia asociada}
\label{tab:res-smart-og}
\end{table}
{\footnotesize Nota. Elaboración propia.}

\subsection{Librería publicada y verificación funcional}
\label{subsec:res-og-libreria}

La Figura~\ref{fig:res-npm} muestra el paquete publicado en el registro público de npm, en su
versión 0.1.0, con licencia MIT, publicada el 30 de agosto de 2026. El 11 de octubre de 2026 se
publicó la versión 0.1.1, hoy la vigente, que incorpora los modelos \texttt{rf\_v11} e
\texttt{if\_v10} y la corrección del defecto de recodificación de la consulta
(Cuadro~\ref{tab:pares-modelos}). La publicación se realiza de forma automática desde GitHub
Actions al fusionar la rama \texttt{develop} en \texttt{main}, siempre que la versión local difiera
de la publicada.

\begin{figure}[H]
\centering
\includegraphics[width=0.92\textwidth]{resultados/obj0-a.png}
\caption[Paquete logsguardian publicado en npm]{Página del paquete \texttt{logsguardian} en el registro público de npm (versión 0.1.0, captura previa a la publicación de la 0.1.1)}
\label{fig:res-npm}
\end{figure}
{\footnotesize Nota. Captura de npmjs.com/package/logsguardian. Fecha de publicación según el registro de npm.}

La Figura~\ref{fig:res-tests} muestra la estructura del monorepo (paquetes \texttt{core},
\texttt{extractor} y \texttt{mlops}) y la ejecución de la suite de pruebas del extractor de
características: 11 suites y 74 pruebas aprobadas, entre ellas las de determinismo,
decodificación recursiva y Command Injection compuesto.

\begin{figure}[H]
\centering
\includegraphics[width=0.92\textwidth]{resultados/obj0-b.png}
\caption[Paquetes del monorepo y pruebas del extractor]{Paquetes del monorepo y resultado de la suite de pruebas del extractor de características}
\label{fig:res-tests}
\end{figure}
{\footnotesize Nota. Captura de terminal, rama \texttt{feat/goals\_ciber}, repositorio \texttt{logSguarDian}.}

La Figura~\ref{fig:res-cli} muestra la CLI de la librería ejecutada dentro del contenedor de la
aplicación de referencia: \texttt{attacks list} cataloga los eventos clasificados por tipo (294
SQLi, 181 Path Traversal, 90 XSS y 23 Command Injection en la base de eventos de esa sesión), y
\texttt{attacks inspect sqli} muestra las métricas del modelo, las características de mayor
importancia y ejemplos de payloads detectados. En esa captura, de \texttt{rf\_v3}, la característica
de mayor importancia global del Random Forest es la longitud de la cabecera \texttt{User-Agent}
(\texttt{ua\_length}, 0.0796). En el análisis SHAP \cite{lundberg2017unified,valenzuela2026ia} de \texttt{rf\_v11}, \texttt{ua\_length} es también
la característica de mayor contribución media en las clases legítima (0.0866) y SQLi (0.0734).

\begin{figure}[H]
\centering
\includegraphics[width=\textwidth]{resultados/obj3-d.png}
\caption[CLI de logsguardian: attacks list e inspect]{Salida de los comandos \texttt{attacks list} y \texttt{attacks inspect sqli} de la CLI de \texttt{logsguardian}}
\label{fig:res-cli}
\end{figure}
{\footnotesize Nota. Captura de terminal, 22 de septiembre de 2026. Las métricas que muestra
\texttt{attacks inspect} en esta captura corresponden a \texttt{rf\_v3}; el archivo de métricas se
actualizó después a \texttt{rf\_v11}, y la versión 0.1.1 del paquete ya lo incluye
(Cuadro~\ref{tab:res-rf}).}

La Figura~\ref{fig:res-e2e} muestra la suite de pruebas extremo a extremo, que envía por HTTP real
100 payloads por clase a una aplicación Express con el middleware y los modelos ONNX publicados.
Las ocho pruebas se aprobaron: 100~\% de bloqueo en SQLi, 98~\% en XSS, 99~\% en Path Traversal,
96~\% en Command Injection y 2~\% de falsos positivos sobre peticiones legítimas. Como se indicó en
la Sección~\ref{subsec:met-fuga}, este conjunto comparte payloads con el entrenamiento; por ello
se presenta como verificación del funcionamiento extremo a extremo de la librería y no como
medición de detección sobre payloads nuevos, que se reporta en la Sección~\ref{subsec:res-oe3-cobertura}.

\begin{figure}[H]
\centering
\includegraphics[width=0.62\textwidth]{resultados/obj3-e.png}
\caption[Suite de pruebas extremo a extremo]{Resultado de la suite de detección extremo a extremo (\texttt{pnpm run test:e2e})}
\label{fig:res-e2e}
\end{figure}
{\footnotesize Nota. Captura de terminal, repositorio \texttt{logSguarDian}.}

\subsection{Paridad de la serialización ONNX}

El Cuadro~\ref{tab:res-paridad} presenta la paridad entre las predicciones de los modelos en
Python y sus versiones ONNX, calculada sobre 1{,}000 muestras para los modelos vigentes. Las
diferencias máximas son del orden de $10^{-7}$, por debajo del criterio de 0.1~\%. La
Figura~\ref{fig:res-paridad} muestra el reporte de paridad y la prueba automatizada que compara
\texttt{onnxruntime-node} con Python, con diferencia máxima de 0 en ambos modelos en esa captura, previa al reentrenamiento.

\begin{table}[H]
\centering
\small
\begin{tabular}{|L{4.6cm}|L{3.2cm}|L{3.2cm}|}
\hline
\textbf{Indicador} & \textbf{Random Forest} & \textbf{Isolation Forest} \\ \hline
Versión del modelo & \texttt{rf\_v11} & \texttt{if\_v10} \\ \hline
Características de entrada & 69 & 63 \\ \hline
Diferencia máxima Python--ONNX & $1.01\times10^{-7}$ & $2.42\times10^{-7}$ \\ \hline
Muestras comparadas & 1{,}000 & 1{,}000 \\ \hline
Criterio ($<$0.1~\%) & Cumple & Cumple \\ \hline
\end{tabular}
\caption[Paridad entre modelos Python y ONNX]{Paridad entre las predicciones de los modelos en Python y en ONNX}
\label{tab:res-paridad}
\end{table}
{\footnotesize Nota. Elaboración propia con datos de \texttt{training/models/parity\_report.json}.}

\begin{figure}[H]
\centering
\includegraphics[width=0.72\textwidth]{resultados/obj3-k.png}
\caption[Reporte y prueba de paridad ONNX]{Reporte de paridad ONNX y prueba automatizada de paridad entre Node.js y Python}
\label{fig:res-paridad}
\end{figure}
{\footnotesize Nota. Captura de terminal, repositorio \texttt{logSguarDian}. La captura es previa al
reentrenamiento del 26 de septiembre de 2026; los valores vigentes son los del
Cuadro~\ref{tab:res-paridad}.}

\subsection{Rendimiento: extracción de características y latencia}

El Cuadro~\ref{tab:res-extractor} presenta el tiempo de extracción del vector de características
por petición en dos series: la extracción directa en el hilo principal y la ida y vuelta completa
a un \texttt{worker\_thread} que ejecuta el extractor, que corresponde a la forma en que el
Alcance define el criterio. Cada serie se midió de forma serial sobre 2{,}000 iteraciones por tipo
de petición, en tres corridas independientes. En tráfico mixto, el p95 fue de 0.0126~ms en la
serie directa y de 0.0273~ms en la ida y vuelta al hilo; el mayor p95 de cualquier tipo de petición
en la ida y vuelta fue de 0.0305~ms (XSS). Todos los valores están por debajo del criterio de 5~ms.

\begin{table}[H]
\centering
\small
\begin{tabular}{|L{2.9cm}|L{1.5cm}|L{1.5cm}|L{1.5cm}|L{1.5cm}|L{1.5cm}|L{1.5cm}|}
\hline
 & \multicolumn{3}{L{5.34cm}|}{\textbf{Directa, hilo principal (ms)}} & \multicolumn{3}{L{5.34cm}|}{\textbf{Ida y vuelta a \texttt{worker\_thread} (ms)}} \\ \hline
\textbf{Tipo de petición} & \textbf{p50} & \textbf{p95} & \textbf{p99} & \textbf{p50} & \textbf{p95} & \textbf{p99} \\ \hline
Legítima & 0.0103 & 0.0169 & 0.0409 & 0.0213 & 0.0304 & 0.0467 \\ \hline
SQLi & 0.0092 & 0.0120 & 0.0340 & 0.0210 & 0.0275 & 0.0405 \\ \hline
XSS & 0.0102 & 0.0135 & 0.0340 & 0.0226 & 0.0305 & 0.0430 \\ \hline
Path Traversal & 0.0076 & 0.0097 & 0.0302 & 0.0162 & 0.0264 & 0.0377 \\ \hline
Command Injection & 0.0090 & 0.0102 & 0.0314 & 0.0212 & 0.0283 & 0.0405 \\ \hline
\textbf{Mixto} & \textbf{0.0097} & \textbf{0.0126} & \textbf{0.0350} & \textbf{0.0206} & \textbf{0.0273} & \textbf{0.0413} \\ \hline
\end{tabular}
\caption[Tiempo de extracción de características]{Tiempo de extracción del vector de características por petición, directa y con ida y vuelta a un \texttt{worker\_thread}}
\label{tab:res-extractor}
\end{table}
{\footnotesize Nota. Elaboración propia con datos de \texttt{benchmarks/results/extractor-2026-10-09T23-34-*.json}
(repositorio \texttt{logSguarDian}, commit \texttt{7b036bb}). Mediana de tres corridas; Apple M5,
Node.js v25.6.0; 200 iteraciones de calentamiento descartadas por serie.}

El Cuadro~\ref{tab:res-latencia} presenta la latencia del servidor anfitrión medida en Linux
nativo (GitHub Actions, \texttt{ubuntu-latest}, 26 de septiembre de 2026), con cinco repeticiones
por condición y la mediana como valor reportado. En los escenarios normal y de ataque, el
incremento absoluto del p95 fue de 12.08~ms y 12.85~ms respectivamente, ambos por debajo del
criterio de 15~ms. En el escenario de volumen, el criterio de 15~ms no se cumple: el p95 con middleware fue de 6{,}559.3~ms y el
7.00~\% de las peticiones (1{,}824 de 26{,}048) terminó por tiempo de espera excedido; en las
demás combinaciones de escenario y condición no hubo tiempos de espera excedidos. Los tiempos de
espera se concentraron en los últimos intervalos de cada repetición.

\begin{table}[H]
\centering
\small
\begin{tabular}{|L{2.3cm}|L{1.6cm}|L{1.6cm}|L{1.7cm}|L{1.7cm}|L{2.6cm}|}
\hline
\textbf{Escenario} & \textbf{p95 base (ms)} & \textbf{p95 activo (ms)} & \textbf{$\Delta p95_{\mathrm{abs}}$ (ms)} & \textbf{$\Delta p95_{\mathrm{rel}}$} & \textbf{Tiempos de espera excedidos} \\ \hline
Normal & 2.49 & 14.57 & 12.08 & +485.1~\% & 0 \\ \hline
Ataque & 4.62 & 17.47 & 12.85 & +278.1~\% & 0 \\ \hline
Volumen & 63.84 & 6{,}559.30 & 6{,}495.46 & +10{,}174.6~\% & 1{,}824 / 26{,}048 (7.00~\%) \\ \hline
\multicolumn{3}{|L{6.3cm}|}{Criterio ($\Delta p95_{\mathrm{abs}}\leq$15~ms)} & \multicolumn{3}{L{6.8cm}|}{Cumple en normal y en ataque; no cumple en volumen} \\ \hline
\end{tabular}
\caption[Latencia p95 con y sin middleware]{Latencia p95 del servidor anfitrión con y sin middleware, Linux nativo}
\label{tab:res-latencia}
\end{table}
{\footnotesize Nota. Elaboración propia con datos del flujo \texttt{latency-benchmark.yml}
(repositorio \texttt{logSguarDian-vulnerable-project}). Mediana de cinco repeticiones; el p95 de
cada repetición es el promedio ponderado de los p95 por intervalo posteriores al calentamiento.
Corrida 36276690236 de GitHub Actions (2 vCPU, 7.8~GiB), par de modelos B; evidencia cruda en
\texttt{docs/evidence/latency-linux-2026-09-26.tgz}.}

Como valores iniciales, la Figura~\ref{fig:res-lat-hist} (paneles~\ref{fig:res-lat-hist-rel} y~\ref{fig:res-lat-hist-abs})
presenta las mediciones previas en Docker Desktop sobre macOS, para cuatro variantes de la
arquitectura de escritura y de espera del Isolation Forest. La variante publicada (híbrida)
registró un incremento absoluto de 7 a 9~ms y relativo de +142~\% a +178~\%.

\begin{figure}[H]
\centering
\begin{subfigure}[t]{0.52\textwidth}
\centering
\includegraphics[width=\textwidth]{resultados/obj3-h1.png}
\caption{Incremento relativo}
\label{fig:res-lat-hist-rel}
\end{subfigure}
\hfill
\begin{subfigure}[t]{0.46\textwidth}
\centering
\includegraphics[width=\textwidth]{resultados/obj3-h2.png}
\caption{Incremento absoluto}
\label{fig:res-lat-hist-abs}
\end{subfigure}
\caption[Latencia inicial en Docker Desktop por variante de arquitectura]{Incremento de latencia p95 por variante de arquitectura, mediciones iniciales en Docker Desktop sobre macOS}
\label{fig:res-lat-hist}
\end{figure}
{\footnotesize Nota. Capturas de \texttt{docs/vulnerable-app-evaluation/config2-latency-evaluation.md}.
Mediciones previas a la metodología en Linux nativo del Cuadro~\ref{tab:res-latencia}.}

\subsection{Memoria y CPU}

El Cuadro~\ref{tab:res-memoria} presenta la memoria y el uso de CPU del contenedor de la
aplicación en la corrida de Linux nativo, con las muestras asignadas a cada repetición y el periodo
de calentamiento excluido. La memoria se expresa en MiB, la unidad que reporta \texttt{docker
stats}. Sin middleware, el contenedor usó en promedio entre 63 y 83~MiB; con el middleware activo,
entre 395 y 459~MiB. El incremento pareado de memoria fue de 332.02~MiB bajo carga normal,
331.86~MiB bajo carga de ataque y 375.92~MiB en el escenario de volumen. El criterio del protocolo
es un incremento de 150~MB (143.1~MiB) como máximo, y el valor medido lo supera más de dos veces en
los tres escenarios. El incremento de CPU fue de 88.93 puntos porcentuales bajo carga normal
(+614~\%) y de 85.12 bajo carga de ataque (+530~\%), por encima del criterio de 15~\%. El incremento
de memoria fue prácticamente igual en los escenarios normal y de ataque.

\begin{table}[H]
\centering
\small
\begin{tabular}{|L{1.8cm}|L{1.8cm}|L{1.8cm}|L{2.0cm}|L{2.0cm}|L{2.3cm}|}
\hline
\textbf{Escenario} & \textbf{Memoria base (MiB)} & \textbf{Memoria activo (MiB)} & \textbf{$\Delta$ memoria (MiB)} & \textbf{CPU base / activo (\%)} & \textbf{$\Delta$ CPU (pp; relativo)} \\ \hline
Normal & 65.18 & 397.51 & +332.02 & 14.48 / 103.32 & +88.93 (+614~\%) \\ \hline
Ataque & 63.22 & 394.78 & +331.86 & 16.07 / 101.25 & +85.12 (+530~\%) \\ \hline
Volumen & 83.29 & 459.15 & +375.92 & 82.34 / 115.29 & +32.99 (+40~\%) \\ \hline
\multicolumn{3}{|L{6.24cm}|}{Criterio: $\Delta$ memoria $\leq$150~MB (143.1~MiB); $\Delta$ CPU $\leq$15~\%} & \multicolumn{3}{L{7.14cm}|}{No cumple en ninguno de los escenarios} \\ \hline
\end{tabular}
\caption[Memoria y CPU con y sin middleware]{Memoria y CPU del contenedor de la aplicación con y sin middleware, Linux nativo}
\label{tab:res-memoria}
\end{table}
{\footnotesize Nota. Elaboración propia con datos de \texttt{docs/evidence/resources\_summary.md}
(repositorio \texttt{logSguarDian-vulnerable-project}), ventana A: periodo medido de cada
repetición, primeros 15~s excluidos; media de cinco repeticiones por condición e incremento como
media de las diferencias pareadas. Corrida 36276690236, par de modelos B, 2 vCPU. La memoria es
la del contenedor completo, no solo la del proceso Node.js. Los valores de CPU superiores a 100~\%
corresponden a más de un núcleo.}

\subsection{Validación en aplicaciones de terceros}
\label{subsec:res-og-terceros}

El Cuadro~\ref{tab:res-terceros} presenta los resultados de integrar el paquete publicado en npm, entonces la versión 0.1.0 (atribución
probable al par~A; el modelo exacto de cada medición no consta en el repositorio), en modo \texttt{block} y con la configuración por defecto, en tres aplicaciones vulnerables de
terceros. En cada una se confirmó primero la explotabilidad sin protección. Las tasas de bloqueo en
ataques se obtuvieron con el conjunto de payloads de la suite extremo a extremo
(Sección~\ref{sec:met-validez}), y las tasas sobre tráfico legítimo con peticiones de cada
aplicación. En modo \texttt{block}, la librería bloqueó entre el 92.5~\% y el 100~\% de los ataques y una proporción igual
de alta del tráfico legítimo con \texttt{User-Agent} de navegador, aunque las muestras de tráfico legítimo difieren en tamaño y
diversidad: en node-api-goat, las 50 peticiones se dirigieron a una sola ruta
(\texttt{/hexToRgb}), y la medición de DVNA, con 5 peticiones, es preliminar.

\begin{table}[H]
\centering
\small
\begin{tabular}{|L{2.0cm}|L{2.4cm}|L{1.3cm}|L{1.3cm}|L{1.4cm}|L{1.3cm}|L{1.8cm}|}
\hline
\textbf{Aplicación} & \textbf{Arquitectura} & \textbf{SQLi} & \textbf{XSS} & \textbf{Path Trav.} & \textbf{CMDi} & \textbf{Legítimo bloqueado} \\ \hline
DVNA \cite{appsecco2026dvna} & Express, EJS, MySQL & 38/40 & 39/40 & No aplica & 37/40 & 5/5 (100~\%)\textsuperscript{a} \\ \hline
node-api-goat \cite{layro012026nodeapigoat} & API pura sin base de datos & No aplica & 49/50 & 50/50 & 48/50 & 50/50 (100~\%)\textsuperscript{b} \\ \hline
dvws-node \cite{snoopysecurity2026dvwsnode} & Express, MongoDB/MySQL, SOAP/GraphQL & 40/40 & 37/40 & 40/40 & 40/40 & 37/40 (92.5~\%) \\ \hline
\end{tabular}
\caption[Bloqueo en aplicaciones de terceros]{Ataques y peticiones legítimas bloqueados en aplicaciones vulnerables de terceros}
\label{tab:res-terceros}
\end{table}
{\footnotesize Nota. Elaboración propia con datos de \texttt{docs/vulnerable-projects-tests/}. "No
aplica": la aplicación no expone un endpoint vulnerable a ese vector (verificado en el código).
En node-api-goat, sin base de datos, la categoría SQLi no se evaluó como vulnerabilidad.
\textsuperscript{a}Medición preliminar, $n=5$. \textsuperscript{b}50 valores distintos enviados a una
sola ruta (\texttt{/hexToRgb}), todos con el \texttt{User-Agent} de Chrome 128; una comprobación con
cinco peticiones y el mismo \texttt{User-Agent} bloqueó 4 de 5. Commits base evaluados, clonados el 30 de septiembre de 2026 y modificados con los cambios de integración y de compatibilidad registrados en \texttt{docs/vulnerable-projects-tests/}, sin commitear: DVNA \cite{appsecco2026dvna}, commit \texttt{9ba473a} (8 de noviembre de
2023); node-api-goat \cite{layro012026nodeapigoat}, commit \texttt{868ee6c} (12 de agosto de 2024);
dvws-node \cite{snoopysecurity2026dvwsnode}, commit \texttt{eb6e288} (7 de septiembre de 2026).}

Dos resultados complementan el cuadro. En dvws-node, la tasa de bloqueo de tráfico legítimo fue
de 75~\% (30/40) con el \texttt{User-Agent} por defecto de Node.js y de 92.5~\% (37/40) con un
\texttt{User-Agent} de navegador. En OWASP Juice Shop, evaluada por el componente de IA en modo
\texttt{monitor} con una copia local del paquete que contiene el par~C (\texttt{25b407e6...} y
\texttt{5be28ee8...}), y no con el paquete del registro, el modelo habría bloqueado 303 de 312
peticiones legítimas con \texttt{User-Agent} de navegador (97.1~\%) y 123 de 312 con cURL
(39.4~\%). El Cuadro~\ref{tab:res-terceros-ua} reúne, para cada medición de bloqueo de tráfico legítimo, el
modelo, el modo y el \texttt{User-Agent}, porque las cifras no son comparables entre sí si se omite
alguna de esas tres condiciones. En modo \texttt{monitor} la librería no bloquea: la cifra es la
proporción de peticiones a las que el modelo asignó veredicto de bloqueo.

\begin{table}[H]
\centering
\small
\begin{tabular}{|L{2.6cm}|L{2.2cm}|L{1.8cm}|L{3.0cm}|L{3.0cm}|}
\hline
\textbf{Aplicación} & \textbf{Modelo} & \textbf{Modo} & \textbf{User-Agent} & \textbf{Legítimo con veredicto de bloqueo} \\ \hline
Juice Shop & Par~C & \texttt{monitor} & Navegador / cURL & 303/312 (97.1~\%) / 123/312 (39.4~\%) \\ \hline
dvws-node & No consta & \texttt{block} & Navegador / Node.js por defecto & 37/40 (92.5~\%) / 30/40 (75~\%) \\ \hline
node-api-goat & No consta & \texttt{block} & Chrome 128 & 50/50 en una ruta; 4/5 en una comprobación \\ \hline
DVNA & No consta & \texttt{block} & No consta & 5/5, preliminar \\ \hline
Aplicación de referencia (Artillery) & Par~B & \texttt{block} & Artillery & 0 bloqueos (0 de 6{,}300 por repetición) \\ \hline
Aplicación de referencia (navegadores) & Par~A & \texttt{block} & 15 combinaciones de navegadores reales & 13/15 \\ \hline
\end{tabular}
\caption[Bloqueo de tráfico legítimo por condición]{Bloqueo de tráfico legítimo según aplicación, modelo, modo y \texttt{User-Agent}}
\label{tab:res-terceros-ua}
\end{table}
{\footnotesize Nota. Elaboración propia con datos de \texttt{docs/vulnerable-projects-tests/} y de las corridas de latencia y de User-Agent del proyecto. ``No consta'': el modelo exacto de la medición no se registró; la atribución probable de las tres aplicaciones de terceros es el par~A.}

En las tres aplicaciones de la tabla, la clase predicha de los ataques bloqueados
difirió con frecuencia de la clase real; por ejemplo, en DVNA, 27 de los 37 Command Injection
bloqueados se clasificaron como SQLi. Una quinta aplicación, OWASP NodeGoat, no se integró porque
su imagen base (Node.js 12) es incompatible con los requisitos de la librería (Node.js $\geq$20).

% ==============================================================================
\section{Objetivo específico 1: evaluación de los vectores de ataque}
\label{sec:res-oe1}

\begin{table}[H]
\centering
\small
\begin{tabular}{|L{2.2cm}|L{5.2cm}|L{5.6cm}|}
\hline
\textbf{Criterio} & \textbf{Definición en el objetivo} & \textbf{Evidencia obtenida} \\ \hline
Específico & Cuatro vectores intra-request: SQLi, XSS, Path Traversal/LFI y Command Injection &
Cinco clases canónicas (cuatro de ataque y una legítima) en la taxonomía del dataset \\ \hline
Medible & Mapeo de cada vector a OWASP Top 10 2021 y MITRE ATT\&CK & 4/4 vectores mapeados a
categoría OWASP, CWE y técnica ATT\&CK, verificados en fuente oficial \\ \hline
Alcanzable & Marcos públicos y gratuitos & Páginas oficiales de OWASP y MITRE \\ \hline
Relevante & Estructurar el dataset del OE2 y los criterios del OE3 & Las cinco clases definen la
etiqueta de M1 y las categorías de todas las métricas del OE3 \\ \hline
Temporal & Fase 1 (mayo de 2026) & Taxonomía versionada en junio de 2026 \\ \hline
\end{tabular}
\caption[Criterios SMART del objetivo específico 1]{Criterios SMART del objetivo específico 1 y evidencia asociada}
\label{tab:res-smart-oe1}
\end{table}
{\footnotesize Nota. Elaboración propia.}

El Cuadro~\ref{tab:res-smart-oe1} resume el objetivo. El mapeo de los cuatro vectores a CWE, OWASP
Top 10 2021 y MITRE ATT\&CK se presenta en el Cuadro~\ref{cuadro:owasp-mitre-cwe} del Marco
Teórico: SQL Injection, XSS y Command Injection corresponden a A03:2021 (Injection) y Path
Traversal/LFI a A01:2021 (Broken Access Control). El Cuadro~\ref{tab:res-oe1} presenta cómo cada
vector se materializó en el proyecto: su clase canónica, su origen en la fuente CAPEC del corpus, el
endpoint de la aplicación de referencia donde se evaluó y la herramienta que generó sus payloads de
evaluación.

\begin{table}[H]
\centering
\small
\begin{tabular}{|L{2.3cm}|L{2.0cm}|L{2.3cm}|L{3.4cm}|L{2.6cm}|}
\hline
\textbf{Vector} & \textbf{Clase canónica} & \textbf{Columna CAPEC} & \textbf{Endpoint de evaluación} & \textbf{Generación de payloads} \\ \hline
SQL Injection & \texttt{sqli} & 66 & \texttt{GET /posts/search?q=}, \texttt{POST /login} & sqlmap \\ \hline
XSS & \texttt{xss} & 242 (Code Injection) & \texttt{POST /posts}, \texttt{POST /profile} & Payloads clásicos (ZAP no automatizado) \\ \hline
Path Traversal/LFI & \texttt{path\_traversal} & 126 & \texttt{GET /posts/:id/ attachment?file=} & Generador propio \\ \hline
Command Injection & \texttt{cmdi} & 88, 248 & \texttt{POST /admin/ping} & Generador propio \\ \hline
\end{tabular}
\caption[Materialización de los vectores en el proyecto]{Materialización de los cuatro vectores en la taxonomía del dataset y en el entorno de evaluación}
\label{tab:res-oe1}
\end{table}
{\footnotesize Nota. Elaboración propia con datos de \texttt{training/label\_map.yaml} y del
repositorio \texttt{logSguarDian-vulnerable-project}. La asignación de CAPEC-242 a \texttt{xss} es
una decisión de mapeo del dataset, no una equivalencia conceptual.}

La Figura~\ref{fig:res-labelmap} muestra la taxonomía de cinco clases y el mapeo de las columnas
CAPEC en el archivo versionado de etiquetas. La Figura~\ref{fig:res-featurespec} muestra un
fragmento de la especificación del vector de características, que define cada característica, su
fórmula y la categoría de ataque que discrimina; el vector completo tiene 76 características, de
las cuales el Random Forest consume 69 y el Isolation Forest 63.

\begin{figure}[H]
\centering
\includegraphics[width=0.68\textwidth]{resultados/obj1-a.png}
\caption[Taxonomía de clases en label\_map.yaml]{Taxonomía canónica de cinco clases y mapeo de columnas CAPEC en \texttt{label\_map.yaml}}
\label{fig:res-labelmap}
\end{figure}
{\footnotesize Nota. Captura de \texttt{training/label\_map.yaml}, repositorio \texttt{logSguarDian}.}

\begin{figure}[H]
\centering
\includegraphics[width=\textwidth]{resultados/obj1-b.png}
\caption[Fragmento de la especificación de características]{Fragmento de la especificación del vector de características}
\label{fig:res-featurespec}
\end{figure}
{\footnotesize Nota. Captura de \texttt{docs/feature-spec.md}, repositorio \texttt{logSguarDian}.}

% ==============================================================================
\section{Objetivo específico 2: dataset etiquetado}
\label{sec:res-oe2}

\begin{table}[H]
\centering
\small
\begin{tabular}{|L{2.2cm}|L{5.2cm}|L{5.6cm}|}
\hline
\textbf{Criterio} & \textbf{Definición en el objetivo} & \textbf{Evidencia obtenida} \\ \hline
Específico & Dataset de tráfico HTTP etiquetado con tráfico legítimo y las cuatro categorías &
Corpus M1 con cinco clases presentes en las tres particiones \\ \hline
Medible & Mínimo 100{,}000 muestras etiquetadas & 383{,}202 muestras (Cuadro~\ref{tab:res-particiones}) \\ \hline
Alcanzable & Fuentes públicas y conjuntos sintéticos & Siete fuentes públicas auditadas en el
corpus, más dos listas públicas de payloads y conjuntos propios
(Cuadro~\ref{tab:met-origenes}, Figura~\ref{fig:res-auditoria}) \\ \hline
Relevante & Base de entrenamiento y validación representativa & Partición estratificada por clase y
fuente con verificación de integridad \\ \hline
Temporal & Fase 1 (mayo de 2026) & Corpus inicial en junio de 2026; regenerado hasta septiembre de
2026 \\ \hline
\end{tabular}
\caption[Criterios SMART del objetivo específico 2]{Criterios SMART del objetivo específico 2 y evidencia asociada}
\label{tab:res-smart-oe2}
\end{table}
{\footnotesize Nota. Elaboración propia.}

El Cuadro~\ref{tab:res-smart-oe2} resume el objetivo. El Cuadro~\ref{tab:res-particiones} presenta
el tamaño de cada partición del corpus vigente: 383{,}202 muestras etiquetadas, 3.8 veces el mínimo
de 100{,}000 establecido. Las particiones de validación y de prueba tienen el mismo tamaño
(57{,}481 filas) por la proporción 70/15/15.

\begin{table}[H]
\centering
\small
\begin{tabular}{|L{4.0cm}|L{2.8cm}|L{2.8cm}|}
\hline
\textbf{Partición} & \textbf{Filas} & \textbf{Proporción} \\ \hline
Entrenamiento & 268{,}240 & 70~\% \\ \hline
Validación & 57{,}481 & 15~\% \\ \hline
Prueba & 57{,}481 & 15~\% \\ \hline
\textbf{Total} & \textbf{383{,}202} & \textbf{100~\%} \\ \hline
Criterio ($\geq$100{,}000) & \multicolumn{2}{L{6.0cm}|}{Cumple} \\ \hline
\end{tabular}
\caption[Tamaño de las particiones del corpus]{Tamaño de las particiones del corpus etiquetado}
\label{tab:res-particiones}
\end{table}
{\footnotesize Nota. Elaboración propia con datos de \texttt{training/results/v11\_test\_results.json}
y conteo directo de \texttt{training/splits/*.parquet}.}

El Cuadro~\ref{tab:res-clases} presenta la distribución por clase en el conjunto de prueba. El
corpus no está balanceado en conteos: SQLi representa el 59.3~\% de la partición y Command
Injection el 2.6~\%. El desbalance se trató con ponderación de clases en el entrenamiento. Como
referencia, antes de la deduplicación y el filtrado las fuentes crudas sumaban 974{,}718 filas
(57.8~\% legítimas, 30.3~\% SQLi, 7.1~\% Path Traversal, 3.8~\% XSS y 0.9~\% Command Injection).

\begin{table}[H]
\centering
\small
\begin{tabular}{|L{4.0cm}|L{2.8cm}|L{2.8cm}|}
\hline
\textbf{Clase} & \textbf{Filas} & \textbf{Proporción} \\ \hline
\texttt{sqli} & 34{,}107 & 59.3~\% \\ \hline
\texttt{benign} & 14{,}886 & 25.9~\% \\ \hline
\texttt{xss} & 4{,}491 & 7.8~\% \\ \hline
\texttt{path\_traversal} & 2{,}529 & 4.4~\% \\ \hline
\texttt{cmdi} & 1{,}468 & 2.6~\% \\ \hline
\textbf{Total} & \textbf{57{,}481} & \textbf{100~\%} \\ \hline
\end{tabular}
\caption[Distribución por clase del conjunto de prueba]{Distribución por clase del conjunto de prueba}
\label{tab:res-clases}
\end{table}
{\footnotesize Nota. Elaboración propia con datos de \texttt{training/results/v11\_test\_results.json}.}

La Figura~\ref{fig:res-train-dist} muestra la distribución por clase de la partición de
entrenamiento en la generación del corpus previa a su regeneración del 12 de septiembre de 2026,
con el mismo orden de magnitud por clase que la partición de prueba vigente.

\begin{figure}[H]
\centering
\includegraphics[width=0.62\textwidth]{resultados/obj2-b.png}
\caption[Distribución por clase de la partición de entrenamiento]{Carga de características y distribución por clase de la partición de entrenamiento}
\label{fig:res-train-dist}
\end{figure}
{\footnotesize Nota. Captura de notebook de entrenamiento, repositorio \texttt{logSguarDian}.
Generación previa a la regeneración del 12 de septiembre de 2026 (268{,}064 filas de entrenamiento y
57{,}443 de validación); el tamaño vigente es el del Cuadro~\ref{tab:res-particiones}.}

La Figura~\ref{fig:res-auditoria} muestra la tabla de auditoría de las nueve fuentes candidatas,
con su licencia verificada, DOI, formato, número de registros, categorías cubiertas y decisión de
inclusión. Se aprobaron ocho y se excluyó AIT-LDSv2; siete de las aprobadas aportan filas al corpus
(Cuadro~\ref{tab:met-origenes}). La Figura~\ref{fig:res-parsers}
muestra los nueve analizadores por fuente y el validador del esquema canónico, junto con los
documentos de metodología de etiquetado, deduplicación y muestreo.

\begin{figure}[H]
\centering
\includegraphics[width=\textwidth]{resultados/obj1-c.png}
\caption[Auditoría de fuentes del dataset]{Tabla de auditoría de las fuentes del dataset}
\label{fig:res-auditoria}
\end{figure}
{\footnotesize Nota. Captura de \texttt{docs/dataset-audit.md}, repositorio \texttt{logSguarDian}.}

\begin{figure}[H]
\centering
\includegraphics[width=\textwidth]{resultados/obj2-e.png}
\caption[Analizadores por fuente y documentación del etiquetado]{Analizadores por fuente y documentación de la metodología de etiquetado}
\label{fig:res-parsers}
\end{figure}
{\footnotesize Nota. Captura de terminal, repositorio \texttt{logSguarDian}.}

La Figura~\ref{fig:res-lock} muestra el historial de commits de bloqueo del conjunto de prueba, desde el inicial del 11 de junio
hasta el de \texttt{rf\_v11} del 30 de agosto de 2026. No hay un único bloqueo por generación: el valor del bloqueo cambió nueve
veces en total, la última el 23 de septiembre de 2026 (hash \texttt{7486c258...}), con el mismo
cambio repartido en tres commits, y la lectura del 26 de septiembre verificó ese hash antes de cualquier
inferencia.

\begin{figure}[H]
\centering
\includegraphics[width=0.78\textwidth]{resultados/obj2-f.png}
\caption[Historial de bloqueos del conjunto de prueba]{Historial de commits de bloqueo del conjunto de prueba}
\label{fig:res-lock}
\end{figure}
{\footnotesize Nota. Captura de terminal, repositorio \texttt{logSguarDian}. El hash mostrado
(\texttt{a3140ab3...}) corresponde al bloqueo del 30 de agosto de 2026, previo al bloqueo vigente.}

% ==============================================================================
\section{Objetivo específico 3: efectividad del sistema de detección}
\label{sec:res-oe3}

\begin{table}[H]
\centering
\small
\begin{tabular}{|L{2.2cm}|L{5.2cm}|L{5.6cm}|}
\hline
\textbf{Criterio} & \textbf{Definición en el objetivo} & \textbf{Evidencia obtenida} \\ \hline
Específico & Precisión, recall y F1 por categoría de ataque & Reporte de clasificación de
\texttt{rf\_v11} sobre el conjunto de prueba bloqueado (Cuadro~\ref{tab:res-rf}) \\ \hline
Medible & F1 $\geq$0.80 en al menos 3 de 4 categorías; cobertura $\geq$80~\% por categoría & F1
entre 0.9239 y 0.9956 (4/4); cobertura entre 92.5~\% y 100~\% en 3 categorías con el paquete 0.1.1 (85.0~\% en Path Traversal con la 0.1.0) \\ \hline
Alcanzable & Conjunto de prueba bloqueado y entorno Docker reproducible & Lectura de prueba con
verificación de hash; corpus limpio reproducible con semilla fija \\ \hline
Relevante & Balance entre falsos negativos y falsos positivos sin analista dedicado & Falsos
positivos medidos en el conjunto de prueba, en la suite E2E y con navegadores reales \\ \hline
Temporal & Fases 4 y 5 (agosto a septiembre de 2026) & Lectura de prueba el 26 de septiembre;
corpus limpio el 24 de septiembre de 2026 \\ \hline
\end{tabular}
\caption[Criterios SMART del objetivo específico 3]{Criterios SMART del objetivo específico 3 y evidencia asociada}
\label{tab:res-smart-oe3}
\end{table}
{\footnotesize Nota. Elaboración propia.}

El Cuadro~\ref{tab:res-smart-oe3} resume el objetivo según los criterios SMART.

\subsection{Métricas de clasificación sobre el conjunto de prueba}

El Cuadro~\ref{tab:res-rf} presenta las métricas del Random Forest \texttt{rf\_v11} (par C) sobre
el conjunto de prueba bloqueado (57{,}481 filas), obtenidas en la lectura del 26 de septiembre de
2026, que verificó el hash del conjunto y las sumas de verificación de los binarios evaluados. Fue la tercera lectura
del conjunto de prueba para la generación \texttt{rf\_v11}; las dos anteriores se describen en la
Sección~\ref{subsec:met-f2}. Las métricas asignan a cada petición la clase de mayor probabilidad,
sin aplicar el umbral de bloqueo de 0.35. Las cuatro categorías de ataque superaron el
criterio de F1 $\geq$0.80. El F1 más bajo fue el de Command Injection (0.9239) y el más alto el de
SQLi (0.9956). La exactitud global fue de 0.9932 y el F1 macro de 0.9776. En la clase legítima, 23
de 14{,}886 peticiones (0.15~\%) se clasificaron como ataque.

\begin{table}[H]
\centering
\small
\begin{tabular}{|L{3.2cm}|L{1.8cm}|L{1.8cm}|L{1.8cm}|L{1.8cm}|L{1.6cm}|}
\hline
\textbf{Clase} & \textbf{Precisión} & \textbf{Recall} & \textbf{F1} & \textbf{AUC-ROC} & \textbf{Soporte} \\ \hline
\texttt{cmdi} & 0.9015 & 0.9475 & 0.9239 & 0.99901 & 1{,}468 \\ \hline
\texttt{path\_traversal} & 0.9853 & 0.9798 & 0.9826 & 0.99899 & 2{,}529 \\ \hline
\texttt{sqli} & 0.9955 & 0.9956 & 0.9956 & 0.99982 & 34{,}107 \\ \hline
\texttt{xss} & 0.9950 & 0.9804 & 0.9877 & 0.99930 & 4{,}491 \\ \hline
\texttt{benign} & 0.9984 & 0.9985 & 0.9984 & 0.99998 & 14{,}886 \\ \hline
\textbf{Macro} & 0.9751 & 0.9804 & \textbf{0.9776} & --- & 57{,}481 \\ \hline
\multicolumn{3}{|L{7.6cm}|}{Criterio (F1 $\geq$0.80 en $\geq$3 de 4)} & \multicolumn{3}{L{6.0cm}|}{Cumple (4 de 4)} \\ \hline
\end{tabular}
\caption[Métricas del Random Forest sobre el conjunto de prueba]{Métricas por clase del Random Forest (\texttt{rf\_v11}) sobre el conjunto de prueba bloqueado}
\label{tab:res-rf}
\end{table}
{\footnotesize Nota. Elaboración propia con datos de \texttt{training/results/v11\_test\_results.json}
(26 de septiembre de 2026). AUC-ROC en esquema uno contra el resto.}

Como valor inicial, la Figura~\ref{fig:res-rfv3} muestra las métricas de \texttt{rf\_v3} sobre el
conjunto de prueba de su generación (59{,}947 filas), con un F1 macro de 0.9682 y un F1 de 0.8954
en Command Injection. Entre \texttt{rf\_v3} y \texttt{rf\_v11}, el F1 de Command Injection pasó de
0.8954 a 0.9239, el de Path Traversal de 0.9671 a 0.9826, el de XSS de 0.9844 a 0.9877 y el de SQLi
de 0.9955 a 0.9956. Las particiones de ambas generaciones no son idénticas.

\begin{figure}[H]
\centering
\includegraphics[width=0.66\textwidth]{resultados/obj3-a.png}
\caption[Métricas iniciales de rf\_v3]{Métricas por clase de \texttt{rf\_v3} sobre el conjunto de prueba de su generación (valor inicial)}
\label{fig:res-rfv3}
\end{figure}
{\footnotesize Nota. Captura de \texttt{docs/decision-policy.md}, sección 2.1, repositorio \texttt{logSguarDian}.}

La Figura~\ref{fig:res-cmdi-evol} muestra una de las iteraciones de la Fase~5 sobre la categoría
de menor desempeño: la incorporación de 2{,}470 payloads reales y diversos de Command Injection en
\texttt{v9}, que elevó el F1 de esa clase en validación de 0.89--0.90 a 0.9369, después de que el
sobremuestreo sintético (SMOTE) produjera una mejora de solo 0.015.

\begin{figure}[H]
\centering
\includegraphics[width=0.7\textwidth]{resultados/obj2-d.png}
\caption[Evolución del F1 de Command Injection con datos reales diversos]{Evolución del desempeño en Command Injection al incorporar datos reales diversos (\texttt{v9})}
\label{fig:res-cmdi-evol}
\end{figure}
{\footnotesize Nota. Captura de \texttt{docs/limitations.md}, sección 1.1, repositorio \texttt{logSguarDian}.}

La Figura~\ref{fig:res-confusion} presenta la matriz de confusión del Random Forest sobre la
partición de validación. La confusión más frecuente fue la de SQLi clasificada como Command
Injection (106 casos), seguida de XSS clasificada como SQLi (56 casos).

\begin{figure}[H]
\centering
\includegraphics[width=0.74\textwidth]{resultados/rf_confusion_matrix.png}
\caption[Matriz de confusión del Random Forest]{Matriz de confusión del Random Forest sobre la partición de validación}
\label{fig:res-confusion}
\end{figure}
{\footnotesize Nota. Generada con scikit-learn \cite{pedregosa2011}; archivo
\texttt{training/results/rf\_confusion\_matrix.png}, 26 de septiembre de 2026.}

El Cuadro~\ref{tab:res-if} presenta las métricas del Isolation Forest sobre el mismo conjunto de
prueba, con el umbral fijado en sus metadatos (0.00421), seleccionado sobre validación. La tasa de detección de anomalías fue de
91.25~\% y la de falsos positivos de 5.46~\%; ambos criterios se cumplen. La
Figura~\ref{fig:res-ifcurve} muestra la curva de recall frente a la tasa de falsos positivos y el
punto de operación seleccionado.

\begin{table}[H]
\centering
\small
\begin{tabular}{|L{4.6cm}|L{2.6cm}|L{2.6cm}|L{2.2cm}|}
\hline
\textbf{Indicador} & \textbf{Valor} & \textbf{Criterio} & \textbf{Estado} \\ \hline
Detección de anomalías (recall) & 91.25~\% (38{,}868 / 42{,}595) & $\geq$50~\% & Cumple \\ \hline
Tasa de falsos positivos & 5.46~\% (813 / 14{,}886) & $\leq$10~\% & Cumple \\ \hline
Precisión & 97.95~\% & --- & --- \\ \hline
Recall por clase: \texttt{cmdi} / \texttt{path\_traversal} & 95.64~\% / 99.68~\% & --- & --- \\ \hline
Recall por clase: \texttt{sqli} / \texttt{xss} & 89.54~\% / 98.02~\% & --- & --- \\ \hline
\end{tabular}
\caption[Métricas del Isolation Forest sobre el conjunto de prueba]{Métricas del Isolation Forest (\texttt{if\_v10}) sobre el conjunto de prueba bloqueado}
\label{tab:res-if}
\end{table}
{\footnotesize Nota. Elaboración propia con datos de \texttt{training/results/v11\_test\_results.json}.}

\begin{figure}[H]
\centering
\includegraphics[width=0.8\textwidth]{resultados/if_recall_fp_curve.png}
\caption[Curva de recall frente a falsos positivos del Isolation Forest]{Recall frente a tasa de falsos positivos del Isolation Forest y punto de operación seleccionado}
\label{fig:res-ifcurve}
\end{figure}
{\footnotesize Nota. Archivo \texttt{training/results/if\_recall\_fp\_curve.png}, 26 de septiembre de 2026.}

\subsection{Cobertura sobre el corpus limpio}
\label{subsec:res-oe3-cobertura}

El Cuadro~\ref{tab:res-corpus} presenta la construcción del corpus limpio de la Ronda~5 y el
Cuadro~\ref{tab:res-cobertura} la cobertura obtenida con la Configuración~2. El corpus contiene 601
peticiones: 200 por cada una de las categorías SQLi, Path Traversal y Command Injection, y un
payload de XSS confirmado con OWASP ZAP. El filtro de exclusión descartó 3 candidatos de Path Traversal por
coincidencia exacta con una fuente de entrenamiento; en SQLi y Command Injection no hubo descartes.
Con el paquete 0.1.0 (par~A), la cobertura fue de 100~\% en SQLi, 91.5~\% en Command Injection
y 85.0~\% en Path Traversal, con intervalos de Wilson \cite{wilson1927probable} al 95~\% de 98.1--100.0~\%, 86.8--94.6~\% y
79.4--89.3~\%, respectivamente. Con el tarball de validación 22550ea (par~C) y el mismo corpus, la
cobertura fue de 99.5~\%, 96.0~\% y 95.0~\%, con intervalos de 97.2--99.9~\%, 92.3--98.0~\% y
91.0--97.3~\%. Las tres categorías superan el criterio de 80~\% con ambos pares, pero con el par~A
el límite inferior del intervalo de Path Traversal queda por debajo de ese umbral. Con el paquete
0.1.1, publicado después, instalado desde el registro de npm y medido con el ejecutor corregido, la
cobertura fue de 99.5~\% en SQLi, 100.0~\% en Path Traversal y 92.5~\% en Command Injection, con
intervalos de 97.2--99.9~\%, 98.1--100.0~\% y 88.0--95.4~\%. La medición del par~C no es una
evaluación estrictamente ciega, y la del paquete 0.1.1 se hizo sobre el corpus que reveló los
defectos de medición (Sección~\ref{sec:met-validez}). Cada medición corresponde a una sola corrida.
Dos peticiones se contaron como no detectadas por fallos y no por el veredicto del modelo:
\texttt{cmdi\_94} terminó por tiempo de espera en la corrida del par~A y \texttt{sqli\_75} devolvió
HTTP 500 en la del par~C. En XSS, el \textit{Active Scan} de ZAP, dirigido por su API REST y con
sesión autenticada, produjo 538 alertas, de las cuales una sola es de XSS: \texttt{POST /profile},
campo \texttt{bio}, con el payload \texttt{</div><scrIpt>alert(1);</scRipt><div>} y riesgo alto.
Los campos \texttt{content} y \texttt{title} de \texttt{POST /posts} se escanearon con la misma
política y no resultaron vulnerables. Aunque la regla de ZAP se llama \textit{Reflected}, la
vulnerabilidad es de XSS almacenado, porque \texttt{bio} se guarda y se vuelve a mostrar en cada
vista del perfil. Ese payload fue bloqueado por ambos pares de modelos (1 de 1), de modo que la
cobertura de XSS tiene n=1 y no es comparable en peso estadístico con las otras tres categorías.

Los fallos de Path Traversal se concentraban en las categorías de codificación. El desglose por
estilo de codificación, calculado con \texttt{training/audit/breakdown\_config2.py} sobre los
resultados de la Configuración~2 con el ejecutor anterior, muestra que el paquete 0.1.0 (par~A) no
bloqueó 30 payloads: 21 de los 43 con codificación URL o mixta, 7 de los 23 con doble codificación
y 2 de los 19 con UTF-8 sobrelargo. Con el modelo vigente (par~C) los fallos eran 10: 9 de los 23
con doble codificación y 1 de los 19 con UTF-8 sobrelargo.

Esos fallos tenían dos causas independientes, y ninguna era una limitación de detección. La primera
era un defecto del middleware, que reconstruía la cadena de consulta a partir de \texttt{req.query},
ya decodificada una vez por Express, y le añadía una capa de codificación. La segunda era un defecto
del ejecutor de pruebas, que volvía a codificar los payloads de Path Traversal y Command Injection
que el generador ya había codificado. El Cuadro~\ref{tab:res-pt-causas} muestra que ninguna
corrección bastó por sí sola: con el paquete 0.1.1, que incluye la corrección del middleware, y el
ejecutor anterior, los 10 fallos de Path Traversal fueron los mismos que con el tarball sin esa
corrección (payload por payload), y desaparecieron solo al corregir también el ejecutor. Con ambas
correcciones, el paquete 0.1.1, instalado desde el registro de npm, bloqueó 200 de 200 payloads de
Path Traversal (intervalo de Wilson al 95~\%: 98.1--100.0~\%).

\begin{table}[H]
\centering
\small
\begin{tabular}{|L{5.4cm}|L{2.4cm}|L{2.8cm}|L{2.4cm}|}
\hline
\textbf{Paquete y ejecutor} & \textbf{Fallos de Path Traversal} & \textbf{Fallos de Command Injection} & \textbf{Fallos de SQLi} \\ \hline
0.1.0 (par~A), ejecutor anterior & 30 & 17 & 0 \\ \hline
Tarball 22550ea (par~C), ejecutor anterior & 10 & 8 & 1 \\ \hline
0.1.1 (par~C), ejecutor anterior & 10 & 8 & 1 \\ \hline
0.1.1 (par~C), ejecutor corregido & 0 & 15 (1 por tiempo de espera) & 1 \\ \hline
\end{tabular}
\caption[Fallos por paquete y versión del ejecutor]{Payloads no bloqueados de cada categoría, por paquete y por versión del ejecutor de pruebas}
\label{tab:res-pt-causas}
\end{table}
{\footnotesize Nota. Elaboración propia con datos de los archivos de resultados de la Configuración~2
(repositorio \texttt{logSguarDian-vulnerable-project}). No se midió el paquete 0.1.0 con el
ejecutor corregido, de modo que la mejora de 85.0~\% a 100.0~\% en Path Traversal no puede
repartirse entre el cambio de modelo, la corrección del middleware y la del ejecutor.}

Con el ejecutor corregido, la cobertura de Command Injection pasó de 96.0~\% a 92.5~\% (185 de 200;
intervalo de Wilson al 95~\%: 88.0--95.4~\%). Esto no es un empeoramiento de la detección. Los
generadores producen 105 de los 200 payloads de Command Injection con separadores ya codificados
(\texttt{\%26}, \texttt{\%7c}, \texttt{\%0a}, \texttt{\%3B}), y el ejecutor anterior los codificaba
una segunda vez. De esos, 10 se bloqueaban en su forma corrompida y no se bloquean al enviarse
correctamente (todos reciben HTTP 200); otros 3 hicieron el recorrido contrario, y 1 terminó por
tiempo de espera. El resultado neto es de 7 payloads bloqueados menos. El 96.0~\% estaba inflado
por ese defecto, y el 92.5~\% es la primera medición correcta de esta técnica de evasión: la
librería no detecta 10 payloads de Command Injection con separadores codificados, un límite real de
la detección.

La corrección del ejecutor afecta también a las demás mediciones de Path Traversal y Command
Injection de este capítulo con el ejecutor anterior, que no se repitieron.

El filtro de exclusión descarta únicamente las coincidencias exactas del campo completo (y, en
Command Injection, las de plantilla), por lo que no garantiza la ausencia de solapamiento por
forma normalizada. La medición de solapamiento del trabajo de graduación de Diego Valenzuela \cite{valenzuela2026ia}
(sección VIII.3.6), con \texttt{training/leakage\_check.py} (commit \texttt{22550ea}) sobre las
particiones de entrenamiento, validación y prueba, da los siguientes resultados. Con igualdad
normalizada del campo completo (decodificación de URL, minúsculas y dígitos sustituidos), 4 de los
200 payloads de SQLi coinciden con alguna partición (4 con entrenamiento, 1 con validación y 2
con prueba) y ninguno de Path Traversal ni de Command Injection. Con la regla de subcadena, un
criterio permisivo que su autor califica de límite superior y no de medida de solapamiento, Path
Traversal alcanza 7 coincidencias y Command Injection 2. La cobertura de SQLi reportada incluye,
por lo tanto, esos 4 payloads. Las cifras anteriores a esta ronda no son comparables: el corpus
de la Ronda~4 tiene 200 de 200 payloads de Command Injection y 7 de 113 de XSS que coinciden por
igualdad normalizada, y ninguno de SQLi ni de Path Traversal; por subcadena, 26 de SQLi y 13 de
XSS.

\begin{table}[H]
\centering
\small
\begin{tabular}{|L{2.7cm}|L{4.2cm}|L{1.8cm}|L{1.8cm}|L{1.4cm}|}
\hline
\textbf{Categoría} & \textbf{Fuente} & \textbf{Candidatos} & \textbf{Descartados} & \textbf{Finales} \\ \hline
SQLi & sqlmap v1.10.9 & 1{,}995 & 0 & 200 \\ \hline
Path Traversal & Generador propio & 400 & 3 & 200 \\ \hline
Command Injection & Generador propio (no derivado de SecLists) & 400 & 0 & 200 \\ \hline
XSS & OWASP ZAP (\textit{Active Scan}, API REST) & 1 alerta confirmada & 0 & 1 \\ \hline
\end{tabular}
\caption[Construcción del corpus limpio]{Construcción del corpus limpio de evaluación (Ronda~5, semilla 20260923)}
\label{tab:res-corpus}
\end{table}
{\footnotesize Nota. Elaboración propia con datos de \texttt{docs/config3b-results.md} y
\texttt{attack-sim/exclusion\_report.json} (repositorio \texttt{logSguarDian-vulnerable-project}).}

\begin{table}[H]
\centering
\small
\begin{tabular}{|L{2.2cm}|L{1.9cm}|L{2.0cm}|L{2.0cm}|L{2.2cm}|L{1.9cm}|}
\hline
\textbf{Categoría} & \textbf{Ronda 4 (par A)} & \textbf{Ronda 5, 0.1.0 (par A)\textsuperscript{a}} & \textbf{Ronda 5, tarball 22550ea (par C)\textsuperscript{a}} & \textbf{Ronda 5, 0.1.1 (par C), ejecutor corregido\textsuperscript{b}} & \textbf{Criterio ($\geq$80~\%)} \\ \hline
SQLi & 76/77 (98.7~\%) & 200/200 (100.0~\%) & 199/200 (99.5~\%) & 199/200 (99.5~\%) & Cumple \\ \hline
Path Traversal & 197/200 (98.5~\%) & 170/200 (85.0~\%) & 190/200 (95.0~\%) & 200/200 (100.0~\%) & Cumple \\ \hline
Command Injection & 200/200 (100.0~\%) & 183/200 (91.5~\%) & 192/200 (96.0~\%) & 185/200 (92.5~\%) & Cumple \\ \hline
XSS & 110/113 (97.3~\%) & 1/1 (n=1) & 1/1 (n=1) & 1/1 (n=1) & Cumple, con n=1 \\ \hline
\end{tabular}
\caption[Cobertura por categoría con logsguardian]{Cobertura por categoría con \textit{logSguarDian} activo (Configuración~2)}
\label{tab:res-cobertura}
\end{table}
{\footnotesize Nota. Elaboración propia con datos de \texttt{docs/config3b-results.md} y
\texttt{attack-sim/results\_config2-develop-tarball-20261009.json} (repositorio
\texttt{logSguarDian-vulnerable-project}). El corpus de la Ronda~4 tiene fuga confirmada en Command Injection (200/200 payloads coinciden por plantilla con
una fuente de entrenamiento). \textsuperscript{a}Medido con la versión del ejecutor anterior a la
corrección descrita en el texto, que añadía una capa de codificación a los payloads de Path
Traversal y Command Injection ya codificados. \textsuperscript{b}Paquete instalado desde el registro de
npm (\texttt{attack-sim/results\_config2-published-0.1.1-harness-fixed-20261010.json}); es la única
columna medida con el ejecutor corregido. Las columnas de la Ronda~5 con el ejecutor anterior no
son comparables con ella.}

\subsection{Comparación con un WAF y defensa en profundidad}

Las Rondas 1 a 3, con un corpus pequeño construido a mano, produjeron tres resultados: ModSecurity
con OWASP CRS en nivel de paranoia 1 bloqueó los 9 ataques originales y las 14 variantes de evasión
(0 de 14 llegaron a la aplicación); con el WAF en modo de solo detección, \textit{logSguarDian} bloqueó inicialmente 16 de los 23
payloads; los 7 no detectados eran de XSS enviados como \texttt{multipart/form-data}, que la
librería no inspeccionaba porque el analizador de ese tipo de cuerpo se montaba por ruta, después
del middleware. Tras montar un analizador global antes del middleware, \textit{logSguarDian}
bloqueó los 23 (23/23), y esa misma corrección evidenció una vulnerabilidad de XSS almacenado en la
aplicación de referencia.

La Figura~\ref{fig:res-ronda4} presenta la Ronda~4, que comparó las cuatro configuraciones sobre
590 payloads de SecLists. La configuración sin protección no bloqueó ningún ataque (0/590), el WAF
solo bloqueó 550/590 (93.2~\%), \textit{logSguarDian} solo 583/590 (98.8~\%) y ambas capas 587/590
(99.5~\%). De los 40 ataques que el WAF no bloqueó, \textit{logSguarDian} bloqueó 37. Los resultados
de Command Injection de esta ronda están afectados por la fuga descrita en la nota del
Cuadro~\ref{tab:res-cobertura}; la comparación se repitió después con el corpus limpio de la Ronda~5 (Cuadro~\ref{tab:res-ronda5}). La
Configuración~2 se repitió sobre los mismos 590 payloads el 9 de septiembre de 2026 con el par~B y
también bloqueó 583 de 590, aunque 10 payloads cambiaron de estado entre ambas corridas. En esa
repetición, los 200 payloads de Command Injection se registraron con la clase \texttt{sqli}; la
confusión afecta la clase reportada en el registro, no la decisión de bloqueo.

\begin{figure}[H]
\centering
\includegraphics[width=0.88\textwidth]{resultados/obj3-g.png}
\caption[Comparación de configuraciones en la Ronda 4]{Ataques bloqueados por configuración y análisis de defensa en profundidad (Ronda~4)}
\label{fig:res-ronda4}
\end{figure}
{\footnotesize Nota. Captura de \texttt{docs/config3b-results.md}, repositorio
\texttt{logSguarDian-vulnerable-project}.}

El Cuadro~\ref{tab:res-ronda5} repite la comparación con el corpus limpio de la Ronda~5, que no
tiene la fuga de la Ronda~4. Sus cifras de Path Traversal y de Command Injection se midieron con el
ejecutor anterior, que alteraba los payloads ya codificados y afecta a las cuatro configuraciones,
incluidas las del WAF; solo la Configuración~2 se repitió con el ejecutor corregido
(Cuadro~\ref{tab:res-cobertura}), de modo que esas dos filas deben leerse como provisionales. La Configuración~1 bloqueó 0 de 601 payloads, lo que confirma que el
corpus y el ejecutor miden lo que declaran. El WAF solo (3a) bloqueó el 99.5~\% de SQLi, una cifra equivalente a la de
\textit{logSguarDian} solo (2, 100~\%), y el 98.5~\% de Path Traversal, más que la librería
(85.0~\%), pero solo el 68.0~\% de Command Injection, frente al 91.5~\% de la librería. La combinación de ambas capas (3b)
bloqueó el 100~\% de SQLi y de Path Traversal y el 98.5~\% de Command Injection. De los 68 ataques
que el WAF solo no bloqueó (1 de SQLi, 3 de Path Traversal y 64 de Command Injection), la
Configuración~3b bloqueó 65 y la Configuración~2, 64. La Configuración~2 se midió con un tarball que contiene el par~A y de la 3b no consta el paquete
exacto; la 3a no incorpora la librería y se midió en una rama distinta de la aplicación. En XSS,
las tres configuraciones protegidas bloquearon el único payload, con n=1.

\begin{table}[H]
\centering
\small
\begin{tabular}{|L{2.7cm}|L{2.0cm}|L{2.5cm}|L{2.5cm}|L{2.5cm}|}
\hline
\textbf{Categoría} & \textbf{Config.~1 (sin protección)} & \textbf{Config.~2 (solo \textit{logSguarDian})} & \textbf{Config.~3a (solo WAF)} & \textbf{Config.~3b (ambas capas)} \\ \hline
SQLi & 0/200 (0.0~\%) & 200/200 (100.0~\%) & 199/200 (99.5~\%) & 200/200 (100.0~\%) \\ \hline
Path Traversal & 0/200 (0.0~\%) & 170/200 (85.0~\%) & 197/200 (98.5~\%) & 200/200 (100.0~\%) \\ \hline
Command Injection & 0/200 (0.0~\%) & 183/200 (91.5~\%) & 136/200 (68.0~\%) & 197/200 (98.5~\%) \\ \hline
XSS (n=1) & 0/1 & 1/1 & 1/1 & 1/1 \\ \hline
\end{tabular}
\caption[Cuatro configuraciones sobre el corpus limpio de la Ronda 5]{Payloads bloqueados por configuración sobre el corpus limpio de la Ronda~5 (601 peticiones)}
\label{tab:res-ronda5}
\end{table}
{\footnotesize Nota. Elaboración propia con datos de
\texttt{attack-sim/results\_config\{1,2-published-with-real-xss,3a,3b\}-round5-20261010.json}
(repositorio \texttt{logSguarDian-vulnerable-project}); conteos recalculados a partir de los
archivos de resultados. Configuración~2: tarball con el par~A (ver Sección~\ref{subsec:met-f2}); de la Configuración~3b no consta el paquete exacto (los resultados no lo registran), y se atribuye al par~A. Medido con el ejecutor
anterior a la corrección; las Configuraciones 3a y 3b no se repitieron con el ejecutor corregido.}

\subsection{Falsos positivos con navegadores reales}

El Cuadro~\ref{tab:res-ua} presenta el estudio de falsos positivos con cadenas
\texttt{User-Agent} de cinco navegadores reales sobre tres rutas sin consulta ni cuerpo de la
aplicación de referencia. De las 15 combinaciones legítimas, 13 (86.7~\%) se bloquearon. Solo
Firefox en macOS evitó el bloqueo en dos rutas (\texttt{/login} y \texttt{/posts}). Como contraste,
la suite E2E registró 2 de 100 peticiones legítimas bloqueadas (Figura~\ref{fig:res-e2e}).

\begin{table}[H]
\centering
\small
\begin{tabular}{|L{3.6cm}|L{2.2cm}|L{2.2cm}|L{2.2cm}|L{2.0cm}|}
\hline
\textbf{\texttt{User-Agent}} & \textbf{\texttt{/login}} & \textbf{\texttt{/posts}} & \textbf{\texttt{/profile}} & \textbf{Bloqueadas} \\ \hline
Chrome (macOS) & 403 & 403 & 403 & 3/3 \\ \hline
Chrome (Windows) & 403 & 403 & 403 & 3/3 \\ \hline
Safari (macOS) & 403 & 403 & 403 & 3/3 \\ \hline
Safari (iOS) & 403 & 403 & 403 & 3/3 \\ \hline
Firefox (macOS) & 200 & 200 & 403 & 1/3 \\ \hline
\textbf{Total} & 4/5 & 4/5 & 5/5 & \textbf{13/15} \\ \hline
\end{tabular}
\caption[Falsos positivos con User-Agent de navegadores reales]{Respuesta del middleware a peticiones legítimas con \texttt{User-Agent} de navegadores reales}
\label{tab:res-ua}
\end{table}
{\footnotesize Nota. Elaboración propia con datos de \texttt{attack-sim/ua\_fp\_study\_results.json}
(repositorio \texttt{logSguarDian-vulnerable-project}). \texttt{/login} sin sesión; \texttt{/posts} y
\texttt{/profile} con sesión iniciada. Ejecutado con el par~A; atribución por la fecha del commit, ya
que el archivo de resultados no registra el modelo.}

% ==============================================================================
\section{Resumen de indicadores}
\label{sec:res-resumen}

El Cuadro~\ref{tab:res-resumen} consolida los indicadores definidos en el
Cuadro~\ref{tab:met-metricas} de la Metodología, con el valor medido y su estado frente al
criterio.

\begin{longtable}{|L{3.5cm}|L{1.2cm}|L{2.7cm}|L{3.0cm}|L{2.1cm}|}
\hline
\textbf{Indicador} & \textbf{Obj.} & \textbf{Criterio} & \textbf{Valor medido} & \textbf{Estado} \\ \hline
\endfirsthead
\multicolumn{5}{l}{\footnotesize\textit{(continuación)}} \\ \hline
\textbf{Indicador} & \textbf{Obj.} & \textbf{Criterio} & \textbf{Valor medido} & \textbf{Estado} \\ \hline
\endhead
\multicolumn{5}{r}{\footnotesize\textit{Continúa en la página siguiente}} \\
\endfoot
\endlastfoot
Vectores mapeados a OWASP, CWE y ATT\&CK & OE1 & 4/4 & 4/4 & Cumple \\ \hline
Tamaño del corpus & OE2 & $\geq$100{,}000 & 383{,}202 & Cumple \\ \hline
Clases representadas por partición & OE2 & 5/5 & 5/5 & Cumple \\ \hline
F1 por categoría & OE3 & $\geq$0.80 en $\geq$3 de 4 & 0.9239--0.9956 (4/4) & Cumple \\ \hline
AUC-ROC por categoría & OE3 & Referencia & 0.99899--0.99982 & --- \\ \hline
Detección del IF & OE3 & $\geq$50~\% & 91.25~\% & Cumple \\ \hline
Falsos positivos del IF & OE3 & $\leq$10~\% & 5.46~\% & Cumple \\ \hline
Cobertura SQLi / PT / CMDi (paquete 0.1.0, ejecutor anterior) & OE3 & $\geq$80~\% & 100 / 85.0 / 91.5~\% (PT, IC 95~\%: 79.4--89.3~\%) & Cumple; no robusto en PT \\ \hline
Cobertura SQLi / PT / CMDi (modelo vigente, tarball) & OE3 & $\geq$80~\% & 99.5 / 95.0 / 96.0~\% (PT, IC 95~\%: 91.0--97.3~\%) & Cumple \\ \hline
Cobertura SQLi / PT / CMDi (paquete 0.1.1, ejecutor corregido) & OE3 & $\geq$80~\% & 99.5 / 100.0 / 92.5~\% (PT, IC 95~\%: 98.1--100.0~\%; CMDi, IC 95~\%: 88.0--95.4~\%) & Cumple \\ \hline
Cobertura XSS & OE3 & $\geq$80~\% & 1/1 (n=1, ZAP) & Cumple, con n=1 \\ \hline
Falsos positivos con navegador & OE3 & Descriptivo & 13/15 combinaciones de \texttt{User-Agent} & --- \\ \hline
Paridad ONNX & OG & $<$0.1~\% & $\leq2.42\times10^{-7}$ & Cumple \\ \hline
Extracción de características (p95, ida y vuelta al hilo) & OG & $\leq$5~ms & 0.0273~ms & Cumple \\ \hline
$\Delta p95_{\mathrm{abs}}$ normal / ataque & OG & $\leq$15~ms & 12.08 / 12.85~ms & Cumple \\ \hline
$\Delta p95_{\mathrm{rel}}$ normal / ataque & OG & Complementario & +485.1 / +278.1~\% & --- \\ \hline
$\Delta p95_{\mathrm{abs}}$ volumen (con 7.00~\% de tiempos de espera) & OG & $\leq$15~ms (referencia) & 6{,}495.46~ms & No cumple \\ \hline
Memoria adicional bajo carga (normal / ataque) & OG & $\leq$150~MB (143.1~MiB) & +332.02 / +331.86~MiB & No cumple \\ \hline
CPU adicional bajo carga (normal / ataque) & OG & $\leq$15~\% & +614 / +530~\% & No cumple \\ \hline
Contribución en defensa en profundidad (Ronda~5, ejecutor anterior) & OG & Descriptivo & 65 de 68 ataques no bloqueados por el WAF (provisional) & --- \\ \hline
Bloqueo de legítimos en terceros (Cuadro~\ref{tab:res-terceros-ua}) & OG & Descriptivo & 92.5--100~\% (modo \texttt{block}); 97.1~\% veredicto en Juice Shop (modo \texttt{monitor}) & --- \\ \hline
\caption[Resumen de indicadores]{Resumen de indicadores: criterio, valor medido y estado}
\label{tab:res-resumen}
\end{longtable}
{\footnotesize Nota. Elaboración propia con datos de las secciones anteriores de este capítulo.}

\ifdefined\modoensamble
\else
	\section*{Referencias}
	\printbibliography[heading=none]
	\end{document}
\fi

\fi
%
% \newpage
% \ifdefined\parpordefecto
% 	\defaultparformat{j-capitulos}
% \else
% 	\input{j-capitulos}
% \fi
%
\ifdefined\CAPdiscusion
	\newpage
	% ==============================================================================
% ANÁLISIS DE RESULTADOS (antes "Discusión de resultados")
% ==============================================================================
% Interpreta la evidencia de k-resultados.tex. Solo componente de
% ciberseguridad.
%
% Enfoque (decisión del estudiante, 2026-10-09): la discusión argumenta el
% cumplimiento de los objetivos frente a los criterios de aceptación
% definidos en g-objetivos.tex y h-alcance.tex (Alcance técnico: extracción
% <= 5 ms y Delta p95 <= 15 ms en carga normal y de ataque; validación
% principal en entorno controlado, validación con startups complementaria).
% El escenario de volumen se evalúa contra el mismo criterio de 15 ms y no se cumple. Los hallazgos fuera de esos criterios (falsos
% positivos en aplicaciones de terceros) se presentan como condiciones de
% despliegue y trabajo futuro. Ninguna cifra cambia respecto a Resultados.
%
% "Variantes desconocidas de los vectores cubiertos" (objetivo general): la
% evidencia es la Ronda 5, corpus de payloads de las cuatro categorías
% generados con sqlmap y generadores propios y filtrados por exclusión
% contra el entrenamiento; la detección del sistema se cuenta como bloqueo
% o marca de anomalía (Procedimiento 4.4).
%
% kapoor2023: taxonomía verificada en el PDF del artículo (2026-10-09):
% [L2] "Model uses features that are not legitimate"; [L3] "Test set is not
% drawn from the distribution of scientific interest".
%
% Ronda 4: los 40 ataques que el WAF no bloqueó fueron 11 SQLi, 3 XSS y 26
% Path Traversal, ninguno de Command Injection.
\ifdefined\modoensamble
\else
	\input{0-preambulo_standalone}
\fi

\chapter{Análisis de resultados}
\label{cap:discusion}

Este capítulo interpreta la evidencia presentada en el Capítulo~\ref{cap:resultados} frente a los
criterios de aceptación definidos en los Objetivos y en el Alcance. Se analiza, objetivo por
objetivo, qué significan los resultados y en qué medida sustentan su cumplimiento; se contrastan
con los trabajos previos y los fundamentos del Marco Teórico; se analizan las condiciones de
despliegue y las limitaciones que acotan los hallazgos; se valida la solución frente a sus
criterios de aceptación; y se cierra con sugerencias que se desprenden de los resultados y que
dan paso a las conclusiones y a las recomendaciones. En cada apartado se distingue lo comprobado
de lo que es solo una explicación posible.

% ------------------------------------------------------------------------------
\section{Objetivo general: una librería RASP funcional para Node.js}

El objetivo general establece cuatro componentes verificables: una librería npm, un pipeline de
ingeniería de características, un modelo híbrido capaz de detectar y clasificar amenazas conocidas
y variantes desconocidas de los vectores cubiertos, y la operación en tiempo de ejecución sin
comprometer el rendimiento del servidor. Los resultados sustentan los tres primeros en lo esencial
y solo en parte el cuarto, como se detalla al final de esta sección.

\textbf{Librería y pipeline.} \texttt{logsguardian} está publicada en el registro público de
npm (0.1.0 y 0.1.1), se instala con un solo comando y no requiere infraestructura externa. Integra un extractor de
76 características implementado una sola vez en TypeScript, de modo que el entrenamiento y la
producción calculan el mismo vector, y un modelo híbrido ejecutado en hilos de trabajo que no
bloquean el Event Loop. La paridad de $\leq2.42\times10^{-7}$ entre los modelos de Python y sus
versiones ONNX garantiza que la serialización no altera las predicciones. Las evaluaciones
emplearon tres pares de modelos con la misma arquitectura de ejecución
(Cuadro~\ref{tab:pares-modelos}): el paquete 0.1.0 contiene \texttt{rf\_v10}/\texttt{if\_v9}, que
consume 67 y 61 características, y las métricas del conjunto de prueba corresponden a
\texttt{rf\_v11}/\texttt{if\_v10}, que consume 69 y 63 y es el contenido del paquete 0.1.1,
publicado el 11 de octubre de 2026 con la corrección del middleware. Como se discute en la
Sección~\ref{sec:dis-oe3}, las generaciones evaluadas cumplen el criterio de cobertura sobre
payloads no vistos, aunque con el paquete 0.1.0 el margen en Path Traversal es estrecho. Con la 0.1.1 solo se midió la cobertura de la Configuración~2 de la Ronda~5; las aplicaciones de
terceros y el rendimiento no se midieron con ella, de modo que lo que el paquete vigente haría en
esas pruebas queda por medir. Este resultado cubre la brecha identificada en los Antecedentes: las soluciones
RASP comerciales requieren licencias o plataformas externas, y la única alternativa abierta para
Node.js (node-rasp) exigía modificar el runtime y fue abandonada \cite{saparchive2019}.

\textbf{Amenazas conocidas.} La detección y clasificación de las cuatro categorías se sustenta en
el OE3 (Sección~\ref{sec:dis-oe3}): F1 de entre 0.9239 y 0.9956 sobre el conjunto de prueba.

\textbf{Variantes desconocidas de los vectores cubiertos.} La evidencia de esta capacidad es el
corpus limpio de la Ronda~5: 600 payloads de SQL Injection, Path Traversal y Command Injection
generados con sqlmap y con generadores propios, y filtrados para excluir toda coincidencia exacta
del campo completo con el entrenamiento. Son, por construcción, variantes que no coinciden de forma
exacta con las que el modelo vio (4 de los 200 payloads de SQL Injection coinciden con alguna
partición por igualdad normalizada, según la medición de solapamiento del trabajo de graduación de Diego Valenzuela \cite{valenzuela2026ia}, y ninguno de Path Traversal ni de Command Injection), con técnicas de codificación (doble codificación URL, UTF-8 sobrelargo, barra invertida,
byte nulo) y separadores de shell distintos de los del entrenamiento. El paquete 0.1.0 bloqueó
el 100~\% de las de SQL Injection, el 91.5~\% de las de Command Injection y el 85.0~\% de las de
Path Traversal, y el modelo vigente, el 99.5~\%, el 96.0~\% y el 95.0~\%. Esta capacidad corresponde
al Random Forest: la cobertura cuenta solo peticiones bloqueadas y el Isolation Forest no bloquea.
El Isolation Forest cumple sus dos criterios sobre el conjunto de prueba (91.25~\% de detección con
5.46~\% de falsos positivos), pero su contribución a la detección de variantes no vistas solo tiene
evidencia indirecta: en el conjunto de validación del componente de Inteligencia Artificial \cite{valenzuela2026ia} marca 19
de los 39 ataques que el Random Forest deja pasar, y bloquear sus marcas agregaría 866 benignos
bloqueados; esas cifras son optimistas, porque el umbral del Isolation Forest se eligió sobre ese
mismo conjunto, y no se evaluaron sobre el conjunto de prueba. Con el paquete 0.1.0, el límite inferior del intervalo de confianza de Path Traversal
(79.4~\%) quedaba por debajo del criterio, pero esa cifra incluía dos defectos de medición ya
corregidos; con el paquete 0.1.1 y el ejecutor corregido es de 100~\% (98.1--100.0~\%), medido
sobre el mismo corpus que reveló los defectos y, por tanto, sin independencia
(Sección~\ref{sec:met-validez}). Con esas salvedades, el sistema reconoce variantes no vistas
de los vectores cubiertos en una proporción superior al criterio de cobertura del 80~\%.

\textbf{Rendimiento del servidor anfitrión.} El Alcance fija dos criterios cuantitativos: la
extracción de características en 5~ms o menos por petición y un incremento de la latencia p95 de
15~ms como máximo bajo carga normal y bajo carga de ataque. Ambos se cumplen: la extracción,
medida con la ida y vuelta completa a su \texttt{worker\_thread}, tardó 0.0273~ms en el p95, más
de 180 veces por debajo del límite, y el incremento de latencia fue de
12.08~ms y 12.85~ms. El valor relativo de esos mismos incrementos (+485~\% y +278~\%) confirma por
qué el criterio se expresó en milisegundos: la aplicación de referencia responde en 2 a 5~ms, de
modo que cualquier costo fijo produce porcentajes elevados sin que el usuario perciba una
degradación proporcional; en una API con tiempos de respuesta de decenas o cientos de
milisegundos, 12~ms adicionales representan una fracción menor. Con ello, el criterio de latencia se cumple en carga normal y de ataque; en el escenario de
volumen no se cumple (\(\Delta p95\) de 6{,}495.46~ms y 7.00~\% de tiempos de espera).

\textbf{Costo de recursos.} El protocolo definió además dos indicadores de costo, que no se
cumplen: un incremento de memoria de 150~MB (143.1~MiB) como máximo y un incremento de CPU de
15~\% como máximo. Con el middleware activo, el contenedor usó 332.0~MiB más de memoria bajo carga
normal y 331.9~MiB más bajo carga de ataque, más del doble del límite, y el uso de CPU aumentó en
89 y 85 puntos porcentuales. El incremento de memoria es casi idéntico entre escenarios, por lo
que corresponde a un costo fijo de la arquitectura al activarse (el runtime nativo de ONNX y tres
hilos de inferencia con sus modelos cargados) y no a un crecimiento con la carga. El límite de
150~MB se fijó antes de medir la memoria nativa que reserva \texttt{onnxruntime-node} al cargar
modelos de árboles, proporcional a su número de nodos, y pensando en instancias de 1~GB; en una
instancia de ese tamaño, 332~MiB representan cerca de un tercio de la memoria disponible, un costo
que debe considerarse al dimensionar la infraestructura. El incremento de CPU, medido en un
ejecutor de 2~vCPU a 100 peticiones por segundo, equivale aproximadamente a un núcleo dedicado a
la inferencia. Ambos costos son objetivos directos de optimización para la siguiente versión.

\textbf{Veredicto sobre el objetivo general.} El objetivo general se cumple de forma parcial. Se
cumplen la librería publicada, el pipeline, la clasificación de las cuatro categorías, la cobertura
de variantes no vistas con el Random Forest (con las reservas indicadas) y el criterio de latencia
en carga normal y de ataque. No se cumplen los indicadores de memoria y CPU, ni el criterio de
latencia en el escenario de volumen, y el aporte del Isolation Forest a la detección de variantes
desconocidas solo tiene evidencia indirecta. Atribuir ese resultado a la arquitectura y no a un
modelo en particular es razonable para el costo de recursos, que es casi constante entre
escenarios, pero no está comprobado: el costo por componente no se midió.

% ------------------------------------------------------------------------------
\section{Objetivo específico 1: selección de vectores}

Los cuatro vectores quedaron mapeados a OWASP Top 10 2021 y a MITRE ATT\&CK con fuentes
oficiales verificadas, y ese mapeo cumplió la función que el objetivo le asigna: definió las cinco
clases canónicas del dataset y las categorías sobre las que se calculan todas las métricas del
OE3. La precisión del mapeo se reforzó durante el proyecto: Path Traversal/LFI se ubicó en
A01:2021 (Broken Access Control) y no en A03:2021 (Injection), lo que muestra que la selección de
los cuatro vectores responde a un criterio operativo (todos materializan su payload en la petición
HTTP y son detectables con las mismas características) y no a su pertenencia a una sola categoría
de OWASP. El OE1 se cumple.

% ------------------------------------------------------------------------------
\section{Objetivo específico 2: dataset etiquetado}

El corpus de 383{,}202 peticiones etiquetadas supera 3.8 veces el mínimo de 100{,}000 muestras. Su
representatividad se sustenta en la diversidad de origen: siete fuentes públicas auditadas (más de
la mitad de las filas proviene de la colección CAPEC y una cuarta parte de modsec-learn), entre
ellas registros reales de un servidor en producción protegido por ModSecurity \cite{lucz2025}
(1.0~\% de las filas), complementadas con dos listas públicas de payloads y con conjuntos propios,
con las cinco clases presentes en las tres particiones. Construir el corpus a partir de fuentes
públicas, en lugar de generarlo en su totalidad en el entorno controlado, aportó un volumen y una
variedad de técnicas de ataque que el entorno no habría podido producir en el tiempo disponible;
el entorno controlado se empleó para la validación, que es donde su reproducibilidad aporta más
valor. El desbalance de clases se trató con ponderación en el entrenamiento, y su efecto queda
acotado por los resultados: incluso la clase minoritaria, Command Injection (2.6~\% del conjunto de
prueba), alcanzó un F1 de 0.9239. El objetivo pide balancear tráfico legítimo y ataques: el corpus contiene
ambos en las cinco clases y en las tres particiones, pero no los equilibra en conteos (el 59.3~\%
del conjunto de prueba es SQLi), de modo que el criterio se cumple como presencia de ambos tipos de
tráfico y no como proporciones iguales. El etiquetado determinista y la partición verificada por hash
permiten reconstruir cada partición a partir de sus fuentes públicas, con la excepción de los
conjuntos sintéticos de Command Injection, cuyo generador no está versionado. El OE2 se cumple.

% ------------------------------------------------------------------------------
\section{Objetivo específico 3: efectividad de la detección}
\label{sec:dis-oe3}

El criterio del OE3 exige un F1 de al menos 0.80 en tres de las cuatro categorías. El sistema lo
alcanzó en las cuatro, con un margen mínimo de 0.12 sobre el criterio (Command Injection, 0.9239) y
un F1 macro de 0.9776. El AUC-ROC de entre 0.99899 y 0.99982 en las categorías de ataque indica que la capacidad de discriminar
cada categoría no depende de un umbral particular. La precisión de 0.9984 en la clase legítima
muestra que, dentro de la distribución del corpus, el modelo cumple el balance entre falsos
negativos y falsos positivos que el objetivo plantea para un entorno sin analista dedicado.

La cobertura sobre payloads no vistos confirma que estas métricas no dependen de memorizar el
entrenamiento. La auditoría de fuga fue decisiva para esa conclusión: en Command Injection, el
100~\% de la Ronda~4 correspondía en parte a plantillas presentes en el entrenamiento, y la cifra
atribuible a detección real, medida con el corpus limpio, es 91.5~\% con el paquete 0.1.0 y
96.0~\% con el tarball del modelo vigente, y 92.5~\% con el paquete 0.1.1 y el ejecutor corregido.
Esta última cifra es más baja que la anterior, pero no significa que la detección haya empeorado:
el ejecutor anterior volvía a codificar los payloads con separadores ya codificados, y 10 de ellos,
corrompidos así, se bloqueaban, mientras que enviados correctamente no se bloquean (HTTP 200); es la
primera medición correcta de esa técnica de evasión, y un límite real de la librería. En Path
Traversal, el paquete 0.1.0 alcanzó 85.0~\% con un intervalo de confianza cuyo límite inferior
(79.4~\%) quedaba por debajo del criterio, pero esa cifra no describe solo al paquete: sus fallos se
concentraban en las categorías de codificación y tenían dos causas independientes, un defecto del
middleware, que añadía una capa de codificación a la consulta, y un defecto del ejecutor de pruebas,
que volvía a codificar payloads ya codificados. Ninguna corrección bastó por sí sola: con el
middleware corregido y el ejecutor anterior los fallos fueron los mismos 10 que sin la corrección, y
desaparecieron al corregir también el ejecutor. Con el paquete 0.1.1 instalado desde el registro de
npm, la cobertura de Path Traversal es de 200 de 200 (intervalo de Wilson al 95~\%: 98.1--100.0~\%),
por encima del criterio con un margen amplio. Esa cifra tiene dos límites. No se midió el paquete
0.1.0 con el ejecutor corregido, por lo que no puede repartirse la mejora entre el modelo, el
middleware y el ejecutor. Y se obtuvo sobre el mismo corpus que sirvió para encontrar los
defectos, de modo que es optimista por construcción y debe confirmarse con un corpus nuevo. En XSS,
la detección se sustenta en el F1 de 0.9877 sobre el conjunto de prueba, en el 98~\% de la suite
extremo a extremo, en el 97.3~\% de la Ronda~4 y en el único caso que ZAP confirmó sobre la
aplicación de referencia (\texttt{POST /profile}, bloqueado por todos los paquetes evaluados), un
resultado real pero con n=1 que no tiene el peso estadístico de las otras tres categorías. El OE3
se cumple, sin reservas en las métricas de clasificación y con dos reservas en la cobertura: el
n=1 de XSS y que la medición vigente de Path Traversal y de Command Injection no es independiente
del corpus que motivó las correcciones.

% ------------------------------------------------------------------------------
\section{Defensa en profundidad}
\label{sec:dis-defensa}

En la Ronda~4, ModSecurity con OWASP CRS en nivel de paranoia 1 bloqueó el 93.2~\% de los ataques,
\textit{logSguarDian} el 98.8~\% y ambas capas el 99.5~\%; de los 40 ataques que el WAF dejó pasar,
la librería bloqueó 37. Ese corpus tiene fuga confirmada en Command Injection, pero ninguno de los
40 ataques que el WAF no bloqueó pertenece a esa categoría (fueron 11 de SQL Injection, 3 de XSS y
26 de Path Traversal). Por igualdad normalizada, 7 de los 113 payloads de XSS y ninguno de los 77
de SQL Injection ni de los 200 de Path Traversal coinciden con alguna partición; por subcadena,
26 de SQL Injection y 13 de XSS, un criterio permisivo que el trabajo de graduación de Diego
Valenzuela \cite{valenzuela2026ia} (sección VIII.3.5) califica de límite superior. Con esa salvedad, el resultado apunta al
posicionamiento de \textit{logSguarDian} definido en el Marco Teórico: una capa complementaria al
WAF, que detiene la mayor parte de lo que supera la primera capa, y no un sustituto.

La comparación sobre el corpus limpio de la Ronda~5, que no tiene la fuga de la Ronda~4, apunta
en la misma dirección, con una reserva importante: sus cifras de Path Traversal y de Command
Injection se midieron con el ejecutor anterior, que alteraba los payloads ya codificados y afecta
también al WAF, y no se repitieron con el ejecutor corregido, de modo que esas dos categorías son
provisionales. Con esa reserva, el WAF solo y la librería sola tienen perfiles distintos: en SQL Injection son
prácticamente equivalentes (99.5~\% y 100~\%, una diferencia de un payload); en Path Traversal el
WAF bloquea más (98.5~\% frente a 85.0~\%); y en Command Injection la librería bloquea más
(91.5~\% frente a 68.0~\%), la categoría en la que las firmas del WAF son más débiles.
Las dos capas combinadas bloquean el 100~\% de SQL Injection y de Path Traversal y el 98.5~\% de
Command Injection, más que cualquiera por separado en esa última, y de los 68 ataques que el WAF
no bloqueó, la combinación bloqueó 65. Esto indica que las capas detectan cosas distintas y que
una no subsume a la otra, con la salvedad de que la ventaja del WAF en Path Traversal queda
descartada para la librería con el paquete 0.1.1 (200 de 200), y de que la desventaja del WAF en
Command Injection (68.0~\%) debe confirmarse repitiendo las configuraciones con el ejecutor
corregido, porque la alteración de los separadores codificados pudo afectarla. Las Rondas 1 a 3 aportaron
además un requisito de integración: la librería debe montarse después de los analizadores de
cuerpo de la aplicación, porque con \texttt{multipart/form-data} procesado por ruta no veía el
contenido de la petición.

% ------------------------------------------------------------------------------
\section{Contraste con los antecedentes y el marco teórico}

Las métricas de clasificación de este trabajo están al nivel de las reportadas en los
Antecedentes: F1 de 0.9956 en SQL Injection frente a un F1 de 0.98 del Random Forest de Peralta-Garcia et al.
\cite{peralta2024sql}, y F1 de 0.9239 en Command Injection frente a una exactitud de 99.3~\% y un
\textit{recall} de 98.2~\% de Wang et al. \cite{seyyar2024}\footnote{Wang et al. reportan 99.3~\% en
el resumen y 99.21~\% en la comparación de la sección de resultados.}, cuyo modelo de aprendizaje
profundo se evaluó en una GPU y no reporta latencia de inferencia ni despliegue como middleware. La comparación numérica directa tiene límites, porque los
conjuntos de datos, las métricas y las unidades de análisis difieren. Las contribuciones de este
trabajo están en lo que los anteriores no reportan. Frente al WAF de Durmuşkaya y Bayraklı \cite{durmuskaya2025waf}, el más cercano en arquitectura, \textit{logSguarDian} opera dentro del proceso sin un salto de red
adicional; y, a diferencia de los trabajos revisados en los Antecedentes, este proyecto mide la latencia agregada
bajo carga sostenida, evalúa la detección sobre payloads filtrados contra el entrenamiento y prueba
la librería publicada en aplicaciones de terceros.

Frente al Marco Teórico, los resultados confirman dos de sus planteamientos y obligan a matizar un
tercero. Confirman el posicionamiento de un RASP como capa complementaria al WAF y no como su
sustituto: la librería, que analiza la petición dentro del proceso, y el WAF perimetral detectan
conjuntos parcialmente distintos (Sección~\ref{sec:dis-defensa}). Confirman también la descripción
del Isolation Forest como detector que modela el tráfico legítimo y marca lo atípico: detecta el
91.25~\% de los ataques del conjunto de prueba con 5.46~\% de falsos positivos. Obligan a matizar
la expectativa de que ese detector aporte por sí solo a las variantes desconocidas, porque sus
marcas sobre ataques que el Random Forest deja pasar existen (19 de 39 en validación \cite{valenzuela2026ia}) pero
costarían 866 bloqueos indebidos, y por eso solo registra. La advertencia de Kapoor y Narayanan
\cite{kapoor2023} sobre características no legítimas y conjuntos de prueba de la misma distribución
que el entrenamiento resultó pertinente: la cabecera \texttt{User-Agent} y la saturación de
casi-duplicados son ejemplos concretos en este corpus (Secciones~\ref{sec:dis-despliegue} y
\ref{sec:dis-limitaciones}).

% ------------------------------------------------------------------------------
\section{Condiciones de despliegue}
\label{sec:dis-despliegue}

Dos evaluaciones adicionales, una de ellas con un criterio que no se cumple, delimitan las
condiciones en que debe desplegarse.

\textbf{Volumen de datos.} En un escenario no contemplado en el protocolo, pero evaluado contra el mismo criterio de 15~ms, con la base de datos
sembrada y una ruta de listado que ya tardaba 63.84~ms sin middleware, el 7~\% de las peticiones
terminó por tiempo de espera. Los tiempos de espera se concentraron al final de cada repetición,
un patrón consistente con una cola que crece cuando el tiempo de servicio supera la tasa de
llegada. El criterio de 15~ms no se cumple en este escenario ($\Delta p95$ de 6{,}495.46~ms), lo que
limita el cumplimiento del objetivo general a las condiciones de carga normal y de ataque con
respuestas de la aplicación cortas; identifica el límite operativo de la versión actual y
orienta su siguiente iteración.

\textbf{Calibración con tráfico real.} En tres aplicaciones de terceros, con la configuración por
defecto, la librería bloqueó entre el 92.5~\% y el 100~\% del tráfico legítimo probado con un
\texttt{User-Agent} de navegador, con muestras pequeñas o poco diversas (5 peticiones en DVNA; 50
dirigidas a una sola ruta en node-api-goat); en Juice Shop, en modo \texttt{monitor}, el modelo asignó
veredicto de bloqueo al 97.1~\%, y en el estudio con navegadores reales se bloquearon 13 de 15
combinaciones de \texttt{User-Agent} (Cuadro~\ref{tab:res-terceros-ua}). La exclusión por fuente del trabajo del componente de Inteligencia Artificial
\cite{valenzuela2026ia}, con modelos distintos del evaluado, apunta en el mismo
sentido: sin \texttt{modsec\_learn}, \texttt{payloads\_csv} o \texttt{xss\_dataset} en el entrenamiento, el
bloqueo de sus filas legítimas sube de 0.1~\%, 0.4~\% y 0.2~\% a 76.0~\%, 27.0~\% y 94.5~\%. El Alcance definió la validación en el entorno
controlado como criterio principal y la validación externa como escenario complementario; estas
pruebas, realizadas por recomendación del asesor, cumplen ese papel complementario y aportan el
dato que el entorno controlado no podía dar. La hipótesis principal es un sesgo de representación
del corpus: solo el 0.4~\% de las filas legítimas tiene una cabecera \texttt{User-Agent} de más de
40 caracteres, frente al 33.3--97.7~\% de las clases de ataque, y esa característica es la de mayor
contribución en el modelo vigente (\texttt{rf\_v11}) según el análisis SHAP \cite{lundberg2017unified,valenzuela2026ia}, y lo era también en
\texttt{rf\_v3}. Esa hipótesis no es la única explicación: con el \texttt{User-Agent} por defecto de
Node.js, dvws-node bloqueó el 75~\% del tráfico legítimo, y Juice Shop el 39.4~\% con cURL, de modo
que el \texttt{User-Agent} amplifica el bloqueo sin explicarlo por completo, y no se verificó con un
corpus legítimo alternativo. Kapoor y Narayanan describen este tipo
de característica, que codifica el origen de los datos y no la propiedad de interés, como no
legítima, y advierten que un conjunto de prueba de la misma distribución que el entrenamiento no
representa la distribución de interés \cite{kapoor2023}. Si la hipótesis se confirma, el hallazgo sería corregible con datos y
calibración, sin cambiar la arquitectura, y define la configuración recomendada para la población
objetivo: instalación en modo \texttt{monitor} y calibración del umbral con el tráfico propio antes
de activar el bloqueo.

% ------------------------------------------------------------------------------
\section{Cumplimiento de los objetivos}

El Cuadro~\ref{tab:dis-cumplimiento} resume el cumplimiento de cada objetivo frente a su criterio
de aceptación y la evidencia que lo sustenta.

\begin{table}[H]
\centering
\small
\begin{tabular}{|L{2.4cm}|L{1.9cm}|L{8.5cm}|}
\hline
\textbf{Objetivo} & \textbf{Estado} & \textbf{Evidencia principal} \\ \hline
OE1 & Cumplido & Cuatro vectores mapeados a OWASP Top 10 2021 y MITRE ATT\&CK con fuente oficial; cinco clases canónicas que estructuran el dataset y la evaluación. \\ \hline
OE2 & Cumplido & 383{,}202 muestras etiquetadas (3.8 veces el mínimo) de siete fuentes auditadas, dos listas públicas y conjuntos propios, con etiquetado determinista y partición verificada. \\ \hline
OE3 & Cumplido, con reservas en la cobertura & F1 de 0.9239 a 0.9956 en 4 de 4 categorías; Isolation Forest con 91.25~\% de detección y 5.46~\% de falsos positivos; cobertura sobre payloads no vistos de 92.5 a 100~\% con el paquete vigente 0.1.1 y el ejecutor corregido (Path Traversal 200/200, intervalo 98.1--100.0~\%). Reservas: XSS con n=1, y que esa medición se hizo sobre el corpus que reveló los defectos de medición. \\ \hline
Objetivo general & Cumplido de forma parcial & Librería publicada en npm; variantes no vistas sobre el criterio de 80~\% con el Random Forest; extracción de 0.0273~ms y $\Delta p95$ de 12.08 y 12.85~ms. No se cumplen memoria (+332~MiB), CPU (+89~pp) ni la carga de volumen; el aporte del Isolation Forest a las variantes solo tiene evidencia indirecta (19 de 39 en validación). \\ \hline
\end{tabular}
\caption[Cumplimiento de los objetivos]{Cumplimiento de los objetivos frente a sus criterios de aceptación}
\label{tab:dis-cumplimiento}
\end{table}
{\footnotesize Nota. Elaboración propia con base en el Capítulo~\ref{cap:resultados}.}

% ------------------------------------------------------------------------------
\section{Hipótesis y validación de la solución}
\label{sec:dis-validacion}

El protocolo y el capítulo de Objetivos no formulan una hipótesis estadística contrastable, de
modo que no corresponde aceptarla o rechazarla. El objetivo general cumple ese papel como
afirmación de factibilidad: que un middleware de aprendizaje automático distribuido como paquete
npm puede detectar y clasificar ataques web sin comprometer el rendimiento del servidor anfitrión.
La evidencia respalda la primera parte (la detección y clasificación) y la mitad de la segunda
(la latencia en carga normal y de ataque), y no respalda el resto: los recursos de memoria y CPU,
el escenario de volumen y el comportamiento con tráfico legítimo real sin calibración previa. La
afirmación queda, por tanto, respaldada de forma parcial.

La validación de la solución combinó cuatro niveles de prueba. A nivel funcional, pruebas
automatizadas en integración continua (incluida la paridad numérica entre Python y ONNX) y la
instalación del paquete empaquetado en un directorio limpio contra una aplicación Express real.
A nivel de clasificación, la lectura del conjunto de prueba bloqueado por hash. A nivel de sistema,
las cuatro configuraciones de defensa sobre la aplicación de referencia, con un corpus filtrado
contra el entrenamiento, y el benchmark de rendimiento en Linux nativo con cinco repeticiones por
condición. Y a nivel externo, la integración en aplicaciones de terceros. Los criterios de
aceptación se aplicaron tal como los fijaron el Alcance y la Metodología, y el
Cuadro~\ref{tab:dis-cumplimiento} registra el estado de cada uno. Los dos niveles que más
información aportaron fueron el corpus
filtrado, que redujo la cobertura atribuible de Command Injection de 100~\% a 91.5~\%, y las
aplicaciones de terceros, una validación complementaria, que mostraron que las métricas del corpus no predicen el comportamiento
con tráfico legítimo real. Dos explicaciones quedan como posibles y no comprobadas: que el sesgo de
la cabecera \texttt{User-Agent} sea la causa principal de los bloqueos indebidos, y que el costo de
memoria y CPU se deba a la arquitectura de hilos y no a los modelos.

% ------------------------------------------------------------------------------
\section{Limitaciones}
\label{sec:dis-limitaciones}

Las amenazas a la validez descritas en la Sección~\ref{sec:met-validez} acotan el alcance de estas
conclusiones. El conjunto de prueba se volvió a bloquear en cada generación del modelo y se leyó
tres veces en la generación vigente; el corpus tiene una saturación de casi-duplicados que puede
inflar las métricas, sobre todo las de SQL Injection (aproximadamente 395 veces más pares casi
duplicados que filas, en una clase que es el 59.3~\% del conjunto de prueba) y que no se
cuantificó; el origen de los datos restaurados no está registrado, y el
17.8~\% de las filas de Command Injection del conjunto de prueba proviene de conjuntos sintéticos
sin generador versionado. Las evaluaciones usaron tres pares de modelos con la misma arquitectura
de ejecución pero no con la misma entrada (67 y 61 características en el par A; 69 y 63 en los otros dos), y el paquete con el que se hizo la medición principal de cobertura y se probaron las aplicaciones de terceros, el 0.1.0, no contiene el modelo cuyas
métricas de prueba se reportan, que solo está en la versión 0.1.1, posterior a esas mediciones.
La cobertura del modelo vigente sobre el corpus limpio no proviene de una evaluación estrictamente
ciega, y cada medición de cobertura corresponde a una sola corrida. La cobertura de
Path Traversal y de Command Injection con el paquete 0.1.1 se midió sobre el corpus que reveló los
defectos y solo para la Configuración~2; las mediciones con el ejecutor anterior, entre ellas la
comparación con el WAF, no se repitieron; la de XSS descansa en un único caso confirmado, y quedó
pendiente la validación con una startup en producción. Estas
limitaciones no modifican las cifras reportadas, pero cada una indica una verificación que
fortalecería la evidencia en trabajos posteriores.

% ------------------------------------------------------------------------------
\section{Sugerencias derivadas de los resultados}
\label{sec:dis-sugerencias}

Los resultados sugieren las siguientes líneas, que se presentan como indicios que dan paso a las
conclusiones y a las recomendaciones, no como compromisos ni instrucciones. Cada una indica la
evidencia de la que parte.
\begin{itemize}
    \item \textbf{Instalar primero en modo \texttt{monitor}.} Los bloqueos de tráfico legítimo en
    aplicaciones de terceros (92.5~\% en dvws-node con un \texttt{User-Agent} de navegador) sugieren
    que la calibración con el tráfico propio debería preceder al bloqueo.
    \item \textbf{Diversificar el corpus legítimo.} El 0.4~\% de filas legítimas con un
    \texttt{User-Agent} largo, frente a hasta 97.7~\% en las clases de ataque, sugiere que la
    representación de clientes HTTP es la primera variable que conviene probar.
    \item \textbf{Medir el costo por componente.} Un incremento casi constante de 332~MiB entre
    escenarios sugiere un costo fijo de la arquitectura, pero no permite repartirlo entre el Random
    Forest, el Isolation Forest y el entorno de ejecución.
    \item \textbf{Proteger al servidor ante saturación.} El patrón de tiempos de espera al final de
    cada repetición del escenario de volumen sugiere una cola que crece; un mecanismo de
    \textit{fail-open} por longitud de cola podría evitarlo.
    \item \textbf{Repetir la medición con un corpus nuevo.} Que el 100~\% de Path Traversal se
    obtuviera sobre el corpus que reveló los defectos sugiere confirmarlo con otra semilla antes de
    atribuirlo a la librería.
    \item \textbf{Aislar la contribución del Isolation Forest.} Que marque 19 de 39 ataques que el
    Random Forest omite, a un costo de 866 bloqueos indebidos, sugiere evaluarlo fuera del conjunto
    de validación y, quizá, con un umbral distinto.
\end{itemize}

\ifdefined\modoensamble
\else
	\section*{Referencias}
	\printbibliography[heading=none]
	\end{document}
\fi

\fi

\ifdefined\CAPconclusiones
	\newpage
	% ==============================================================================
% CONCLUSIONES
% ==============================================================================
% Solo componente de ciberseguridad (decisión del estudiante, 2026-10-09):
% no se incluyen conclusiones sobre entrenamiento o selección de modelos,
% que corresponden al componente de IA.
% Cada conclusión se deriva de una cifra de k-resultados.tex y de su
% interpretación en l-discusion.tex.
\ifdefined\modoensamble
\else
	\input{0-preambulo_standalone}
\fi

\chapter{Conclusiones}
\label{cap:conclusiones}

Las conclusiones siguen el orden de los objetivos: primero el objetivo general, después los
objetivos específicos y, al final, la hipótesis y la contribución al problema. Cada una indica el
objetivo al que responde, el grado de cumplimiento y el alcance de la evidencia que la sostiene.

\begin{enumerate}
    \item \textbf{Objetivo general: cumplido de forma parcial.}\label{conc:og} Se diseñó, implementó y publicó en
    el registro público de npm la librería \texttt{logsguardian} (versión 0.1.0, 30 de agosto de
    2026, y versión 0.1.1, 11 de octubre de 2026), un middleware RASP para Express que integra un
    extractor de 76 características y un modelo híbrido (Random Forest e Isolation Forest)
    serializado en ONNX y ejecutado en hilos de trabajo, con una diferencia máxima de
    $2.42\times10^{-7}$ entre las predicciones de Python y de ONNX. La librería intercepta las
    peticiones, clasifica las cuatro categorías de ataque y bloquea con el Random Forest. No se
    cumplen los criterios de memoria, CPU ni carga de volumen (conclusión~\ref{conc:rendimiento}), y
    el aporte del Isolation Forest a la detección de variantes desconocidas solo tiene evidencia
    indirecta (conclusión~\ref{conc:variantes}). La versión 0.1.0 contiene la generación de modelo
    \texttt{rf\_v10}/\texttt{if\_v9}; la versión 0.1.1, vigente, contiene \texttt{rf\_v11}/\texttt{if\_v10},
    que es la evaluada sobre el conjunto de prueba, y la corrección de un defecto del middleware.

    \item \textbf{Objetivo específico 1: cumplido.} Los cuatro vectores de ataque intra-request se
    mapearon a OWASP Top 10 2021 (SQL Injection, XSS y Command Injection en A03:2021; Path
    Traversal/LFI en A01:2021) y a MITRE ATT\&CK (T1190 y T1189), y ese mapeo estructuró las cinco
    clases canónicas del dataset y las categorías sobre las que se calculan las métricas del
    objetivo específico~3.

    \item \textbf{Objetivo específico 2: cumplido.} Se construyó un corpus de 383{,}202 peticiones
    HTTP etiquetadas, 3.8 veces el mínimo de 100{,}000, a partir de siete fuentes públicas
    auditadas, dos listas públicas de payloads y conjuntos propios, con etiquetado determinista y
    una partición verificada y bloqueada por hash. El desbalance de clases (59.3~\% SQLi frente a
    2.6~\% Command Injection) no impidió que la clase minoritaria alcanzara un F1 de 0.9239.

    \item \textbf{Objetivo específico 3, clasificación: cumplido.}\label{conc:clasificacion} Sobre el conjunto de prueba
    bloqueado, el sistema obtuvo un F1 de entre 0.9239 (Command Injection) y 0.9956 (SQL Injection)
    en las cuatro categorías, frente al criterio de 0.80 en al menos tres, y el Isolation Forest
    detectó el 91.25~\% de los ataques con 5.46~\% de falsos positivos, frente a los criterios de
    al menos 50~\% y como máximo 10~\%. Estas cifras miden la separación entre clases dentro del
    corpus y no la generalización a tráfico externo: el conjunto de prueba se leyó tres veces en la
    generación vigente y el corpus tiene saturación de casi-duplicados, en particular en SQL
    Injection (unos 395 pares casi duplicados por fila), un efecto que este trabajo no cuantificó.

    \item \textbf{Objetivo específico 3, cobertura sobre payloads no vistos: cumplido, con reservas.}\label{conc:cobertura}
    Sobre 600 payloads filtrados contra el entrenamiento, la cobertura fue, con el paquete 0.1.0
    (tarball de la Ronda~5), de 100~\% en SQL Injection, 85.0~\% en Path Traversal y 91.5~\% en
    Command Injection; con el tarball del modelo vigente, de 99.5~\%, 95.0~\% y 96.0~\%; y con el
    paquete 0.1.1 y el ejecutor de pruebas corregido, de 99.5~\%, 100.0~\% (200 de 200; intervalo de
    Wilson al 95~\% de 98.1 a 100.0~\%) y 92.5~\%. Los fallos de Path Traversal de la primera
    medición tenían dos causas independientes, un defecto del middleware y otro del ejecutor, ambos
    corregidos; la cobertura de Command Injection bajó de 96.0~\% a 92.5~\% porque la cifra anterior
    incluía 10 payloads que solo se bloqueaban en su forma corrompida. Las reservas son dos: la
    cobertura de XSS descansa en un único caso confirmado con OWASP ZAP (1 de 1 en las
    Configuraciones 2, 3a y 3b, n=1), que no permite afirmar el criterio del 80~\% para esa
    categoría y que carece de un corpus limpio; y la medición con el paquete 0.1.1 se hizo solo en
    la Configuración~2 y sobre el corpus que reveló los defectos, por lo que no es independiente.

    \item \textbf{Objetivo general, rendimiento del servidor anfitrión: no cumplido en su
    totalidad.}\label{conc:rendimiento} El extractor cumple el criterio en aislamiento: el p95 con
    \texttt{worker\_threads} es de 0.0273~ms, frente al límite de 5~ms. El sobrecosto del middleware
    completo es de 12.08~ms en carga normal y de 12.85~ms en carga de ataque (límite de 15~ms), y la
    extracción no lo explica. El extractor inicial de 72 características y el paso a
    \texttt{worker\_threads} los implementó el autor de este trabajo; el análisis y la justificación
    de cada característica, el contrato de entrada, las cuatro características adicionales y las
    correcciones del extractor corresponden al componente de Inteligencia Artificial. El criterio
    de latencia no se cumple en el escenario de volumen, donde el $\Delta p95$ fue de 6{,}495.46~ms
    y el 7.00~\% de las peticiones agotó el tiempo de espera con el middleware activo. Los
    indicadores de recursos del protocolo tampoco se cumplen: el middleware agregó entre 331.9 y
    375.9~MiB de memoria según el escenario, más del doble del límite de 150~MB (143.1~MiB), y
    +529.6~\% a +614.4~\% de CPU, frente a un límite de 15~\%. El incremento de memoria es casi
    constante entre escenarios, pero su reparto por componente no se midió.

    \item \textbf{Objetivo general, variantes desconocidas e Isolation Forest.}\label{conc:variantes}
    La evidencia directa de que la librería bloquea variantes desconocidas de los vectores cubiertos
    es la cobertura del Random Forest descrita en la conclusión~\ref{conc:cobertura} (4 de los 200 payloads de SQL
    Injection coinciden por igualdad normalizada con alguna partición; ninguno de Path Traversal ni
    de Command Injection). La contribución del Isolation Forest solo tiene evidencia indirecta: en
    el conjunto de validación del componente de Inteligencia Artificial \cite{valenzuela2026ia} marca 19 de los 39 ataques
    que el Random Forest deja pasar, pero bloquear sus marcas agregaría 866 benignos bloqueados de
    más, y por eso solo registra. Esas cifras son optimistas, porque el umbral del Isolation Forest
    se eligió sobre el mismo conjunto de validación, y no se evaluaron sobre el conjunto de prueba.

    \item \textbf{Objetivo general, defensa en profundidad.}\label{conc:waf} En la Ronda~4, ModSecurity con OWASP CRS
    bloqueó el 93.2~\% de 590 payloads, la librería el 98.8~\% y ambas capas el 99.5~\%; la librería
    bloqueó 37 de los 40 ataques que el WAF dejó pasar, ninguno de ellos de Command Injection, la
    categoría con fuga de datos confirmada en ese corpus. En la Ronda~5, sobre un corpus sin esa fuga,
    el WAF solo bloqueó el 68.0~\% de Command Injection frente al 91.5~\% de la librería, y ambas
    capas el 98.5~\%. Las cifras de Path Traversal y de Command Injection de la Ronda~5 se midieron
    con el ejecutor anterior a la corrección, que alteraba los payloads ya codificados y afecta
    también al WAF, y no se repitieron con el ejecutor corregido; son provisionales. La evidencia
    respalda el uso de la librería como capa complementaria a un WAF y no como su sustituto.

    \item \textbf{Objetivo general, falsos positivos y generalización.}\label{conc:fp} Las métricas del corpus no
    predijeron el comportamiento con tráfico legítimo real, y el corpus benigno es la debilidad
    principal. Con un \texttt{User-Agent} de navegador, la librería en modo \texttt{block} bloqueó el
    92.5~\% del tráfico legítimo en dvws-node (37 de 40; el modelo exacto de esa medición no consta),
    con muestras preliminares en DVNA (5 de 5) y de una sola ruta en node-api-goat (50 de 50). En
    OWASP Juice Shop, en modo \texttt{monitor}, el modelo asignó veredicto de bloqueo al 97.1~\% de
    las peticiones legítimas (303 de 312). En la aplicación de referencia, con otro modelo, ninguna
    petición legítima de Artillery recibió un 403, y con cadenas de \texttt{User-Agent} de
    navegadores reales se bloquearon 13 de 15 combinaciones, frente al 0.15~\% del conjunto de
    prueba. La exclusión por fuente del componente de Inteligencia Artificial \cite{valenzuela2026ia} apunta en el mismo
    sentido: al excluir \texttt{modsec\_learn}, \texttt{payloads\_csv} o \texttt{xss\_dataset} del
    entrenamiento, el bloqueo de filas legítimas de esas fuentes sube de 0.1~\%, 0.4~\% y 0.2~\% a
    76.0~\%, 27.0~\% y 94.5~\%. El sesgo de representación de la cabecera \texttt{User-Agent} es la
    explicación principal, pero no la única (dvws-node bloqueó el 75~\% con el \texttt{User-Agent}
    por defecto de Node.js y Juice Shop dio 39.4~\% con cURL) y no se verificó con un corpus legítimo
    alternativo. En su estado actual, la librería no es apta para instalarse en modo de bloqueo sin
    calibración previa con el tráfico de la aplicación.

    \item \textbf{Trazabilidad y validez de las métricas.}\label{conc:trazabilidad} La auditoría de fuga entre los corpus de
    evaluación y de entrenamiento redujo la cobertura atribuible de Command Injection de 100~\% a
    91.5~\%. El filtro de exclusión, el bloqueo del conjunto de prueba con verificación de hash y el
    manifiesto de sumas de verificación fueron condiciones necesarias para que las cifras reflejaran
    capacidad de detección y no memorización. Esos controles no evitaron que el conjunto de prueba
    se leyera tres veces en la generación vigente, que el origen de los datos restaurados quedara
    sin registrar, ni que, hasta la publicación de la versión 0.1.1, el paquete publicado y las
    métricas de prueba correspondieran a generaciones de modelo distintas, lo que limita la
    trazabilidad de las métricas.

    \item \textbf{Hipótesis.} El trabajo no formuló una hipótesis contrastable; el objetivo general
    cumple ese papel como afirmación de factibilidad: que un middleware de aprendizaje automático
    distribuido como paquete npm puede detectar y clasificar ataques web sin comprometer el
    rendimiento del servidor. La evidencia la respalda en la detección y clasificación y en la
    latencia bajo carga normal y de ataque, y es insuficiente o contraria en memoria, CPU, carga de
    volumen y comportamiento con tráfico legítimo real sin calibración, por lo que queda respaldada
    de forma parcial.

    \item \textbf{Contribución al problema.} El problema de partida es la falta de protección
    accesible a nivel de aplicación en organizaciones sin equipo de seguridad dedicado. El trabajo
    aporta una librería de instalación directa desde npm, sin infraestructura externa ni licencias
    de pago, con un extractor de 76 características y un modelo híbrido de procedencia documentada,
    y establece con evidencia dónde la solución funciona (clasificación de las cuatro categorías,
    cobertura de variantes y latencia bajo carga normal y de ataque) y dónde no (recursos, volumen y
    falsos positivos en aplicaciones de terceros). No demuestra que la herramienta sea adecuada para
    producción, y la validación con una startup quedó fuera de la evidencia del proyecto.
\end{enumerate}

\ifdefined\modoensamble
\else
	\section*{Referencias}
	\printbibliography[heading=none]
	\end{document}
\fi

\fi

\ifdefined\CAPrecomendaciones
	\newpage
	% ==============================================================================
% RECOMENDACIONES
% ==============================================================================
% Solo componente de ciberseguridad. Cada recomendación responde a una
% conclusión de m-conclusiones.tex o a una restricción declarada en
% j-metodologia.tex. La recomendación de validación con una startup se
% incorporó por pedido explícito del estudiante (2026-10-09).
\ifdefined\modoensamble
\else
	\input{0-preambulo_standalone}
\fi

\chapter{Recomendaciones}
\label{cap:recomendaciones}

Las recomendaciones se derivan de las conclusiones (Capítulo~\ref{cap:conclusiones}), de las
limitaciones (Sección~\ref{sec:dis-limitaciones}) y de las restricciones declaradas en la
Metodología (Sección~\ref{sec:met-restricciones}). Formalizan las sugerencias del capítulo de
Análisis de resultados, que se presentaron allí como indicios. Cada una indica la acción propuesta,
su justificación en la evidencia del trabajo y el resultado que se espera obtener. Se agrupan por
destinatario y por tipo de acción.

\section{Para quienes adopten la librería}

\begin{enumerate}
    \item \textbf{Instalar primero en modo \texttt{monitor}.} \textit{Acción:} instalar la librería
    en modo \texttt{monitor}, revisar durante un período de observación los eventos registrados con
    la CLI (\texttt{attacks list} y \texttt{endpoints top}), ajustar el umbral con el tráfico propio
    de la aplicación y solo después activar el modo \texttt{block}. \textit{Justificación:} con un
    \texttt{User-Agent} de navegador, la librería bloqueó el 92.5~\% del tráfico legítimo de
    dvws-node (Conclusión~\ref{conc:fp}). \textit{Resultado esperado:} conocer la tasa de bloqueos
    indebidos de la propia aplicación antes de afectar a usuarios reales.

    \item \textbf{Usar la librería como capa complementaria detrás de un WAF.}
    \textit{Acción:} desplegarla detrás de un WAF y no como su sustituto. \textit{Justificación:}
    bloqueó 37 de los 40 ataques que ModSecurity con OWASP CRS dejó pasar en la Ronda~4, y en
    Command Injection el WAF solo bloqueó el 68.0~\% frente al 91.5~\% de la librería en la
    Ronda~5, una comparación provisional (Conclusión~\ref{conc:waf}). \textit{Resultado esperado:}
    mayor cobertura que la de cada capa por separado.

    \item \textbf{Montar el middleware después de todos los analizadores de cuerpo.}
    \textit{Acción:} registrar el middleware después de los analizadores de cuerpo de la aplicación,
    incluidos los de \texttt{multipart/form-data}. \textit{Justificación:} en las Rondas 1 a 3, la
    librería no inspeccionó 7 de 23 payloads de XSS enviados como \texttt{multipart/form-data}
    mientras el analizador se montaba por ruta, y bloqueó los 23 al montarlo de forma global antes
    del middleware. \textit{Resultado esperado:} que la librería inspeccione el contenido de todas
    las peticiones.

    \item \textbf{No desplegar la versión actual en rutas de alto volumen de datos.}
    \textit{Acción:} excluir de la protección, o paginar, las rutas que devuelven grandes volúmenes
    de datos por petición, hasta que se resuelva la saturación (recomendación~\ref{rec:cola}).
    \textit{Justificación:} en el escenario de volumen, el $\Delta p95$ fue de 6{,}495.46~ms y el
    7.00~\% de las peticiones agotó el tiempo de espera (Conclusión~\ref{conc:rendimiento}).
    \textit{Resultado esperado:} evitar una degradación de la disponibilidad del servidor anfitrión.
\end{enumerate}

\section{Mejoras técnicas de la librería}

\begin{enumerate}
    \setcounter{enumi}{4}
    \item \textbf{Usar el modo \texttt{monitor} por defecto.} \textit{Acción:} cambiar el modo por defecto del paquete, hoy \texttt{block}, a \texttt{monitor} en su próxima versión y documentar en el README cómo
    activar el bloqueo. \textit{Justificación:} las pruebas en aplicaciones de terceros se hicieron
    con la configuración por defecto y arrojaron bloqueos indebidos altos
    (Conclusión~\ref{conc:fp}). \textit{Resultado esperado:} que una instalación sin configuración
    no bloquee tráfico legítimo.

    \item \textbf{Verificar y, si se confirma, corregir el sesgo de la cabecera
    \texttt{User-Agent}.} \textit{Acción:} reunir tráfico legítimo de navegadores y clientes HTTP
    diversos (incluidos cURL y los \texttt{User-Agent} por defecto de Node.js), evaluarlo con la
    configuración vigente, recalibrar el umbral de bloqueo con tráfico de varias aplicaciones y
    repetir las mediciones del Cuadro~\ref{tab:res-terceros-ua}. Si la hipótesis se confirma, la
    ampliación del corpus de entrenamiento corresponde al componente de Inteligencia Artificial y
    queda fuera de este documento. \textit{Justificación:} solo el 0.4~\% de las filas legítimas
    tiene un \texttt{User-Agent} largo, frente al 33.3--97.7~\% de las clases de ataque, y esa
    característica es la de mayor contribución del modelo (Conclusión~\ref{conc:fp}); la hipótesis
    no se verificó con un corpus legítimo alternativo. \textit{Resultado esperado:} una reducción
    medible de los bloqueos indebidos o, si no se produce, el descarte de esa hipótesis.

    \item \label{rec:cola}\textbf{Proteger al servidor ante saturación.} \textit{Acción:}
    instrumentar la cola de los hilos de inferencia durante el escenario de volumen y evaluar un
    mecanismo que deje pasar las peticiones (\textit{fail-open}) cuando la cola supere un límite.
    \textit{Justificación:} el 7.00~\% de tiempos de espera se concentró al final de cada
    repetición, un patrón consistente con una cola que crece, aunque la hipótesis no se comprobó
    (Conclusión~\ref{conc:rendimiento}). \textit{Resultado esperado:} confirmar o descartar la
    saturación y reducir los tiempos de espera del escenario de volumen.

    \item \textbf{Reducir el costo fijo de memoria y CPU.} \textit{Acción:} medir por separado el
    costo de cada componente de la librería (solo el Random Forest, el Random Forest con un hilo
    del Isolation Forest y la configuración actual) y evaluar \texttt{IF\_POOL\_SIZE=1} en lugar
    del valor actual de 2, repitiendo el benchmark con el mismo análisis pareado por repetición.
    \textit{Justificación:} el middleware agregó entre 331.9 y 375.9~MiB, más del doble del límite
    de 143.1~MiB, y alrededor de un núcleo de CPU, y ese costo no se atribuyó a ningún componente
    (Conclusión~\ref{conc:rendimiento}). \textit{Resultado esperado:} identificar el componente que
    domina el costo y la reducción que se logra con cada cambio de configuración.

    \item \textbf{Reducir la confusión entre categorías en los eventos registrados.}
    \textit{Acción:} revisar cómo se asigna la clase que se guarda en el evento y comprobarla con
    una matriz de confusión del registro. \textit{Justificación:} en la repetición de la
    Configuración~2 con el par~B, los 200 payloads de Command Injection se registraron como SQLi, y
    en DVNA, 27 de los 37 bloqueados de esa categoría se clasificaron como SQLi, aunque la decisión
    de bloqueo no cambió. \textit{Resultado esperado:} que el catálogo de ataques de la CLI sea
    confiable como herramienta de análisis.
\end{enumerate}

\section{Validaciones pendientes}

\begin{enumerate}
    \setcounter{enumi}{9}
    \item \textbf{Validar la librería en el entorno productivo de una startup guatemalteca.}
    \textit{Acción:} instalarla en modo \texttt{monitor} en una aplicación real y medir los
    bloqueos indebidos y el $\Delta p95$ con tráfico de usuarios. \textit{Justificación:} esta
    validación formaba parte del alcance complementario del protocolo y no se concretó por falta de
    tiempo para contactar a una startup durante el período de pruebas (Conclusión~\ref{conc:fp}).
    \textit{Resultado esperado:} medir los falsos positivos y el impacto en latencia con tráfico
    real y no solo con aplicaciones vulnerables de terceros.

    \item \textbf{Repetir las mediciones con el paquete vigente y con un corpus nuevo.}
    \textit{Acción:} con el paquete 0.1.1 del registro y el ejecutor corregido, repetir las cuatro
    configuraciones de la Ronda~5 (incluidas la 3a y la 3b), la cobertura con el paquete 0.1.0,
    las aplicaciones de terceros y el benchmark de rendimiento; repetir también la cobertura con
    un corpus nuevo, generado con otra semilla y el mismo filtro de exclusión. Registrar el SHA-256
    del archivo publicado (el de la 0.1.1 es \texttt{abf8d5c4...}) y de sus modelos, y medir el
    desajuste entre la cadena de consulta con la que se entrenó y la que el middleware entrega al
    extractor. \textit{Justificación:} con la 0.1.1 solo se midió la Configuración~2, sobre el
    corpus que reveló los defectos, y de las Configuraciones 1, 3a y 3b no consta el paquete
    (Conclusión~\ref{conc:cobertura}). \textit{Resultado esperado:} cifras independientes y atadas
    al archivo exacto evaluado, que confirmen o corrijan el 100.0~\% en Path Traversal y el
    92.5~\% en Command Injection, y una comparación con el WAF definitiva.

    \item \textbf{Construir un corpus limpio de XSS.} \textit{Acción:} generar payloads de XSS
    con OWASP ZAP u otras herramientas sobre más endpoints o aplicaciones, filtrarlos contra el
    entrenamiento, definir para cada payload su forma exacta en el cable para que el ejecutor no la
    altere, y repetir la comparación con el WAF sobre ambos corpus. \textit{Justificación:} la
    cobertura de XSS descansa en un único payload confirmado por ZAP (1 de 1), que no permite
    afirmar el criterio del 80~\% (Conclusión~\ref{conc:cobertura}). \textit{Resultado esperado:}
    una cobertura de XSS con intervalo de confianza comparable al de las otras categorías.

    \item \textbf{Monitorear la memoria en un escenario sostenido de condición única.}
    \textit{Acción:} ejecutar un escenario de 30~minutos o más con una sola condición estable.
    \textit{Justificación:} la corrida continua de 41.9~minutos mezcla condiciones y no permite
    descartar una fuga de memoria dentro de un proceso estable (Sección~\ref{sec:met-restricciones}).
    \textit{Resultado esperado:} confirmar o descartar una fuga de memoria.

    \item \textbf{Medir la contribución del Isolation Forest con las pruebas de penetración.}
    \textit{Acción:} registrar en cada ronda de pruebas el veredicto \texttt{pass\_anomaly} de cada
    payload que el Random Forest no bloquea y de cada petición legítima, y reportar cuántos
    ataques no bloqueados quedan marcados como anómalos y cuántas alertas generan las peticiones
    legítimas. \textit{Justificación:} la contribución del detector de anomalías a la detección de
    variantes solo tiene evidencia indirecta: en un conjunto de validación ajeno a este trabajo marca
    19 de los 39 ataques que el Random Forest deja pasar, pero bloquear sus marcas agregaría 866
    benignos bloqueados \cite{valenzuela2026ia} (Conclusión~\ref{conc:variantes}).
    \textit{Resultado esperado:} estimar, con el tráfico de ataque de la evaluación del sistema,
    cuánto aporta el detector y cuántas alertas adicionales genera.
\end{enumerate}

\section{Reproducibilidad y entrega de datos}

\begin{enumerate}
    \setcounter{enumi}{14}
    \item \textbf{Publicar los vectores derivados y depurar el conjunto de pruebas extremo a
    extremo.} \textit{Acción:} publicar los vectores de características derivados como recurso
    descargable de una versión del repositorio, sin la columna \texttt{\_row\_hash} y después de
    revisar que ninguna característica permita reconstruir el texto de la petición; y sustituir, en
    el conjunto de pruebas extremo a extremo versionado, los payloads cuyo texto proviene de
    fuentes sin licencia especificada. \textit{Justificación:} el Alcance entrega los scripts y los
    vectores derivados, y las fuentes \texttt{XSS\_dataset.csv} y \texttt{command injection.csv}
    (3.4~\% de las filas) no tienen licencia especificada por su origen
    (Sección~\ref{subsec:met-f1}). \textit{Resultado esperado:} un repositorio público coherente
    con lo que afirma la auditoría de fuentes.

    \item \textbf{Versionar el generador de los conjuntos sintéticos de Command Injection.}
    \textit{Acción:} incorporar al repositorio el script que los genera, o declarar esa porción del
    corpus como no reproducible. \textit{Justificación:} el 17.8~\% de las filas de Command
    Injection del conjunto de prueba (261 de 1{,}468) proviene de esos conjuntos y su generador no
    está versionado (Conclusión~\ref{conc:trazabilidad}). \textit{Resultado esperado:}
    reproducibilidad completa de la partición de prueba.
\end{enumerate}

\section{Ampliación del alcance}

\begin{enumerate}
    \setcounter{enumi}{16}
    \item \textbf{Desacoplar la detección de Node.js y de npm.} \textit{Acción:} separar un
    núcleo de detección reutilizable y escribir adaptadores para otros ecosistemas (por ejemplo,
    \textit{frameworks} de Python o de Java, o un \textit{proxy} local), verificando en cada uno la
    paridad de las predicciones con la de referencia, como se hizo entre Python y ONNX.
    \textit{Justificación:} el Alcance excluyó otros entornos de ejecución, pero el extractor, los
    modelos en ONNX y el umbral no dependen del lenguaje; lo que está atado a Node.js es el
    adaptador de Express y la ejecución en hilos de trabajo (Sección~\ref{subsec:fuera-alcance}).
    \textit{Resultado esperado:} ampliar la población de aplicaciones que pueden usar la detección
    sin reescribir el modelo.

    \item \textbf{Exponer la librería como servidor del Protocolo de Contexto de Modelo (MCP).}
    \textit{Acción:} publicar un servidor MCP que permita a un asistente de programación consultar
    el catálogo de ataques y los \textit{endpoints} más atacados que ya ofrece la CLI, definiendo
    antes qué operaciones son de solo lectura y cuáles podrían cambiar el modo o el umbral.
    \textit{Justificación:} la CLI ya expone esa información, y un asistente con acceso de
    escritura ampliaría la superficie de ataque de la propia herramienta de protección.
    \textit{Resultado esperado:} consulta asistida de los eventos sin ampliar esa superficie.
\end{enumerate}

\section{Para trabajos de investigación similares}

\begin{enumerate}
    \setcounter{enumi}{18}
    \item \textbf{Definir un conjunto de evaluación externo desde el inicio.} \textit{Acción:}
    reunir tráfico legítimo de varias aplicaciones y no limitarse a una partición de la misma
    distribución que el entrenamiento. \textit{Justificación:} las métricas del corpus (F1 de
    0.9239 a 0.9956) no predijeron los bloqueos con tráfico legítimo real (Conclusión~\ref{conc:fp}).
    \textit{Resultado esperado:} métricas que reflejen el comportamiento en producción y no solo la
    separación entre clases dentro del corpus.

    \item \textbf{Medir y tratar los casi-duplicados antes de particionar.} \textit{Acción:}
    medirlos (por ejemplo, con MinHash), agrupar por plantilla al dividir en entrenamiento,
    validación y prueba, y repetir la evaluación con el corpus deduplicado. \textit{Justificación:}
    SQLi tiene unas 395 veces más pares casi duplicados que filas y es el 59.3~\% del conjunto de
    prueba, un efecto que este trabajo no cuantificó (Conclusión~\ref{conc:clasificacion}).
    \textit{Resultado esperado:} cuantificar cuánto de las métricas actuales depende de esa
    saturación.

    \item \textbf{Bloquear el conjunto de prueba una sola vez y versionar cada medición.}
    \textit{Acción:} bloquear el conjunto de prueba una vez para todo el proyecto, atar cada cifra
    al hash de los datos y del modelo que la produjeron, y guardar cada medición de rendimiento
    como un artefacto versionado con el commit, el hardware, la versión del entorno de ejecución y
    las muestras crudas separadas por condición, como se hizo con el benchmark del extractor y con
    la evidencia de la corrida de latencia. \textit{Justificación:} el valor del bloqueo cambió
    nueve veces y el conjunto de prueba se leyó tres veces en la generación vigente
    (Conclusión~\ref{conc:trazabilidad}). \textit{Resultado esperado:} cifras atables a un único
    conjunto de datos y a un único modelo.

    \item \textbf{Construir los corpus de evaluación con un generador versionado y un filtro de
    exclusión.} \textit{Acción:} versionar el generador y aplicar el filtro contra el
    entrenamiento desde la primera ronda, y reportar siempre la tasa de detección junto con la de
    falsos positivos. \textit{Justificación:} la Ronda~4 tuvo fuga confirmada en Command Injection,
    que redujo la cobertura atribuible de 100~\% a 91.5~\% al corregirse
    (Conclusión~\ref{conc:trazabilidad}). \textit{Resultado esperado:} cifras de detección que no
    estén infladas por solapamiento con el entrenamiento.
\end{enumerate}

\ifdefined\modoensamble
\else
	\section*{Referencias}
	\printbibliography[heading=none]
	\end{document}
\fi

\fi

% BIBLIOGRAFÍA
% ------------------------------------------------------------------------------
\ifdefined\CAPbibliografia
	\newpage
    \cleardoublepage\phantomsection
	\chapter{\bibname}
    \printbibliography[heading=none]
\fi

% ANEXOS
% ------------------------------------------------------------------------------
\ifdefined\CAPanexos
	\newpage
	% ==============================================================================
% ANEXOS
% ==============================================================================
% Datos tomados de los repositorios del proyecto:
%   - docs/model-provenance.md (logSguarDian): hashes de los pares de modelos.
%   - SHA-256 de los archivos .tgz de npm: calculados con npm pack y shasum el
%     2026-10-12.
%   - docs/cybersecurity-objectives-compliance.md: historial del bloqueo del
%     conjunto de prueba (git log --all -- training/splits/test.lock.sha256).
%   - docs/dataset-audit.md: auditoria de fuentes.
%   - docs/evidence/README.md (logSguarDian-vulnerable-project): corrida de
%     latencia en Linux.
\ifdefined\modoensamble
\else
	\input{0-preambulo_standalone}
\fi

\chapter{Anexos}
\label{cap:anexos}

\section{Procedencia de los modelos y de los paquetes}
\label{anx:procedencia}

El Cuadro~\ref{tab:anx-hashes} lista los archivos \texttt{rf.onnx} e \texttt{if.onnx} de cada par
de modelos descrito en el Cuadro~\ref{tab:pares-modelos}, con su suma SHA-256 completa. Las sumas
se calcularon sobre el contenido del objeto de git de cada archivo y, para los paquetes, sobre el
contenido extraído del archivo publicado en npm.

\begin{table}[H]
\centering
\small
\begin{tabular}{|L{3.6cm}|L{9.9cm}|}
\hline
\textbf{Archivo} & \textbf{SHA-256} \\ \hline
Par A, \texttt{rf.onnx} (\texttt{rf\_v10}) & {\scriptsize\texttt{180837255dffd596a38bf7c29032ccc313005cd70eca5c6238e23e35bdc0d4ec}} \\ \hline
Par A, \texttt{if.onnx} (\texttt{if\_v9}) & {\scriptsize\texttt{5de374a7f1d19176d7c66d3354e972d37b28a7306a87efa0ce8639278f489150}} \\ \hline
Par B, \texttt{rf.onnx} (\texttt{rf\_v11}, primer build) & {\scriptsize\texttt{a7c015f5a5fa3fabfd91b79ec9338dc7da6a68c392c19ded4b04cc6f0fb903b2}} \\ \hline
Par B, \texttt{if.onnx} (\texttt{if\_v10}, primer build) & {\scriptsize\texttt{3cf8e45c285de19bc86e6dc284e502cd2c24becb888747b9645bed4cfb224417}} \\ \hline
Par C, \texttt{rf.onnx} (\texttt{rf\_v11}, reentrenamiento) & {\scriptsize\texttt{25b407e649283606a231665a5a8c8ea1c2c7ad4fd0374a8945b3f119651b0870}} \\ \hline
Par C, \texttt{if.onnx} (\texttt{if\_v10}, reentrenamiento) & {\scriptsize\texttt{5be28ee8437643b2f13708eee1ac3175348a350ea6fb8f8daeb09bd5717f797e}} \\ \hline
\end{tabular}
\caption[Sumas SHA-256 de los pares de modelos]{Sumas SHA-256 de los archivos ONNX de cada par de modelos}
\label{tab:anx-hashes}
\end{table}
{\footnotesize Nota. Elaboración propia con datos de \texttt{docs/model-provenance.md} (repositorio
\texttt{logSguarDian}). El par~A es el contenido del paquete \texttt{logsguardian@0.1.0}; el par~C,
el del paquete \texttt{logsguardian@0.1.1}.}

El Cuadro~\ref{tab:anx-paquetes} reúne los dos paquetes publicados en el registro de npm y el umbral
del Isolation Forest que cada uno aplica.

\begin{table}[H]
\centering
\small
\begin{tabular}{|L{2.3cm}|L{2.5cm}|L{6.2cm}|L{2.5cm}|}
\hline
\textbf{Paquete} & \textbf{Publicación} & \textbf{SHA-256 del archivo \texttt{.tgz}} & \textbf{Umbral del IF} \\ \hline
\texttt{0.1.0} & 30 de agosto de 2026 & {\scriptsize\texttt{c5c1a1da57cfb873350ac0c8ab3e7bad395e8cbf507b18cc3272df447880dcb0}} & 0.002486 (constante) \\ \hline
\texttt{0.1.1} & 11 de octubre de 2026 (UTC) & {\scriptsize\texttt{abf8d5c47036d9c3a289bc36158f1caaa91b9504689aa97193e51c6808258d15}} & 0.004206 (leído de \texttt{parity\_report.json}) \\ \hline
\end{tabular}
\caption[Paquetes publicados en npm]{Paquetes \texttt{logsguardian} publicados en el registro de npm}
\label{tab:anx-paquetes}
\end{table}
{\footnotesize Nota. Elaboración propia. Las sumas de los archivos se calcularon descargando cada
versión del registro (\texttt{npm pack}); los umbrales provienen de \texttt{docs/decision-policy.md}
del repositorio \texttt{logSguarDian}. Los otros dos archivos que se mencionan en el texto, el
tarball de la Ronda~5 original (\texttt{e09120006dd1...}) y el tarball de validación
(\texttt{cf1238deb333...}), no se publicaron en el registro.}

\section{Historial del bloqueo del conjunto de prueba}
\label{anx:lock}

El bloqueo del conjunto de prueba es el SHA-256 de los identificadores por fila del conjunto,
ordenados y unidos por saltos de línea. El valor cambió nueve veces; el Cuadro~\ref{tab:anx-lock}
las lista. No hay un único bloqueo por generación de modelo, y el cambio del 23 de septiembre de
2026 entró en tres commits.

\begin{table}[H]
\centering
\small
\begin{tabular}{|L{2.2cm}|L{2.4cm}|L{4.8cm}|L{2.2cm}|L{2.2cm}|}
\hline
\textbf{Fecha} & \textbf{Commit} & \textbf{Motivo} & \textbf{Antes} & \textbf{Después} \\ \hline
11 jun 2026 & \texttt{cc00c3c} & Bloqueo inicial & --- & \texttt{3beeaba4} \\ \hline
23 jul 2026 & \texttt{df62619} & Nuevo conjunto, antes de la generación v3 & \texttt{3beeaba4} & \texttt{d1a2a284} \\ \hline
24 jul 2026 & \texttt{6867f28} & Antes de la generación v4 & \texttt{d1a2a284} & \texttt{7df5adb8} \\ \hline
24 jul 2026 & \texttt{59990ca} & Antes de la generación v6 & \texttt{7df5adb8} & \texttt{ee4b85db} \\ \hline
25 jul 2026 & \texttt{5dc7c06} & Antes de la generación v8 & \texttt{ee4b85db} & \texttt{251716bb} \\ \hline
26 jul 2026 & \texttt{fa46bdd} & Antes de la generación v9 & \texttt{251716bb} & \texttt{4a5e0b97} \\ \hline
2 ago 2026 & \texttt{b43de9c} & Antes de la generación v10 & \texttt{4a5e0b97} & \texttt{6c5f253d} \\ \hline
30 ago 2026 & \texttt{8cbe6c9} & Antes de la generación v11 & \texttt{6c5f253d} & \texttt{a3140ab3} \\ \hline
23 sep 2026 & \texttt{f15e7cc}, \texttt{72aabd6}, \texttt{6d1fa68} & Rebloqueo para que coincida con la generación v11 restaurada & \texttt{a3140ab3} & \texttt{7486c258} \\ \hline
\end{tabular}
\caption[Historial del bloqueo del conjunto de prueba]{Historial de cambios del bloqueo del conjunto de prueba (hashes truncados a 8 caracteres)}
\label{tab:anx-lock}
\end{table}
{\footnotesize Nota. Elaboración propia con datos de \texttt{docs/cybersecurity-objectives-compliance.md}
(repositorio \texttt{logSguarDian}), generados con \texttt{git log --all -{}- training/splits/test.lock.sha256}.}

\section{Auditoría de las fuentes de datos}
\label{anx:auditoria}

El Cuadro~\ref{tab:anx-auditoria} resume la auditoría de licencias de las nueve fuentes públicas
candidatas, verificada contra la API o la página de cada origen.

\begin{table}[H]
\centering
\small
\begin{tabular}{|L{3.8cm}|L{3.0cm}|L{3.3cm}|L{3.6cm}|}
\hline
\textbf{Fuente} & \textbf{Licencia} & \textbf{DOI} & \textbf{Decisión} \\ \hline
Colección CAPEC (SR-BH 2020) \cite{sureda2022srbh,sureda2022newmultilabel} & CC0-1.0 & 10.7910/DVN/OGOIXX & Incluida \\ \hline
modsec-learn \cite{pralab2024modseclearn,scano2024modseclearn} & MIT & No verificado & Incluida \\ \hline
OWASP ModSecurity, registros de un servidor protegido (Zenodo) \cite{lucz2025} & CC BY 4.0 & 10.5281/zenodo.17178461 & Incluida \\ \hline
Payloads.csv (Kaggle) \cite{shahane2021xss} & CC BY 4.0 & No asignado por Kaggle & Incluida \\ \hline
payload\_full.csv (Kaggle) \cite{thakur2025payloads} & MIT & No asignado por Kaggle & Incluida \\ \hline
XSS\_dataset.csv (Kaggle) \cite{shah2020xss} & No especificada por el origen & No asignado por Kaggle & Incluida; no se redistribuye \\ \hline
command injection.csv (Kaggle) \cite{pawase2024cmdi} & No especificada por el origen & No asignado por Kaggle & Incluida; no se redistribuye \\ \hline
Lista de payloads de Path Traversal de omurugur \cite{ugur2019traversal} & MIT & No asignado & Aprobada; sin analizador, no aporta filas \\ \hline
AIT-LDSv2 (russellmitchell) \cite{landauer2022ait,landauer2023maintainable} & CC-BY-NC-SA-4.0 & 10.5281/zenodo.5789064 & Excluida: sus etiquetas describen el rol del atacante, no el tipo de payload \\ \hline
\end{tabular}
\caption[Auditoría de las fuentes de datos]{Licencia, DOI y decisión de inclusión de cada fuente auditada}
\label{tab:anx-auditoria}
\end{table}
{\footnotesize Nota. Elaboración propia con datos de \texttt{docs/dataset-audit.md} (repositorio
\texttt{logSguarDian}). El número de filas que cada fuente aporta al corpus final está en el
Cuadro~\ref{tab:met-origenes}. Las listas SecLists y PayloadsAllTheThings no constan en esta
auditoría; su licencia (MIT) se consultó en sus repositorios el 10 de octubre de 2026.}

\section{Evidencia de la corrida de latencia en Linux}
\label{anx:latencia}

El Cuadro~\ref{tab:anx-latencia} identifica la corrida de integración continua de la que provienen
las cifras de latencia, memoria y CPU del Capítulo~\ref{cap:resultados}.

\begin{table}[H]
\centering
\small
\begin{tabular}{|L{4.2cm}|L{9.3cm}|}
\hline
\textbf{Campo} & \textbf{Valor} \\ \hline
Flujo de trabajo & \texttt{.github/workflows/latency-benchmark.yml} \\ \hline
Identificador de la ejecución & 36276690236 \\ \hline
Identificador del artefacto & 10917957439 \\ \hline
Commit & \texttt{f15a9e8} \\ \hline
Fecha & 26 de septiembre de 2026 \\ \hline
Ejecutor & Ubuntu (Azure), x86\_64, 2 vCPU, 7.8~GiB de memoria \\ \hline
Docker Engine / Compose & 28.0.4 / v2.38.2 \\ \hline
Modelos & Par~B (\texttt{rf.onnx} \texttt{a7c015f5...}, \texttt{if.onnx} \texttt{3cf8e45c...}) \\ \hline
Contenido del archivo & 30 reportes de Artillery (3 escenarios, 2 condiciones, 5 repeticiones), muestras de \texttt{docker stats} cada $\approx$5~s y análisis de memoria \\ \hline
Copia en el repositorio & \texttt{docs/evidence/latency-linux-2026-09-26.tgz} (SHA-256 {\scriptsize\texttt{b96d901408547331fb4859a5dac4810abdb1caaeb2f15f3378877c2290187cec}}) \\ \hline
\end{tabular}
\caption[Corrida de latencia en Linux]{Identificación de la corrida de latencia, memoria y CPU en Linux nativo}
\label{tab:anx-latencia}
\end{table}
{\footnotesize Nota. Elaboración propia con datos de \texttt{docs/evidence/README.md} (repositorio
\texttt{logSguarDian-vulnerable-project}). El artefacto original de GitHub Actions expira el 26 de
octubre de 2026; la copia del repositorio es la que se conserva.}

\ifdefined\modoensamble
\else
	\section*{Referencias}
	\printbibliography[heading=none]
	\end{document}
\fi

\fi

% GLOSARIO
% ------------------------------------------------------------------------------
\ifdefined\CAPglosario
	\newpage
	\glsaddall
	\printglossary
\fi

\end{document}